\section{Price metrics: exact small cases and necessary restrictions}
\label{sec:price-geometry}

This section uses the definition of a spectral point in the attached Anthropic
manuscript, Definition~A.1 (p.~19), and its definition of a price,
Definition~A.4 (p.~20). These spectral points are not required to be additive
under direct sums. All uses of the manuscript's price-space theorem below are
conditional on its checked local proof chain; no claim about the correctness of
the entire circuit lower-bound manuscript is implicit.

\subsection{A saturation lemma}

\begin{lemma}[PG-SAT: saturation of an upper functional]
\label{pg:saturation}
Let $p$ be an upper functional on tensors on a finite party set $I$, in the
sense of Anthropic \S A.4, and let $C\ne0$ satisfy
$p(C^n)=c^n$ for every integer $n\geq0$, where $c=p(C)>0$.
There is a spectral point $\varphi\leq p$ with $\varphi(C)=c$.
\end{lemma}
\begin{proof}
Set $q(t)=\inf_{n\geq0}p(tC^n)/c^n$. The sequence in this infimum is
nonincreasing. It follows that $q(1)=1$, that $q$ obeys the cover inequality,
and that $q\leq p$. For submultiplicativity use
$p(stC^{a+b})\leq p(sC^a)p(tC^b)$ and take independent infima over
$a,b$. Moreover $q(tC)=c q(t)$, so $q$ respects the triple $(1,C,1/c)$.
Anthropic Proposition~A.5 (pp.~20--21), whose upper-functional argument was
targeted checked in this run, supplies $\varphi\leq q$ respecting that triple.
Consequently $1\leq\varphi(C)/c\leq p(C)/c=1$.
\end{proof}

\subsection{A four-point description conditional on the square}

For a pseudometric $d$, write the weighted quadruple expression as
\begin{align}
 \mathcal Q_{x,y,z,w}(s,t)={}&
 (1-t)(1-s)d(x,y)^2+t(1-s)d(y,z)^2\nonumber\\
 &+tsd(z,w)^2+(1-t)s d(w,x)^2\nonumber\\
 &-t(1-t)d(x,z)^2-s(1-s)d(y,w)^2.
 \label{pg:weighted-quad}
\end{align}
Every CAT(0) metric has $\mathcal Q\geq0$ for all $s,t\in[0,1]$.
The exact finite-metric characterization due to Toyoda (2020),
Definition~1.1 and Theorem~1.3, pp.~114--115, says that on at most
five points these inequalities are also sufficient for CAT(0)
embeddability. For four points this is the earlier Gromov theorem,
recorded there as Theorem~1.7, p.~116. We need the following more
specific realization in order to apply the pendant operation to prices.

\begin{theorem}[PG-FOUR-CORE]
\label{pg:four-core}
A pseudometric $d$ on four points is isometrically embeddable in a
CAT(0) space if and only if there are $x_1,\ldots,x_4\in\mathbb R^2$
and $a_1,\ldots,a_4\geq0$ such that
\[
 d(i,j)=\|x_i-x_j\|+a_i+a_j\qquad(i\ne j).
\]
Equivalently, this representation holds if and only if
\eqref{pg:weighted-quad} is nonnegative for all choices of its arguments.
\end{theorem}
\begin{proof}
A representation of the stated form embeds in a Euclidean plane with
an interval of length $a_i$ attached at $x_i$ for each $i$; the selected
points are the outer endpoints. Repeated use of the CAT(0) gluing theorem
along a single point makes this a CAT(0) space (Bridson and Haefliger,
1999, Theorem~II.11.1, p.~347). Thus it satisfies the weighted quadruple
inequalities. We prove directly the converse needed here.

If the pseudometric identifies two points, its metric quotient has at
most three points and embeds in the Euclidean plane. Pulling this
embedding back gives the desired representation with every $a_i=0$.
Hence assume all distances between distinct points are positive.
On $H=\{z\in\mathbb R^4:\sum_i z_i=0\}$ define
\[
 Q_d(z)=-\frac12\sum_{i,j}z_i z_jd(i,j)^2.
\]
The elementary Gram criterion says that $Q_d$ is positive semidefinite
if and only if $d$ embeds in $\mathbb R^3$. Indeed, fixing a base point
$p$ gives the Gram matrix
$G_{ij}=(d(p,i)^2+d(p,j)^2-d(i,j)^2)/2$ on the other three indices,
and $Q_d(z)=\sum_{i,j\ne p}z_i z_jG_{ij}$ after substituting
$z_p=-\sum_{i\ne p}z_i$. A positive semidefinite $G$ is the Gram
matrix of vectors giving the required distances by this same identity.

First suppose that $d$ embeds in $\mathbb R^3$. If its Gram matrix
has rank at most two, nothing is required. Otherwise fix one point $p$
and put
\[
 L=\frac12\min_{i,j\ne p,\ i\ne j}
       \bigl(d(p,i)+d(p,j)-d(i,j)\bigr)>0.
\]
For $0\leq u\leq L$, subtract $u$ from the three distances incident
to $p$, leaving the other distances unchanged. The resulting distances
remain a pseudometric: the only changing triangle defects are the
ones with middle vertex $p$, and those decrease by $2u$. Also
$L\leq\min_{i\ne p}d(p,i)$, by the original triangle inequalities.
At $u=L$ one triangle is degenerate, so its Gram matrix is not positive
definite. Start with the positive definite Gram matrix at $u=0$ and
increase $u$ to the first value $u_0$ where positive definiteness fails.
Continuity makes the matrix at $u_0$ positive semidefinite and singular.
The new pseudometric therefore embeds in $\mathbb R^2$. It is the
required core, with $a_p=u_0$ and the other pendant lengths zero.

Now suppose that $Q_d$ is not positive semidefinite, and choose $z\in H$
with $Q_d(z)<0$. No coordinate of $z$ can be zero, since every
three-point metric embeds in the Euclidean plane and hence has a
positive semidefinite Gram form. Nor can $z$ have two positive and two
negative coordinates: after dividing by the common total positive
and negative mass, $Q_d(z)$ is a positive scalar multiple of one
of the expressions \eqref{pg:weighted-quad}. Consequently, after
changing the sign of $z$, exactly one coordinate, indexed by $p$,
is negative and the other three are positive.

Relabel those three points as $1,2,3$, and put $r_i=d(p,i)>0$.
Let $\alpha_{ij}\in[0,\pi]$ be their comparison angles at $p$, so
\[
 \cos\alpha_{ij}
 =\frac{r_i^2+r_j^2-d(i,j)^2}{2r_i r_j}.
\]
We claim that
$\alpha_{12}+\alpha_{23}+\alpha_{31}>2\pi$.
Otherwise choose $\beta_{ij}\in[\alpha_{ij},\pi]$ with their sum
equal to $2\pi$. This is possible because the initial sum is at most
$2\pi$ and the sum of the three upper bounds is $3\pi$.
There are planar unit vectors $v_1,v_2,v_3$ with these pairwise angles:
place them successively around the circle with gaps
$\beta_{12},\beta_{23},\beta_{31}$. Their Gram matrix has off-diagonal
entries $\cos\beta_{ij}\leq\cos\alpha_{ij}$. Therefore, for the
positive coefficients $b_i=z_i r_i$,
\[
 Q_d(z)=\sum_{i,j=1}^3 b_i b_j\cos\alpha_{ij}
 \geq\left\|\sum_{i=1}^3 b_i v_i\right\|^2\geq0,
\]
a contradiction. The diagonal entries in this formula are one.

For $0\leq t<1$, try pendant lengths $a_p=0$ and $a_i=t r_i$.
The candidate core has radial distances $(1-t)r_i$ and pair distances
$d(i,j)-t(r_i+r_j)$. Define
\[
 t_* =\min_{1\leq i<j\leq3}
       \frac{d(i,j)-|r_i-r_j|}{2\min(r_i,r_j)}\in[0,1].
\]
The strict angle-sum inequality implies $t_*>0$.
For $0\leq t\leq t_*$ with $t<1$, each radial triangle is valid:
its opposite side, divided by $1-t$, lies in
$[|r_i-r_j|,r_i+r_j]$. Moreover
\[
 \frac{d(i,j)-t(r_i+r_j)}{1-t}
\]
is nonincreasing in $t$. Thus all three comparison angles are
continuous and nonincreasing. If $t_*<1$, at least one angle is zero
at $t_*$, and their sum is then at most $2\pi$. By continuity there
is $t_0\in(0,t_*]$ where their sum is exactly $2\pi$.
Place unit vectors in the plane with these three successive angles,
and give the corresponding vectors the lengths $(1-t_0)r_i$.
The cosine law realizes all the candidate core distances exactly.
Together with the origin representing $p$, this is the required
planar core. No triangle inequality among the three noncentral
points was assumed along this deformation; at $t_0$ it follows from
the explicit Euclidean realization.

If $t_*=1$, every pair has $d(i,j)=r_i+r_j$. The original metric
is a star, and the zero core with pendant lengths $a_i=r_i$, $a_p=0$
gives the required representation.
\end{proof}

\begin{corollary}[PG-FOUR-CONDITIONAL]
\label{pg:four-conditional}
Fix a field and assume the checked price-space chain used throughout
this section. Exactly one of the following two descriptions applies
to its four-party prices.
\begin{enumerate}
\item If the Euclidean square is absent, the prices are exactly the
      tree pseudometrics.
\item If the Euclidean square is present, the prices are exactly the
      pseudometrics satisfying all weighted quadruple inequalities
      \eqref{pg:weighted-quad}, equivalently all four-point metrics
      isometrically embeddable in a CAT(0) space.
\end{enumerate}
\end{corollary}
\begin{proof}
The absent-square case follows from Theorems~\ref{pg:square-gate} and
\ref{pg:tree-all}. In the present-square case every finite planar
Euclidean metric is a price by Theorem~\ref{pg:square-gate}.
Theorem~\ref{pg:four-core} and repeated pendant addition
(Lemma~\ref{pg:pendant}) give every metric in the stated class as
a price. The price-space theorem supplies the reverse inclusion.
\end{proof}

\begin{proposition}[PG-FIVE-CORE-FAIL: the four-point construction stops]
\label{pg:five-core-fail}
There is a five-point Euclidean metric, hence a five-point CAT(0)
metric, which cannot be written as a planar Euclidean metric plus
nonnegative singleton pendant lengths.
\end{proposition}
\begin{proof}
In $\mathbb R^3$ take
\[
 v_1=(1,0,0),\quad v_2=(0,1,0),\quad
 v_3=(-1,0,0),\quad v_4=(0,-1,0),\quad q=(1,0,1).
\]
The first four points form a square with sides $\sqrt2$ and diagonals
$2$. The distances from $q$ to them are respectively
\[
 (r_1,r_2,r_3,r_4)=(1,\sqrt3,\sqrt5,\sqrt3).
\]
Suppose that a planar core with nonnegative pendant lengths realizes
this metric. Realize that representation by attaching intervals to
the plane. The four base points span a flat square in this CAT(0)
space, by the flat quadrilateral theorem (Bridson and Haefliger, 1999,
Theorem~II.2.11, p.~181; equivalently the square equality case recorded
in Anthropic Lemma~C.4(f), p.~29). Every base pendant length must be
zero: at the endpoint of a positive pendant the geodesics to both
adjacent square vertices share that pendant, whereas in a flat square
these two geodesics intersect only at the vertex. A planar rigid motion
therefore puts the four base core points at $(1,0),(0,1),(-1,0),(0,-1)$.

Let $x=(x_1,x_2)$ be the fifth core point and let $a\geq0$ be its
pendant length. Its core distances are $r_i-a$. Squared distances to
the four square vertices satisfy the parallelogram relation
\[
 (r_1-a)^2+(r_3-a)^2=(r_2-a)^2+(r_4-a)^2.
\]
Since $r_1^2+r_3^2=r_2^2+r_4^2=6$, this implies
\[
 a(1+\sqrt5-2\sqrt3)=0.
\]
The parenthesized coefficient is nonzero: equality would imply
$6+2\sqrt5=12$, hence $\sqrt5=3$. Thus $a=0$.
The equality $r_2=r_4$ now gives $x_2=0$, and
$r_3^2-r_1^2=4$ gives $x_1=1$. This makes the distance from $x$ to
the first square vertex zero, contradicting $r_1=1$.
\end{proof}

This obstruction concerns one attempted representation, not price
membership. The five-point metric in Proposition~\ref{pg:five-core-fail}
has not been proved to be a price or a nonprice. Toyoda's weighted
quadruple characterization still gives the exact CAT(0) outer bound
on five-party prices, but the square-present converse requires another
argument.

\begin{proposition}[PG-EDGE-RIGIDITY: a saturated normalized wire]
\label{pg:edge-rigidity}
Let $\psi$ be a normalized ordinary spectral point on a finite party set
$I$. If $d_\psi(a,b)=1$ for distinct $a,b\in I$, then its entire wire
pseudometric is induced by a subset of the unit interval. More precisely,
\[
 d_\psi(i,j)=|d_\psi(a,i)-d_\psi(a,j)|\qquad(i,j\in I).
\]
\end{proposition}
\begin{proof}
For a third party $i$, restrict the functional to tensors supported on
$J=\{a,b,i\}$, with trivial factors elsewhere. This is an enlarged
spectral point. Its logarithmic value $\beta$ on the binary full unit
tensor on $J$ satisfies $\beta\leq1$, because that tensor is a restriction
of the full unit tensor on $I$. It also satisfies $\beta\geq1$, because
it restricts to the saturated wire on $\{a,b\}$. Thus $\beta=1$, and the
normalization lemma makes this restricted functional an ordinary
normalized three-party spectral point.

Alman, Li and Pratt (2026), arXiv:2604.01386v2, Proposition~8.1,
\S8.1, pp.~47--48, proves that a three-party normalized spectral point
whose exponent on one wire equals one has the sum of its other two wire
exponents equal to one. Its proof applies Sch\"onhage's direct-sum
degeneration, quoted there as Lemma~8.2, and works over every field.
Consequently $d_\psi(a,i)+d_\psi(i,b)=1$ for every $i$.

Embed the wire metric into a CAT(0) price space. Equality in the triangle
inequality puts every image point on the unique geodesic segment from
$a$ to $b$. Distance on this segment is the difference of the distances
from $a$, giving the formula.
\end{proof}

The condition in this proposition is a normalization condition, not a
condition one may impose freely by rescaling a price. A nontrivial enlarged
spectral point only permits powers that preserve its cover inequality;
rescaling its metric to have diameter one does not by itself produce an
ordinary normalized spectral point. Thus the proposition does not settle
the square question.

This is a sufficient extension construction for the particular functionals
used below. It does not assert that an arbitrary multiplicative function on
wire networks extends to all tensors without further hypotheses.

\begin{proposition}[PG-STAR: star metrics over every field]
\label{pg:star}
For every field $k$, every finite set $I$, and every family $a_i\geq0$, the
pseudometric
\[
 d(i,j)=a_i+a_j\quad(i\ne j),\qquad d(i,i)=0
\]
is a price over $k$.
\end{proposition}
\begin{proof}
The zero case is immediate. Otherwise write $a=\sum_i a_i>0$ and
$\theta_i=a_i/a$. The Strassen upper support functional $p=\zeta^\theta$
is normalized, monotone under restrictions, additive under direct sums and
submultiplicative; these are the arbitrary-field assertions of Christandl,
Vrana and Zuiddam, arXiv:1709.07851v3, Theorem~2.4, p.~11.
A cover is a restriction of the direct sum of its covering tensors, so $p$
is an upper functional in the present sense.

Every graph tensor $G=\prod_{ij}\langle q_{ij}\rangle_{ij}$ is tight,
hence oblique. For completeness, orient its edges, assign distinct sufficiently
separated integer place values to their labels, and assign opposite signs at
the two ends of each edge. The sum of these signed place values at a vertex
is injective on its local index set, and their total over all vertices is zero
on the support. Ordering the local indices by these injective values makes
the support an antichain. Independent uniform edge labels give uniform local
marginals. The exact entropy formula for oblique tensors, CVZ
Theorem~2.19, pp.~15--16, therefore gives
\[
 p(G)=\prod_{ij}q_{ij}^{\theta_i+\theta_j}.
\]
Take $C$ to be the product of one binary wire for every unordered pair of
distinct parties. Lemma~\ref{pg:saturation} gives $\varphi\leq p$ with
$\varphi(C)=p(C)$. Multiplicativity and the individual upper bounds force
equality on each of the binary wires. Anthropic Lemma~A.2 (p.~19) then gives
$d_\varphi(i,j)=\theta_i+\theta_j$. Multiplication by $a$ gives $d$.
\end{proof}

\begin{theorem}[PG-THREE: exact three-party classification]
\label{pg:three}
On a party set of cardinality at most three, the prices over every field are
exactly the pseudometrics. In particular, the same answer holds over
$\mathbb C$ and $\mathbb F_2$.
\end{theorem}
\begin{proof}
Prices are pseudometrics by Anthropic Lemma~A.2 and closure under scaling and
limits. Conversely, on three parties put
\[
 a_1=\tfrac12(d_{12}+d_{13}-d_{23}),\quad
 a_2=\tfrac12(d_{12}+d_{23}-d_{13}),\quad
 a_3=\tfrac12(d_{13}+d_{23}-d_{12}).
\]
The triangle inequalities give $a_i\geq0$, and $d_{ij}=a_i+a_j$ for
$i\ne j$. Proposition~\ref{pg:star} applies. Fewer parties follow either
directly or by adding zero-distance duplicates.
\end{proof}

\subsection{Tree metrics and the four-party obstruction}

\begin{proposition}[PG-TREE-C: tree metrics over the complex numbers]
\label{pg:tree}
Every metric induced on a finite set of vertices of a weighted tree with
nonnegative edge lengths is a price over $\mathbb C$.
\end{proposition}
\begin{proof}
Delete edges which separate no pair of the selected parties. If the total
remaining edge length is zero the assertion is immediate. Otherwise let
$\theta_e$ be each edge length divided by their total $L$. Each edge
defines a bipartition $b_e$ of the party set. Repeated bipartitions can be
combined, and these splits are pairwise compatible, or noncrossing.
The upper quantum functional $p=F^\theta$ is an upper functional by CVZ,
Theorem~3.5, p.~18, including its cover inequality obtained from restriction
monotonicity and direct-sum subadditivity.

For a single wire between $i$ and $j$, its value is
\[
 p(\langle q\rangle_{ij})=q^{\sum_{e:\,b_e\text{ separates }i,j}\theta_e}.
\]
This also follows directly from CVZ Definition~3.3: all nontrivial bipartition
projectors act on the same two nontrivial tensor legs, whereas trivial
bipartitions contribute zero entropy. Their common entropy supremum is
$\log_2 q$. For a product of wires, submultiplicativity gives the product
of these values as an upper bound. The normalized product of maximally
entangled wire states has, across each split, a uniform nonzero reduced
spectrum of rank equal to the product of dimensions of the wires crossing
that split. The definition of the lower quantum functional, together with
CVZ Theorem~3.24 (p.~25), gives the matching lower bound. Thus $p$ has
exactly these multiplicative values on every graph tensor.

Apply Lemma~\ref{pg:saturation} to the complete product of binary wires.
As in Proposition~\ref{pg:star}, the resulting genuine spectral point has
wire metric $L^{-1}$ times the desired tree metric.
\end{proof}

The argument does not assert that $F^\theta$ itself is a spectral point on
all tensors. Botero et al. (2026), arXiv:2604.18283v1, Theorem~3.14 and
Proposition~4.3, distinguish the failure of universality of nonsingleton
lower quantum functionals from the existence of extensions of their values
on a suitable subsemiring. This distinction matters here.

\begin{proposition}[PG-NONCONVEX: prices need not be closed under addition]
\label{pg:nonconvex}
Assuming the checked price-space theorem, the set of prices on four parties
is not a convex cone over any field.
\end{proposition}
\begin{proof}
Label the parties $00,10,11,01$. The two coordinate cut pseudometrics are
prices: each is the wire metric of a flattening rank. Their sum has cyclic
side lengths $1$ and both diagonals $2$. If it were a price, its embedding
in a CAT(0) price space, Anthropic Proposition~C.3 (p.~28), would give
the roundness inequality
$2^2+2^2\leq1^2+1^2+1^2+1^2$, which is false. Therefore the sum is
not a price. Positive rescaling then also rules out convexity.
\end{proof}

In particular, a positive combination of cut metrics cannot be promoted to
a price merely by using a Crofton representation of a Euclidean metric.
The exact four- and five-party classifications, including the Euclidean
square, are not settled by the preceding inclusions and obstructions.

\subsection{Euclidean configurations do not settle exclusion of the cases}

\begin{proposition}[PG-HILB2: sharp Euclidean rank-step inequality]
\label{pg:hilbert}
For every map $f:\operatorname{Mat}_2(\mathbb F_2)\to H$ into a real or
complex Hilbert space, define $D_f(r)^2$ by averaging
$\|f(A+W)-f(A)\|^2$ over all $A$ and all matrices $W$ of rank $r$.
Then
\[
 D_f(2)^2\leq\frac32 D_f(1)^2.
\]
The factor $3/2$ is optimal, already for maps into $\mathbb R$.
\end{proposition}
\begin{proof}
Let $G=\operatorname{Mat}_2(\mathbb F_2)$, and for $B\in G$ let
$\chi_B(A)=(-1)^{\sum_{ij}B_{ij}A_{ij}}$. For the two uniform rank-step
operators, the character eigenvalues depend only on $\operatorname{rank}B$.
Indeed, invertible left and right multiplication permute each rank class
and transform nonzero $B$ transitively within its rank class. There are nine
rank-one matrices and six invertible matrices. Directly summing the sixteen
characters, or using the explicit enumeration below, gives
\[
\begin{array}{c|cc}
\operatorname{rank}B&\kappa_1(B)&\kappa_2(B)\\\hline
0&1&1\\
1&1/9&-1/3\\
2&-1/3&1/3
\end{array}
\]
For a hand check of the rank-one column, write $W=uv^T$ with
$u,v\in\mathbb F_2^2\setminus\{0\}$. For rank-one $B$, precisely one
nonzero $v$ lies in its kernel; for rank-two $B$, none does. Summing over
$u$ gives respectively $3$ or $-1$, yielding $1/9$ and $-1/3$.
The sum of a nontrivial character over $G$ is zero; subtracting its value
at zero and its nine rank-one values yields the rank-two column.

Expand the finitely many values of $f$ by Fourier characters. Parseval's
identity, applied in their finite-dimensional Hilbert span, gives
\[
 D_f(r)^2=2\sum_{B\in G}(1-\kappa_r(B))\|\widehat f(B)\|^2.
\]
The ratio of the two coefficients is $3/2$ at dual rank one and $1/2$
at dual rank two. Termwise comparison proves the inequality. For sharpness
take $f=\chi_B$ with $\operatorname{rank}B=1$.
\end{proof}

Thus every Euclidean pseudometric on this matrix set obeys a strict version
of the inequality in Anthropic Definition~D.1, with $s=\sqrt{3/2}<\sqrt2$.
The existence of Euclidean prices cannot, by itself, refute contraction.
This does not prove that contraction holds for all prices, and it does not
prove that the two hypotheses actually coexist for either field.

The exact certificate program \path{programs/pg_rank_step_fourier.py}
enumerates all matrices for $n=2,3,4$ and prints rational character averages
and energy ratios. Its recorded output is
\path{outputs/pg_rank_step_fourier.tsv}. The larger finite cases are
computational observations, not substitutes for an all-$n$ theorem.

\subsection{The sharp rank metric is a CAT(0) submetric}

\begin{proposition}[PG-BUILDING: CAT(0) does not force strict contraction]
\label{pg:building}
For every positive integer $n$, the metric
\[
 d(A,B)=\sqrt{\operatorname{rank}(A-B)}
 \quad\text{on }\operatorname{Mat}_n(\mathbb F_2)
\]
embeds isometrically in a complete CAT(0) space. Consequently no strict
inequality $D(n)^2\leq s^2D(1)^2$ with $s^2<n$ holds uniformly for maps
of this matrix set into arbitrary CAT(0) spaces.
\end{proposition}
\begin{proof}
Let $V=\mathbb F_2^{2n}$, and let $\Delta$ be its type-$A_{2n-1}$
spherical building: vertices are nonzero proper subspaces and simplices
are flags. An apartment is given by a basis of $V$. We use the standard
spherical Coxeter metric. The building construction is Bridson and
Haefliger (1999), Example~II.10A.3(5), pp.~343--344; apartments embed
isometrically and the building is CAT(1) by Theorem~II.10A.4(i)--(ii),
p.~344. Its Euclidean cone is CAT(0) by Theorem~II.3.14, p.~188.
Because the underlying field is finite, the building is a finite complex;
its cone is complete.

For each $n$-dimensional subspace $U\subset V$, take the point $p_U$
on its vertex ray at radius $\sqrt n/2$. Any two such subspaces admit
a common apartment, obtained by extending a basis of their intersection
first to bases of each and then of $V$. In that apartment, identify the
cone with the hyperplane $\sum_{j=1}^{2n}x_j=0$ in $\mathbb R^{2n}$.
If $S_U$ is the $n$-element coordinate subset spanning $U$, then
\[
 p_U=\frac1{\sqrt2}\left(\mathbf1_{S_U}-\frac12\mathbf1\right).
\]
It follows in this common apartment that
\[
 d(p_U,p_W)^2
 =\tfrac12\lvert S_U\mathbin\triangle S_W\rvert
 =n-\dim(U\cap W).
\]
For $A\in\operatorname{Mat}_n(\mathbb F_2)$ use its graph
$U_A=\{(x,Ax):x\in\mathbb F_2^n\}$. Then
$\dim(U_A\cap U_B)=\dim\ker(A-B)=n-\operatorname{rank}(A-B)$,
so $A\mapsto p_{U_A}$ is the claimed isometric embedding.
For this metric $D(1)^2=1$ and $D(n)^2=n$ identically.
\end{proof}

No assertion that this building metric is a \emph{price} is made. That
membership is precisely additional algebraic information which the CAT(0)
embedding theorem does not provide. For $n\geq2$ this example is not a
Hilbert submetric: restricting to a $2\times2$ matrix block would contradict
Proposition~\ref{pg:hilbert}.

\subsection{A square is the first obstruction to tree geometry}

\begin{theorem}[PG-SQUARE-GATE: a non-tree price forces all planar prices]
\label{pg:square-gate}
Fix a field and assume the checked price-space theorem. The following
conditions are equivalent:
\begin{enumerate}
\item the Euclidean unit square is a price;
\item every finite subset of the Euclidean plane is a price;
\item some finite price does not embed isometrically into an $\mathbb R$-tree.
\end{enumerate}
Consequently, over every field, failure of the square to be a price is
equivalent to the assertion that, on every finite party set, the prices
are exactly the pseudometrics induced by weighted trees.
\end{theorem}
\begin{proof}
The implication (2)$\Rightarrow$(1) is immediate. A square is not a tree
metric: its three opposite-pair sums are $2,2,2\sqrt2$, contrary to the
four-point condition for tree metrics (Buneman, 1974, pp.~48--49, and
Theorem~2, p.~49). Thus (1)$\Rightarrow$(3).

Suppose (3) holds, and embed this finite price in a complete CAT(0) price
space $X$. Then $X$ is not an $\mathbb R$-tree. There is therefore a
geodesic triangle that is not a tripod; this is the tree characterization
in Bridson--Haefliger, Example~II.1.15(5), p.~167. Write one of its sides
as a unit-speed geodesic $\gamma:[0,L]\to X$, and let $z$ be the third
vertex. The function $h(t)=d(z,\gamma(t))$ is convex and $1$-Lipschitz
(Bridson--Haefliger, Proposition~II.2.2, p.~176).

There exists a differentiability point $t\in(0,L)$ with
$-1<h'(t)<1$. Indeed, otherwise the monotone derivative of this convex
absolutely continuous function takes only the values $-1,1$ almost
everywhere. Then $h$ has the form $c+|t-t_0|$ on the interval, including
the possible case of a single affine branch. Equality in the triangle
inequality at its minimum makes both remaining triangle sides pass
through $\gamma(t_0)$, or through an endpoint in the single-branch case.
Uniqueness of CAT(0) geodesics makes the triangle a tripod, a contradiction.

Set $p=\gamma(t)$. By the first variation formula,
Bridson--Haefliger, Corollary~II.3.6, p.~185, the forward direction of
$\gamma$ and the direction of $[p,z]$ have angle
$\vartheta=\arccos(-h'(t))\in(0,\pi)$. The completion of the space
of directions at $p$ is CAT(1), by their Theorem~II.3.19, p.~191. The
geodesic between these two directions is an isometric interval of length
$\vartheta$. Its Euclidean cone is an isometrically embedded planar
sector with positive angle in the completed tangent cone at $p$.

Every finite metric in this tangent cone is a price. To see this directly,
first approximate its finitely many directions by actual geodesic
directions in $X$. For a finite set of actual directions and radii
$r_i$, choose points at distances $\varepsilon r_i$ from $p$ along
the corresponding geodesics. As $\varepsilon\downarrow0$, their
distances divided by $\varepsilon$ converge to the cone distances by
the definition of Alexandrov angle and the Euclidean cosine formula.
These are pullbacks of rescaled price metrics; completion of the
directions is handled by another limit. Closure of prices under finite
pullbacks, positive scaling and limits proves the assertion.

Every finite planar configuration can be translated and scaled into an
open disk inside the sector. Pulling back and scaling its metric proves
(2). Finally, Proposition~\ref{pg:tree-all} below supplies every
finite tree metric, and the equivalence just proved supplies the reverse
inclusion if the square is absent.
\end{proof}

This theorem reduces the square question to the existence of a single
non-tree price; it does not establish either side of that alternative.

\subsection{One cross-polytope tests an entire Euclidean dimension}

\begin{theorem}[PG-CROSSPOLY-GATE]
\label{pg:crosspoly-gate}
Fix a field and an integer $D\geq2$. The following are equivalent:
\begin{enumerate}
\item the Euclidean metric on $\{\pm e_1,\ldots,\pm e_D\}$ is a price;
\item that metric together with its center $0$ is a price;
\item every finite subset of $\mathbb R^D$ with its Euclidean metric is a price;
\item no strict shortening exists in the specified dimension $D$, for any
width $m\geq1$.
\end{enumerate}
For $D=2$, these conditions are equivalent to the Euclidean unit square
being a price. Consequently a strict planar shortening exists if and only
if every finite price over the field is a tree pseudometric.
\end{theorem}
\begin{proof}
The implications (3)$\Rightarrow$(2)$\Rightarrow$(1) are immediate.
Assume (1), and embed that finite metric isometrically into a complete
CAT(0) price space $X$, using Anthropic Proposition~C.3.
The Kirszbraun extension theorem for curvature bounds gives a
$1$-Lipschitz extension $f:\mathbb R^D\to X$ of the given vertex map.
The precise version used is Alexander, Kapovitch and Petrunin (2011),
\emph{Alexandrov meets Kirszbraun}, \S5, Theorem~5.4, p.~11 of the
author-posted PDF (Theorem~5.1, p.~9, gives the general curvature version).
It applies to a complete CBB(0) domain and a complete CAT(0) target;
$\mathbb R^D$ and $X$ have exactly these properties. There is no properness
or finite-dimensionality requirement on $X$. This is the nonpositive
curvature extension theorem originally due to Lang and Schroeder (1997).

The restriction of $f$ to the unit ball preserves every coordinate
antipodal distance. Lemma~\ref{lem:rigid-ball} therefore makes $f$
isometric on the concentric ball of radius $2^{-(D-1)/2}$. Every finite
Euclidean $D$-dimensional configuration can be translated and scaled
into that ball. Its induced metric is a pullback of the metric of $X$,
hence a price; rescaling proves (3).

The profile-compression argument in Theorem~\ref{thm:selector-price}
shows that (3) excludes a strict shortening whose layouts lie in the
same $D$-dimensional ball, so (3)$\Rightarrow$(4). Conversely, the
duality construction of Proposition~\ref{prop:saturated-diameters} and
Corollary~\ref{cor:ball-price} uses absence of shortening only in its
fixed dimension $D$. Under (4), it supplies a price on the whole unit
ball, bounded by Euclidean distance, and with all coordinate antipodes
at distance two. Applying Lemma~\ref{lem:rigid-ball} gives (3).

At $D=2$, the four vertices are a rotated and rescaled Euclidean unit
square. The final assertion now follows from
Theorem~\ref{pg:square-gate} and Theorem~\ref{pg:tree-all}: absence of
the square is equivalent to every finite price being a tree metric,
whereas the equivalence just proved makes absence of the square
equivalent to existence of a strict shortening in dimension two.
\end{proof}

The assertion that a strict planar shortening exists is still an
undecided alternative. This theorem gives a finite metric test for it;
it does not construct the shortening or assert that the test fails.

\subsection{Every finite tree is a price over every field}

\begin{lemma}[PG-PENDANT: adding a pendant length]
\label{pg:pendant}
For a price $p$ on a finite party set $I$, $i\in I$ and $a\geq0$, the
pseudometric $p+a\delta_i$ is a price, where
$\delta_i(j,k)=1$ when exactly one of $j,k$ equals $i$, and is zero otherwise.
\end{lemma}
\begin{proof}
First suppose the field is infinite and $p=c\,d_\psi$ for a normalized
ordinary spectral point $\psi$, with $c>0$. Use the complete graph on $I$.
Wigderson and Zuiddam (2025), Theorem~12.23, \S12.7, p.~121, states that
the joint spectrum of its binary edge tensors is log-star-convex about
each singleton flattening rank. On these wires, the base-two logarithm of
the flattening rank at $i$ is exactly $\delta_i$. Therefore there is a
normalized spectral point whose wire metric is
\[
 \frac{c}{c+a}d_\psi+\frac{a}{c+a}\delta_i.
\]
Rescaling by $c+a$ gives $p+a\delta_i$. If $c=0$, the desired metric is
a rescaled flattening metric directly. For an arbitrary price, use its
defining approximation by scaled spectral wire metrics and the
normalization result for enlarged spectral points; pass to a limit.

For a finite field $F$, apply the infinite-field argument over $F(t)$.
The same-characteristic invariance of price cones proved in EX-15 permits
both passage to $F(t)$ and descent to $F$. Thus the assertion holds over
every field. The one-party case is vacuous.
\end{proof}

\begin{theorem}[PG-TREE-ALL: tree inclusion over every field]
\label{pg:tree-all}
Every finite pseudometric induced by a weighted tree with nonnegative edge
lengths is a price over every field.
\end{theorem}
\begin{proof}
Put each selected point at a leaf by attaching an edge of length zero if
necessary. Delete unselected leaves, suppress unselected degree-two
vertices by summing their incident edge lengths, and resolve vertices of
degree greater than three into binary trees whose new edges have length
zero. These operations preserve the selected metric. It is thus enough
to consider a binary tree with the selected parties as leaves.

Induct on the number of leaves. With at most two leaves a rescaled cut
metric gives the assertion. With at least three leaves choose a cherry,
that is, two leaves $i,j$ with a common adjacent vertex $v$. Remove those
two pendant edges and use $v$ as the replacement selected leaf. By
induction the resulting tree metric is a price. Pull that metric back
under the map identifying $i$ and $j$ with $v$. The pullback gives zero
distance between $i,j$ and the correct old distances to $v$. If the two
removed edges have lengths $a_i,a_j$, add $a_i\delta_i$ and
$a_j\delta_j$, successively using Lemma~\ref{pg:pendant}. The resulting
distances are exactly those of the original tree, completing the induction.
\end{proof}

This proof supplies the finite-characteristic inclusion without a quantum
functional. It does not assert closure of the price set under addition
of arbitrary cuts: only singleton cuts have been added. The incompatibility
of arbitrary cut addition with CAT(0) remains as shown above.

\subsection{An explicit strict rank-step inequality for tree metrics}

For $f:\operatorname{Mat}_n(\mathbb F_2)\to X$, write
\[
 D_f(r)^2=\mathbb E_{A,Y:\operatorname{rank}Y=r}
 d(f(A),f(A+Y))^2,
\]
where $A$ and $Y$ are independent and uniform on their indicated finite
sets.

\begin{theorem}[PG-TREE-CONTRACTION]
\label{pg:tree-contraction}
Every map from $\operatorname{Mat}_n(\mathbb F_2)$ to an $\mathbb R$-tree
satisfies
\[
 D_f(n)^2\leq4(1-2^{-n})^2D_f(1)^2.
\]
In particular, for $n=4$,
$D_f(4)\leq(15/8)D_f(1)$, and $(15/8)^2<4$.
Consequently, if the Euclidean square is not a price over a field, that
field satisfies Anthropic's contracting hypothesis with the explicit
parameters $n=4$ and $s=15/8$.
\end{theorem}
\begin{proof}
The image of $f$ lies in the finite compact subtree spanned by its finitely
many values. Let $b$ minimize the continuous function
$\mathbb E_A d(b,f(A))^2$ there, and put $r_A=d(b,f(A))$. The case where
every $r_A=0$ is immediate. For a component $B$ of the tree minus $b$
which meets the image, set
\[
 a_B=\mathbb E_A[r_A\mathbf1_{f(A)\in B}],\qquad
 a=\sum_Ba_B=\mathbb E_A r_A.
\]
Travel a sufficiently small distance $t$ from $b$ into the branch $B$.
Because the image is finite, every path from $b$ to an image point in $B$
shares a positive initial segment in that direction. Distances to those
points decrease by $t$, and distances to all other noncentral image points
increase by $t$. Central image points contribute only $t^2$. The right
derivative of the objective at $t=0$ is therefore $2(a-2a_B)$. Minimality
implies $a_B\leq a/2$ for every branch.

The elementary polygon criterion now gives unit vectors $u_B$ in
$\mathbb R^2$ with $\sum_Ba_Bu_B=0$. Here the criterion says that
nonnegative lengths form a closed planar polygon if and only if the
largest length is at most the sum of the others. One direct verification
is to vary the angles of a chain of the remaining segments: its possible
endpoint distances form the interval between
$\max(0,2\max_{C\ne B}a_C-\sum_{C\ne B}a_C)$ and
$\sum_{C\ne B}a_C$; the largest length $a_B$ lies in this interval.
This interval assertion follows by induction from the triangle inequality
and the continuous rotation of the last segment. Close the chain with
the segment of length $a_B$ and divide its edge vectors by their lengths.

Define $v_A=r_Au_B$ when $f(A)\in B$, and $v_A=0$ when $f(A)=b$.
Then $\mathbb E_Av_A=0$ and $\|v_A\|=r_A$. If two image points are in
different branches, their tree distance is the sum of their radii, which
dominates the Euclidean distance of their vectors. If they are in the
same branch, their vectors have distance $|r_A-r_C|$, which is at most
their tree distance by the reverse triangle inequality. Thus
\[
 \|v_A-v_C\|\leq d(f(A),f(C))
 \quad\hbox{for all }A,C.
\]

We compute the spectral gap of the rank-one random walk on the additive
group $G=\operatorname{Mat}_n(\mathbb F_2)$. Its characters are
$\chi_B(A)=(-1)^{\sum_{ij}B_{ij}A_{ij}}$. A rank-one matrix is uniquely
$uv^{\mathsf T}$ for nonzero $u,v\in\mathbb F_2^n$. If
$\operatorname{rank}B=r$, summing first over $v$ gives $2^n-1$ when
$u^{\mathsf T}B=0$, and $-1$ otherwise. There are
$2^{n-r}-1$ nonzero $u$ of the first kind. Hence the eigenvalue is
\[
 \kappa_1(r)
 =\frac{2^{2n-r}-2^{n+1}+1}{(2^n-1)^2}.
\]
This decreases strictly with $r$, so the largest nonconstant eigenvalue
is $\kappa_1(1)$ and the spectral gap is
\[
 \gamma_n=1-\kappa_1(1)
 =\frac{2^{2n-1}}{(2^n-1)^2}.
\]
Parseval's identity, applied to the two coordinates of the mean-zero map
$v$, gives
\[
 D_f(1)^2\geq\mathbb E_{A,Y:\operatorname{rank}Y=1}
 \|v_A-v_{A+Y}\|^2
 \geq2\gamma_n\mathbb E_A r_A^2.
\]
For every full-rank step, the tree triangle inequality and
$(r_A+r_{A+Y})^2\leq2r_A^2+2r_{A+Y}^2$ yield
$D_f(n)^2\leq4\mathbb E_A r_A^2$. Combining gives
\[
 \frac{D_f(n)^2}{D_f(1)^2}
 \leq\frac{2}{\gamma_n}
 =4(1-2^{-n})^2,
\]
with the constant-map case handled separately. For $n=4$ this is
$225/64<4$. If the square is absent, Theorem~\ref{pg:square-gate}
makes every finite price a tree metric, so this uniform bound applies to
all prices on $\operatorname{Mat}_4(\mathbb F_2)$.
\end{proof}

\subsection{A quantitative invariant gate for contraction}

Let $X$ be complete CAT(0), let $\nu=\sum_i t_i\delta_{x_i}$ be a
nonconstant finitely supported probability measure, and let $b$ be its
barycenter. A radial Hilbert realization consists of vectors $v_i$ with
$\|v_i\|=d(b,x_i)$ and $\|v_i-v_j\|\leq d(x_i,x_j)$. The
Izeki--Nayatani invariant is
\[
 \delta_X(\nu)=\inf_{(v_i)}
 \frac{\|\sum_i t_i v_i\|^2}{\sum_i t_i\|v_i\|^2},
 \qquad \delta(X)=\sup_\nu\delta_X(\nu).
\]
This is Izeki and Nayatani (2005), arXiv version~1,
Definition~5.1, p.~22. Their Proposition~5.3, pp.~22--23,
compares the nonlinear graph spectral gap with the ordinary gap by the
factor $1-\delta(X)$. We record its particular consequence for the
rank-one walk, with all normalization constants explicit.

\begin{proposition}[PG-IN-GATE]
\label{pg:in-gate}
If $X$ is complete CAT(0) and $\delta(X)\leq\delta_0<1$, then every
$f:\operatorname{Mat}_n(\mathbb F_2)\to X$ satisfies
\[
 D_f(n)^2\leq
 \frac{4(1-2^{-n})^2}{1-\delta_0}\,D_f(1)^2.
 \tag{PG-IN}
\]
Consequently, if every finite price over a field has an isometric
realization in some complete CAT(0) space $X$ with the same upper bound
$\delta(X)\leq\delta_0<1$, that field is contracting. One may take any
integer $n$ for which
\[
 n>\frac{4(1-2^{-n})^2}{1-\delta_0},
 \qquad
 s=\frac{2(1-2^{-n})}{\sqrt{1-\delta_0}}.
\]
It suffices, more strongly, to choose $n>4/(1-\delta_0)$.
\end{proposition}
\begin{proof}
Let $b$ be the barycenter of the uniform image measure, and set
$V=\mathbb E_A d(b,f(A))^2$. If $V=0$, the conclusion is immediate.
For each $\varepsilon>0$, the definition gives vectors $v_A$, identical
on repeated image points, satisfying
\[
 \|v_A\|=d(b,f(A)),\quad
 \|v_A-v_C\|\leq d(f(A),f(C)),\quad
 \|\mathbb E_Av_A\|^2\leq(\delta_0+\varepsilon)V.
\]
Apply the rank-one Poincar\'e inequality from
Theorem~\ref{pg:tree-contraction}, whose gap is
$\gamma_n=2^{2n-1}/(2^n-1)^2$. This yields
\[
 D_f(1)^2\geq2\gamma_n
 \bigl(\mathbb E\|v_A\|^2-\|\mathbb E v_A\|^2\bigr)
 \geq2\gamma_n(1-\delta_0-\varepsilon)V.
\]
The triangle inequality about $b$ gives $D_f(n)^2\leq4V$.
Let $\varepsilon\downarrow0$ and use
$2/\gamma_n=4(1-2^{-n})^2$. The asserted choice of $s$ has
$s^2<n$; also $s\geq1$ for $n\geq1$ and $\delta_0\geq0$.
Apply this bound to a permitted realization of each price on
$\operatorname{Mat}_n(\mathbb F_2)$.
\end{proof}

\paragraph{Known geometric classes and the scope of the gate.}
Izeki and Nayatani (2005), arXiv version~1, Example~3, p.~25,
and Proposition~5.5, p.~26, give $\delta=0$ for trees and Hilbert
spaces and preservation of the zero bound by finite metric products.
Thus their tree radial realization is the known mechanism used in
Theorem~\ref{pg:tree-contraction}; the rank-step constant is the
specialization calculated here. In particular, finite products of trees
also satisfy the same $n=4$, $s=15/8$ contraction estimate while
including Euclidean spaces of every finite dimension.

Fujiwara and Toyoda (2012), Theorem~1.4 (arXiv version~1, p.~2;
proof via Proposition~3.4, p.~7), prove
$\delta(X)\leq1/2$ for every complete CAT(0) cube complex, with no
dimension bound. Therefore all maps to these spaces satisfy, for $n=8$,
\[
 D_f(8)\leq \frac{255\sqrt2}{128}D_f(1),
 \qquad \left(\frac{255\sqrt2}{128}\right)^2
 =\frac{65025}{8192}<8.
\]
These examples show that even fairly broad CAT(0) classes can contain all
Euclidean configurations and still satisfy a uniform strict rank-step
contraction. They do not exhibit an actual field whose prices have both
properties. Nor has closure of prices under metric products been proved.

\begin{corollary}[PG-IN-BUILDING]
\label{pg:in-building}
Let $X_n$ be the type-$A_{2n-1}$ building cone used in
Proposition~\ref{pg:building}. Then
\[
 \delta(X_n)\geq
 1-\frac{4(1-2^{-n})^2}{n}.
\]
In particular $\delta(X_n)\to1$ as $n\to\infty$.
\end{corollary}
\begin{proof}
The map in Proposition~\ref{pg:building} has
$D_f(1)^2=1$ and $D_f(n)^2=n$. If $\delta(X_n)<1$, substitute these
values in (PG-IN) with $\delta_0=\delta(X_n)$ and rearrange. If
$\delta(X_n)=1$, the inequality is immediate.
\end{proof}

For comparison, Izeki, Kondo, and Nayatani (2012), Proposition~4.4,
p.~730, bound the invariant of tangent cones of the $N$-dimensional
Euclidean building for $\operatorname{PGL}(N+1,\mathbb Q_r)$ by
\[
 \begin{cases}
 \bigl(1+\sqrt{(N-1)/(N+3)}\bigr)/2,&N\text{ odd},\\
 (N+1)/(N+2),&N\text{ even}.
 \end{cases}
\]
The residue prime $r$ does not enter these bounds, but the bounds tend to
one with dimension. These established estimates and
Corollary~\ref{pg:in-building} concern ordinary CAT(0) buildings;
no price membership is asserted.

The useful sufficient hypothesis in Proposition~\ref{pg:in-gate} is
\emph{existence} of uniformly good realizations. A finite price may
embed in an arbitrarily enlarged bad ambient space, so a large invariant
of one ambient realization is not an obstruction to that hypothesis.
A stronger sufficient hypothesis would uniformly bound $\delta$ for all
complete price spaces. Neither hypothesis follows from the currently
checked spectral axioms. This is a precise remaining route for Q1(b),
not a proof of contraction for $\mathbb C$ or $\mathbb F_2$.

\subsection{A five-point obstruction to every Euclidean-core construction}

\begin{proposition}[PG-FIVE-NONL1]
\label{pg:five-nonlone}
There is a five-point CAT(0)-embeddable metric which cannot be written
$d(i,j)=\|x_i-x_j\|+a_i+a_j$ with $a_i\geq0$, in any Euclidean
dimension. Thus replacing the planar core in Theorem~\ref{pg:four-core}
by a higher-dimensional Euclidean core does not extend that theorem to
five points.
\end{proposition}
\begin{proof}
Let $T$ be a tripod with center $o$, and choose one point at distance
$3/\sqrt{10}$ on each of its three branches. In the CAT(0) product
$T\times\mathbb R$, take these three points at real coordinate zero,
and take two further points $(o,\pm1/\sqrt{10})$. Write $L$ for the
three branch points and $R$ for the last two points. Distances within
$L$ equal $6/\sqrt{10}$, the distance within $R$ is $2/\sqrt{10}$,
and all six cross-distances are one.

Every $L_1$ metric on a three-element set $L$ and a two-element set $R$
satisfies the pentagonal inequality
\[
 \sum_{\{i,j\}\subset L}d(i,j)+
 \sum_{\{i,j\}\subset R}d(i,j)
 \leq \sum_{i\in L,j\in R}d(i,j).
 \tag{PG-PENT}
\]
Indeed, it suffices to check cut metrics: a cut containing $a$ members
of $L$ and $b$ members of $R$ has right side minus left side equal to
$(a-b)(a-b-1)\geq0$. For an arbitrary $L_1$ metric, the threshold
identity
$|u-v|=\int_{\mathbb R}|\mathbf1_{u>t}-\mathbf1_{v>t}|\,dt$
and integration over its coordinates prove the same inequality.
Our displayed metric instead has left side $20/\sqrt{10}=2\sqrt{10}>6$.

Every finite Euclidean metric is an $L_1$ metric: rotational invariance
of normalized spherical measure $\sigma$ gives
\[
 \int_{S^{D-1}}|\langle u,x-y\rangle|\,d\sigma(u)
 =c_D\|x-y\|\quad(c_D>0).
\]
A nonnegative pendant contribution $a_i+a_j$, for $i\ne j$, is a sum
of singleton cut metrics and hence is also $L_1$. Such a representation
would therefore satisfy (PG-PENT), a contradiction.
\end{proof}

This is an obstruction to a proposed construction, not a proof that the
five-point metric is a nonprice. The pentagonal inequality is also the
wire-price inequality that would follow from an asymptotic restriction
of the complete bipartite graph tensor $K_{3,2}$ to the disjoint union
of a triangle and an edge. The next result gives a genuine obstruction
to the corresponding \emph{one-copy border restriction}. We first give
the polynomial argument in two small dimensions, and then strengthen
it to every edge dimension using tensor capacity. The asymptotic
implication with extra resources remains open.

\subsection{A GCT obstruction to a pentagonal graph conversion}

All tensors in this subsection are over $\mathbb C$. For an integer
$q\geq2$, let $S_q$ be the graph tensor with a $q$-dimensional wire on
each edge of $K_{3,2}$. Parties $1,2,3$ are its three left vertices,
and parties $4,5$ its right vertices. Let
\[
 M_q=\sum_{i,j,k=1}^q e_{ij}\otimes e_{jk}\otimes e_{ki},
 \qquad
 T_q=M_q\otimes\sum_{\ell=1}^q e_\ell\otimes e_\ell.
\]
Thus $T_q$ is the graph tensor of a triangle on parties $1,2,3$ and
an edge between parties $4,5$. Both source and target have left local
dimensions $q^2$; their right local dimensions are $q^3$ and $q$,
respectively. A border restriction here means membership in the
Zariski closure of tensors obtained by one linear map at each party,
in the fixed target format $(q^2,q^2,q^2,q,q)$.

For ordinary polynomial degree $d$, set
\[
 \mathcal I_q(d)=
 \mathbb C[(\mathbb C^q)^{\otimes3}]^{\operatorname{SL}_q^3}_d.
\]
The attached GCT fragment K-001 records the classical identity
$\mathcal I_3=\mathbb C[I_6,I_9,I_{12}]$. We checked its primary
statement in Bremner and Hu (2013), arXiv version~2, Theorem~2.4,
p.~5. For $q=2$, Bremner, Bickis, and Soltanifar (2012), arXiv
version~1, Theorem~36, p.~14, give
$\mathcal I_2=\mathbb C[\Delta_4]$.

\begin{theorem}[PG-GCT-PENT: small-dimensional border obstruction]
\label{pg:gct-pent}
For $q=2$ and for $q=3$, $T_q$ is not a border restriction of $S_q$.
More generally, the same conclusion holds for any $q$ if there is a
positive-degree homogeneous invariant
\[
 f\in\mathbb C[(\mathbb C^{q^2})^{\otimes3}]
 ^{\operatorname{SL}_{q^2}^3}_d,
 \qquad f(M_q)\ne0,
 \qquad
 \dim\mathcal I_q(d)<\binom{q+d-1}{d}.
 \tag{PG-GCT-condition}
\]
\end{theorem}
\begin{proof}
First fix an invariant $f$ as in (PG-GCT-condition), omitting the last
inequality for the moment. Let
$W_q=(\mathbb C^q)^{\otimes3}$, and write
$\alpha\boxtimes\beta\in(\mathbb C^{q^2})^{\otimes3}$ for the
Kronecker product with corresponding local factors grouped together.
The polynomial
\[
 F(\alpha,\beta)=f(\alpha\boxtimes\beta)
\]
has bidegree $(d,d)$ and is invariant under independent
$\operatorname{SL}_q^3$ actions on $\alpha$ and $\beta$.
For example, an operator $h\in\operatorname{SL}_q$ on a local factor
of $\alpha$ induces $h\otimes\mathrm{id}$ on the grouped factor,
whose determinant is $(\det h)^q=1$. The same applies to $\beta$.
Consequently
\[
 F\in\mathcal I_q(d)\otimes\mathcal I_q(d).
\]
If $r=\dim\mathcal I_q(d)$, choose a basis $P_1,\ldots,P_r$;
then $F(\alpha,\beta)=\sum_{a,b=1}^r c_{ab}P_a(\alpha)P_b(\beta)$.
The assertion is also valid with $F=0$ when $r=0$.

A restriction at the two right parties of $S_q$ produces tensors
$\alpha(u),\beta(v)\in W_q$ depending linearly on variables
$u=(u_1,\ldots,u_q)$ and $v=(v_1,\ldots,v_q)$ dual to the two
right target spaces. Before the three left maps are applied, its
bilinear slice is exactly $\alpha(u)\boxtimes\beta(v)$. Therefore
\[
 f\bigl(\alpha(u)\boxtimes\beta(v)\bigr)
 =\sum_{a,b=1}^r c_{ab}P_a(\alpha(u))P_b(\beta(v)).
 \tag{PG-GCT-separation}
\]
Its coefficient matrix, with rows indexed by degree-$d$ monomials in
$u$ and columns by degree-$d$ monomials in $v$, has rank at most $r$.

The three left maps do not invalidate this bound. Put $m=q^2$.
Invariance under scalar matrices in $\operatorname{SL}_m$ forces
$m\mid d$. For an invertible $g_i\in\operatorname{GL}_m$, write
$g_i=\lambda_i h_i$ with $h_i\in\operatorname{SL}_m$. Homogeneity
and invariance give the polynomial identity
\[
 f\bigl((g_1\otimes g_2\otimes g_3)z\bigr)
 =\prod_{i=1}^3(\det g_i)^{d/m} f(z).
 \tag{PG-GCT-character}
\]
It holds for singular $g_i$ as well, by polynomial continuation.
Thus every restriction $Z$ of $S_q$ has coefficient-matrix rank at
most $r$ for the polynomial $f(Z(u,v))$. The matrix entries are
polynomials in the entries of $Z$. Its $(r+1)$-minors therefore vanish
on the entire Zariski closure of the restrictions.

For the proposed target, however,
\[
 T_q(u,v)=\left(\sum_{\ell=1}^q u_\ell v_\ell\right)M_q,
 \qquad
 f(T_q(u,v))=f(M_q)(u\mathbin{\cdot}v)^d.
\]
The coefficient matrix is diagonal in the monomial bases, with the
nonzero diagonal entries
$f(M_q)\binom d{\lambda_1,\ldots,\lambda_q}$ indexed by
$\lambda_1+\cdots+\lambda_q=d$. Its rank is
$\binom{q+d-1}{d}$. This proves the general obstruction.

It remains to supply $f$ and the dimension inequality for $q=2,3$.
The matrix multiplication tensor $M_q$ is polystable for
$\operatorname{SL}_{q^2}^3$: see B\"urgisser and Ikenmeyer (2017),
arXiv version~2, Corollary~4.9, p.~22. In particular there is a
positive-degree homogeneous invariant nonzero at $M_q$.
For clarity, the last implication only uses separation of its closed
nonzero orbit from zero: take a polynomial equal to one on that orbit
and zero at zero, apply the Reynolds projection to invariants, and
choose a positive homogeneous component nonzero at $M_q$.
No explicit degree or invariant formula is required.

For $q=2$, the cited invariant-ring description gives
$\dim\mathcal I_2(d)\leq1<d+1$ for every $d>0$.
For $q=3$, it gives
\[
 \dim\mathcal I_3(d)
 =\#\{(a,b,c)\in\mathbb Z_{\geq0}^3:6a+9b+12c=d\}
 \leq\binom{\lfloor d/6\rfloor+2}{2}
 <\binom{d+2}{2}.
\]
The first bound counts possible pairs $(a,b)$, since $c$ is then
uniquely determined and $a+b\leq\lfloor d/6\rfloor$.
These are precisely the required strict inequalities.
\end{proof}

\paragraph{What this does and does not certify.}
The cut-rank tests allow the proposed conversion: for a cut containing
$a$ left vertices and $b$ right vertices, the source has
$(a-b)(a-b-1)$ more crossing edges than the target. Thus
Theorem~\ref{pg:gct-pent} gives a border obstruction beyond all these
flattening tests, and uses the supplied GCT calibration concretely.
For comparison, exact restriction is already impossible for every
$q>1$: full target ranks at the three left parties force the left
maps to be invertible, while the two right maps must each have rank
$q$. The source is maximally entangled across the combined left/right
bipartition, so its image then has bipartition rank $q^2$; the target
has rank one across that bipartition.

The coefficient-rank argument uses the small invariant rings only for
$q=2,3$. Its dimension estimate alone does not cover $q^N$ for arbitrary
$N$. The following argument removes this limitation, but does not
remove the separate obstruction created by subexponential extra
resources in the definition of asymptotic restriction.

\subsection{A capacity obstruction in every edge dimension}

For $X\in(\mathbb C^m)^{\otimes3}$, define the \emph{norm capacity}
\[
 c_m(X)=\inf_{g\in\operatorname{SL}_m(\mathbb C)^3}\|gX\|.
\]
We use the ordinary Hilbert tensor norm. This is the square root of
the capacity convention of B\"urgisser, Garg, Oliveira, Walter, and
Wigderson (2018), arXiv version~1, equation~(1.3), p.~4.
The required standard Kempf--Ness facts are: every orbit closure
contains a unique closed orbit; minimum-norm vectors in that closed
orbit form one compact-unitary orbit; and a nonzero tensor is a
minimum-norm vector exactly when all of its normalized one-party
reduced matrices are scalar identities. Precise sources are B\"ohm
and Lafuente (2020), arXiv version~3, Theorem~1.1, p.~2, and
Lemma~5.1, p.~9, and B\"urgisser et al. (2018), Theorem~3.2, p.~13,
with its proof in Appendix~A, pp.~37--38. These statements apply to
the complex special-linear action considered as a real reductive
action with its unitary-invariant norm.

\begin{lemma}[PG-CAP-PRODUCT]
\label{pg:cap-product}
For $A\in(\mathbb C^a)^{\otimes3}$ and
$B\in(\mathbb C^b)^{\otimes3}$,
\[
 c_{ab}(A\boxtimes B)=c_a(A)c_b(B).
 \tag{PG-CAP-product}
\]
For $g_1,g_2,g_3\in\operatorname{GL}_m(\mathbb C)$,
\[
 c_m\bigl((g_1\otimes g_2\otimes g_3)X\bigr)
 =\left(\prod_{i=1}^3|\det g_i|^{1/m}\right)c_m(X).
 \tag{PG-CAP-covariance}
\]
Finally, if $f$ is a homogeneous $\operatorname{SL}_m^3$-invariant
polynomial of positive degree $d$, then
\[
 |f(X)|\leq C_f c_m(X)^d,
 \qquad C_f=\max_{\|Y\|=1}|f(Y)|.
 \tag{PG-CAP-invariant}
\]
\end{lemma}
\begin{proof}
The subgroup acting independently on the two factors is contained
in the larger special-linear group: locally,
$\det(h\otimes k)=(\det h)^b(\det k)^a=1$.
Approximating the two capacity infima therefore proves the upper
bound in (PG-CAP-product), including equality if either factor has
zero capacity.

Suppose both capacities are positive. Choose minimum-norm vectors
$A_0,B_0$ in the respective orbit closures; their norms are the two
capacities. They have scalar one-party reduced matrices, so the
one-party reduced matrices of $A_0\boxtimes B_0$ are scalar too.
Thus $A_0\boxtimes B_0$ is a minimum-norm vector for the larger
$\operatorname{SL}_{ab}^3$ action and its orbit is closed. It belongs
to the orbit closure of $A\boxtimes B$ by the subgroup inclusion
above. The unique-closed-orbit theorem identifies its orbit with the
closed orbit containing a norm-minimizer of $A\boxtimes B$.
The unitary uniqueness of minimum-norm vectors on a closed orbit
then gives
$c_{ab}(A\boxtimes B)=\|A_0\boxtimes B_0\|
=\|A_0\|\|B_0\|$.

For (PG-CAP-covariance), factor each $g_i$ as a scalar times a
special-linear matrix and use invariance and absolute homogeneity
of capacity. For the last inequality, invariance and homogeneity
give $|f(X)|=|f(hX)|\leq C_f\|hX\|^d$ for every
$h\in\operatorname{SL}_m^3$. Take the infimum over $h$.
\end{proof}

\begin{theorem}[PG-CAP-PENT: all-dimensional border obstruction]
\label{pg:cap-pent}
For every integer $q\geq2$, the tensor
$T_q=K_3(q)\otimes\operatorname{EPR}_q$ is not a border restriction
of $S_q=K_{3,2}(q)$ over $\mathbb C$.
Consequently, for every $q\geq2$ and every positive integer $N$,
$T_q^{\otimes N}$ is not a border restriction of $S_q^{\otimes N}$.
\end{theorem}
\begin{proof}
Assume a sequence $Z_j$ of restrictions of $S_q$ converges to $T_q$.
Such a sequence exists if $T_q$ belongs to the Zariski closure in
the definition of border restriction: the restriction image is the
image of a polynomial map over $\mathbb C$, hence constructible,
and its Euclidean and Zariski closures coincide. Put $m=q^2$.
Every one-party left flattening of $T_q$ has rank $m$, so the same
holds for $Z_j$ for all sufficiently large $j$. Its three left
restriction maps $g_{1,j},g_{2,j},g_{3,j}$ must therefore be
invertible $m\times m$ matrices.

As in the preceding theorem, its two right restriction maps determine
linear maps $\alpha_j,\beta_j:\mathbb C^q\to
(\mathbb C^q)^{\otimes3}$ such that
\[
 Z_j(u,v)=(g_{1,j}\otimes g_{2,j}\otimes g_{3,j})
       \bigl(\alpha_j(u)\boxtimes\beta_j(v)\bigr).
\]
By Lemma~\ref{pg:cap-product}, for a positive scalar $\kappa_j$
independent of $u,v$,
\[
 c_m(Z_j(u,v))=\kappa_j c_q(\alpha_j(u))c_q(\beta_j(v)).
\]
Thus the capacity of these slices separates in the two right
variables. In particular, writing $Z_j^{ab}=Z_j(e_a,e_b)$,
\[
 c_m(Z_j^{11})c_m(Z_j^{22})
 =c_m(Z_j^{12})c_m(Z_j^{21}).
 \tag{PG-CAP-minor}
\]

The proposed target has $T_q^{11}=T_q^{22}=M_q$ and
$T_q^{12}=T_q^{21}=0$. Choose a positive-degree homogeneous
invariant $f$ with $f(M_q)\ne0$, whose existence follows from the
matrix multiplication polystability used in
Theorem~\ref{pg:gct-pent}. Polynomial continuity and
(PG-CAP-invariant) imply that both $c_m(Z_j^{11})$ and
$c_m(Z_j^{22})$ are bounded below by a fixed positive number for
all large $j$. On the other hand,
$c_m(Z_j^{12})\leq\|Z_j^{12}\|\to0$ and
$c_m(Z_j^{21})\leq\|Z_j^{21}\|\to0$. This contradicts
(PG-CAP-minor). This proof requires no continuity assertion for
capacity itself.

Finally $S_q^{\otimes N}$ and $T_q^{\otimes N}$ identify, by
regrouping the wire bases, with $S_{q^N}$ and $T_{q^N}$,
respectively. The first assertion therefore applies to every
equal tensor power.
\end{proof}

\paragraph{Exact scope of the obstruction.}
Theorem~\ref{pg:cap-pent} excludes border restrictions at every
equal tensor power with no extra resource. It strengthens the
small-dimensional invariant calculation and remains invisible to all
cut-rank inequalities. It does \emph{not} exclude a conversion with
a subexponential GHZ overhead, or with $o(N)$ additional source
copies: those modifications enlarge the left source spaces, so the
left maps need no longer be square or invertible and the capacity
factorization in the proof is no longer available. No uniform
positive asymptotic rate gap follows from the argument. Accordingly,
the asymptotic $K_{3,2}\to K_3\sqcup K_2$ conversion remains open
here; this result supplies neither a price violating (PG-PENT) nor
a decision of the square or flat/contracting alternatives.

\begin{corollary}[PG-CAP-HYPER: higher hypermetric graph obstruction]
\label{pg:cap-hyper}
For every $r\geq3$ and $q\geq2$, the graph tensor
$K_r(q)\otimes K_{r-1}(q)$ is not a border restriction of
$K_{r,r-1}(q)$ over $\mathbb C$, with the two complete graphs
on the two respective parts. This holds at every equal tensor power,
although every bipartition flattening test allows the conversion.
\end{corollary}
\begin{proof}
The proof of Lemma~\ref{pg:cap-product} applies verbatim to $r$-party
norm capacity and to a product of $r-1$ factors. Each left local
space of source and target has dimension $m=q^{r-1}$. Along a
hypothetical convergent sequence of restrictions, the full target
left flattenings again force all left maps to be invertible for
large indices. Slicing the $r-1$ right parties of the source gives
a Kronecker product of $r-1$ tensors in $(\mathbb C^q)^{\otimes r}$.
Its capacity, after the left maps, is therefore a product of
functions of the separate right variables.

The left target $K_r(q)$ has scalar one-party reduced matrices and
is a nonzero minimum-norm vector; consequently some positive-degree
invariant is nonzero there. Fix product-basis contractions at all
but two right parties of $K_{r-1}(q)$. At each of the remaining two
parties, fix the labels of all other incident edges and allow the
label of their mutual edge to be either $1$ or $2$. These four
slices of the right graph tensor have scalar values $1,0,0,1$.
The resulting capacity matrix therefore leads to exactly the same
contradiction as (PG-CAP-minor), using the fixed invariant to bound
the two diagonal capacities below and the norm to bound the two
off-diagonal capacities above. Equal powers replace $q$ by $q^N$.

For a cut containing $a$ left vertices and $b$ right vertices, the
source has $(r-1)a+rb-2ab$ crossing edges and the target has
$a(r-a)+b(r-1-b)$. Their difference is
$(a-b)(a-b-1)\geq0$, proving the last assertion.
\end{proof}

The same limitations concerning extra asymptotic resources apply to
Corollary~\ref{pg:cap-hyper}; it is not a spectral hypermetric
inequality in either direction.

\subsection{An integral full-space cover by matching-deleted graphs}

The following finite identity retains the full local tensor spaces. It is
stronger than a scalar Cauchy identity after symmetric projection, but it
currently supplies no subquadratic selector bound as the number of parties
grows.

\begin{proposition}[PG-CAUCHY-R3-INTEGRAL]
\label{pg:cauchy-r3-integral}
Over every field, the six-party graph tensor
$H_3=K_3(2)\otimes K_3(2)$ has a cover by $18$ restrictions of
$K_{3,3}(2)$, each using only the six edges outside one perfect matching.
Equivalently,
\[
 H_3\leq K_{3,3}(2)\otimes\mathrm{GHZ}_{18}.
 \tag{PG-CAUCHY-18}
\]
Each term has coefficient $1$ or $-1$. The identity holds over the integers
before reduction to any field; no division by two is involved.
\end{proposition}
\begin{proof}
Here is the complete coefficient certificate, including its port convention.
Use the alternating edge tensor
$\epsilon=e_0\otimes e_1-e_1\otimes e_0$, which is related to the binary
EPR tensor by an invertible map on one endpoint over every field. Order the
parties as $L_0,L_1,L_2,R_0,R_1,R_2$. At each party use two physical slots,
labelled by the other two labels in its own part, in ascending order.
If their bits are $b_0,b_1$, the local basis index is $b_0+2b_1$.
The target is the product of the six $\epsilon$ edges in the two triangles,
with each edge oriented from its lower label to its higher label.

For $\sigma\in S_3$, let $J_\sigma$ be the crown graph obtained by omitting
$L_iR_{\sigma(i)}$ from $K_{3,3}$. Orient every remaining edge from left to
right. The edge $L_iR_{\sigma(\ell)}$, $i\ne\ell$, uses the slot labelled
$\ell$ at $L_i$ and the slot labelled $\sigma(i)$ at $R_{\sigma(\ell)}$.
For an integer $0\leq M<64$, let $\tau_M$ interchange the two slots at
party $v$ precisely when bit $v$ of $M$ is one. With these definitions,
\[
 H_3=\sum_{(\sigma,M,c)\in\mathcal T}c\,\tau_MJ_\sigma,
 \tag{PG-CAUCHY-port}
\]
where the entire list $\mathcal T$ is
\[
\begin{array}{c|r|r@{\qquad}c|r|r@{\qquad}c|r|r}
 \sigma&M&c&\sigma&M&c&\sigma&M&c\\\hline
 012&6&-1&021&0&1&120&38&-1\\
 012&9&-1&102&1&1&120&54&1\\
 012&14&1&102&8&1&201&0&-1\\
 012&44&-1&102&9&-1&210&22&1\\
 012&45&1&120&17&-1&210&33&1\\
 120&18&-1&120&19&1&210&55&-1
\end{array}
\]
Here, for example, $120$ denotes the permutation with
$(\sigma(0),\sigma(1),\sigma(2))=(1,2,0)$.

For completeness, the exact finite check of (PG-CAUCHY-port) is as follows.
For each of the $2^6$ assignments of a bit $a_e$ to the first endpoint of
its six oriented edges, place the complementary bit $1-a_e$ at the second
endpoint and add the coefficient $(-1)^{\sum_e a_e}$ to the resulting
six-party basis coordinate. Do this separately for the target and for each
listed crown, applying its six local swaps. The eighteen signed crown
arrays sum to the target array in all $4^6=4096$ coordinates. Each individual
array has exactly $64$ nonzero coordinates, all equal to $1$ or $-1$.
This rule is implemented using only Python integer arithmetic in
\path{programs/pg_verify_cauchy_r3_integral.py}; the complete input is
\path{inputs/pg_cauchy_r3_integral_cover.json}, and the exact result is
\path{outputs/pg_cauchy_r3_integral_verification.json}.
An independent verifier,
\path{programs/rt_verify_crown_integral_cover.py}, computes each boundary
coefficient directly as a product of six $\epsilon$ entries, rather than
enumerating internal edge assignments. It independently checks all
$4096$ coordinates; its result is
\path{outputs/rt_crown_integral_verification.json}.
Thus the certificate is an identity over $\mathbb Z$.

To realize each crown as a restriction of the full graph, contract the left
endpoint of each omitted alternating edge with $e_0^*$ and its right endpoint
with $e_1^*$; their contraction is one. Apply the listed port swaps, and
absorb $c$ into one party's map. These operations give the eighteen claimed
restrictions. A common $18$-level GHZ selector selects the corresponding
local map at every party, so their sum is a single restriction from the
source with that selector. In particular five binary GHZ copies suffice.
Finally conjugate the per-leg alternating forms back to ordinary EPR forms.
\end{proof}

Each source term in Proposition~\ref{pg:cauchy-r3-integral} uses six binary
wires. For arbitrary positions $x_i,y_j$ of its six parties, its wire cost is
\[
 W_\sigma=\sum_{i\ne\ell}\|x_i-y_{\sigma(\ell)}\|
 =\sum_{i,j}\|x_i-y_j\|-\sum_i\|x_i-y_{\sigma(i)}\|.
\]
The eighteen terms use the six matching classes with multiplicities
$(5,1,3,5,1,3)$ in lexicographic permutation order. Taking $N$ tensor powers
of this displayed identity alone gives $18^N$ terms. It supplies neither
subexponential growth in $N$ nor a family with selector exponent $o(r^2)$.

\begin{corollary}[PG-CAUCHY-R3-DENSE: thirteen six-wire terms]
\label{pg:cauchy-r3-dense}
The eighteen terms in Proposition~\ref{pg:cauchy-r3-integral} can be grouped
into thirteen restrictions of the same crown sources, over every field.
Thus a $13$-level selector, or four binary GHZ copies, suffices when these
party-local linear combinations are used.
\end{corollary}
\begin{proof}
Let $S_v$ interchange the two slots at party $v$. The five terms in class
$012$ are the following sum of three products of party-local maps:
\[
 -S_1S_2(I-S_3)-S_0S_3+S_2S_3S_5(S_0-I).
\]
The three terms in class $102$ become
$S_0(I-S_3)+S_3$. The five terms in class $120$ become
\[
 -S_0S_4(I-S_1)-S_1S_4-S_1S_2S_5(I-S_4).
\]
Leave the other classes unchanged, with respectively one, one, and three
terms. Each displayed summand is a tensor product of local maps, since
its nonidentity factors act at different parties. Expansion gives exactly
the eighteen signed monomials already proved, reducing their count to
$3+1+2+3+1+3=13$ without division. The standalone integer polynomial
expansion is saved in \path{programs/pg_verify_cauchy_r3_dense.py}, with
input \path{inputs/pg_cauchy_r3_dense_cover.json} and result
\path{outputs/pg_cauchy_r3_dense_verification.json}.
\end{proof}

\paragraph{Comparison with a four-term cover using eight wires.}
The six-wire source formats above are a substantive part of the diagnostic;
the full graph itself has the simpler all-field bound
\[
 H_3(q)\ \text{is a sum of at most }q^2\text{ restrictions of }G_3(q).
 \tag{PG-ROUTING-four}
\]
Indeed slice the two target edges $L_0L_2$ and $R_0R_2$ at their respective
basis labels $a,b\in\{1,\ldots,q\}$. For each fixed pair $(a,b)$, route
the four remaining edges along the edge-disjoint paths
\[
 L_0R_0L_1,\quad L_1R_2L_2,\quad
 R_0L_2R_1,\quad R_1L_0R_2.
\]
They use eight distinct source edges; the unused edge is $L_1R_1$.
At each intermediate party contract the two path legs by the equality
form, retain the endpoint legs, and insert the fixed sliced label in its
appropriate target port. Contract both ends of the unused edge to its
first basis vector. The result is precisely the sliced target tensor,
with coefficient one. Summing over $(a,b)$ proves (PG-ROUTING-four),
giving four terms at $q=2$. This uses more source wires per term than the
crown identity, and its overhead exponent $2$ does not satisfy SH-33.

\begin{proposition}[PG-CAUCHY-CLASS: every matching class is compulsory]
\label{pg:cauchy-class}
Let $r\geq2$. Put a copy of $\mathbb F^2$ on every edge, and define
$H_r=K_r(2)\otimes K_r(2)$ using alternating edge tensors as above. For
$\tau\in S_r$, let $J_\tau$ be the matching-deleted graph
$K_{r,r}\setminus\{L_iR_{\tau(i)}:0\leq i<r\}$, with any fixed port
assignment. Suppose a cover identity expresses $H_r$ as a linear
combination of local port permutations of the $J_\tau$. If $a_\tau$ is
the sum of the coefficients in matching class $\tau$, then
\[
 a_\tau=(-1)^{r(r-1)/2}\operatorname{sgn}(\tau).
 \tag{PG-CAUCHY-class}
\]
In particular, no matching class can be omitted over any field. The same
conclusion holds if each local map is a linear combination of port
permutations, with $a_\tau$ interpreted as the induced scalar after symmetric
quotient at every party. This is a restriction on this specified ansatz,
not a restriction on arbitrary local linear maps from the crown graphs.
\end{proposition}
\begin{proof}
Apply at every party the division-free quotient
$(\mathbb F^2)^{\otimes(r-1)}\to\operatorname{Sym}^{r-1}(\mathbb F^2)$.
Represent it by multiplying its $r-1$ linear coordinates. Port permutations
act trivially on this quotient in every characteristic. Write a bit weight
$w$ for the monomial $z_0^{r-1-w}z_1^w$ at one party.
For each $\sigma\in S_r$, inspect the coefficient with left weights
\[
 (0,1,\ldots,r-1)
\]
and right weight $\sigma^{-1}(j)$ at $R_j$.

In a source crown $J_\tau$, reorder the right columns as
$R_{\sigma(0)},\ldots,R_{\sigma(r-1)}$. Let $A_{i\ell}$ be the bit at the
left end of its edge in row $i$ and reordered column $\ell$, and set a
missing edge's entry equal to zero. The prescribed left weights make the
row sums $i$. Since an alternating edge has complementary endpoint bits,
the prescribed right weights make the column sums $r-1-\ell$.
There is exactly one zero-one $r\times r$ matrix with these margins:
\[
 A_{i\ell}=\begin{cases}1,&i>\ell,\\0,&i\leq\ell.\end{cases}
\]
Indeed row zero is zero, and column zero must therefore be one in every
remaining row. Removing that row and column gives the same argument
recursively. The missing edge in row $i$ occupies column
$\sigma^{-1}\tau(i)$ and must lie at a zero of this staircase. Thus
$\sigma^{-1}\tau(i)\geq i$ for every $i$. Summing these inequalities forces
equality in all of them, so $\tau=\sigma$.

For $\tau=\sigma$, the unique assignment has $r(r-1)/2$ left bits equal
to one, and its coefficient is $(-1)^{r(r-1)/2}$. Hence the $r!$ selected
coordinates of the $r!$ projected crowns form this signed identity matrix.
The target's left complete-graph product has coefficient one at weights
$(0,1,\ldots,r-1)$. Its right complete-graph product is the Vandermonde
\[
 \prod_{i<j}(z_{i0}z_{j1}-z_{i1}z_{j0});
\]
the coefficient at weights $\sigma^{-1}(j)$ is
$\operatorname{sgn}(\sigma)$, as follows from its determinant expansion.
Comparing this coordinate gives (PG-CAUCHY-class). Its values are units
also in characteristic two. A linear combination of local port
permutations induces the scalar sum of its coefficients on the quotient,
which proves the stated extension of the ansatz.
\end{proof}

For $r=3$, the six selected symmetric coordinates have ternary local weights
$(0,1,2;0,1,2)$, $(0,1,2;0,2,1)$, $(0,1,2;1,0,2)$,
$(0,1,2;2,0,1)$, $(0,1,2;1,2,0)$, and $(0,1,2;2,1,0)$,
respectively. The source minor is $-I_6$, and the forced class sums are
$(-1,1,1,-1,-1,1)$, agreeing with the integral identity above. The independent
boundary-coordinate verifier also checks this minor. This coefficient
obstruction does not claim that a cut-rank inequality forces each class.
Moreover EX-31 already proves a stronger asymptotic limitation for pure
port-permutation covers: their term count is at least
$2^{r(r-1)/2}$. Each such restriction is a local partial permutation in
the edge basis, so its row and column sparsities in EX-31 are one.
Consequently any subquadratic selector construction must combine many
port permutations inside dense local maps, or leave this ansatz.


\begin{proposition}[PG-CAUCHY-COMMUTANT: six dense matching terms do not suffice]
\label{pg:cauchy-commutant}
Use the six crowns $J_\sigma$ and port convention of
Proposition~\ref{pg:cauchy-r3-integral}. Over every field, there is no identity
\[
 H_3=\sum_{\sigma\in S_3}
       \left(\bigotimes_{v=0}^5 A_{\sigma v}\right)J_\sigma,
 \qquad A_{\sigma v}\in\operatorname{span}\{I,S_v\},
 \tag{PG-CAUCHY-six}
\]
where $S_v$ swaps the two ports at party $v$. Thus one dense local
permutation-algebra term per matching does not lift the scalar Cauchy
identity, even at $r=3$. This does not exclude additional terms in the
same matching class or local maps outside the displayed span.
\end{proposition}
\begin{proof}
Write $A_{\sigma v}=a_{\sigma v}I+b_{\sigma v}S_v$ and
$u_{\sigma v}=a_{\sigma v}+b_{\sigma v}$. The symmetric quotient at all
parties, together with Proposition~\ref{pg:cauchy-class}, forces
\[
 \prod_vu_{\sigma v}=c_\sigma,
 \qquad(c_\sigma)_{\sigma\in S_3}=(-1,1,1,-1,-1,1)
\]
in lexicographic permutation order. In particular every $u_{\sigma v}$
is nonzero. Divide the local map by this scalar and put
$t_{\sigma v}=b_{\sigma v}/u_{\sigma v}$, so the normalized map is
$I+t_{\sigma v}(S_v-I)$.

For each $v$, let $Q_{-v}$ take the symmetric quotient at the other five
parties, retaining the full four-dimensional space at $v$. Put
$J_{\sigma v}=Q_{-v}J_\sigma$ and $H_v=Q_{-v}H_3$. Two integral identities
and rank checks are used:
\[
 H_v=\sum_\sigma c_\sigma J_{\sigma v},\qquad
 \{(S_v-I)J_{\sigma v}:\sigma\in S_3\}
       \text{ are linearly independent over every field}.
 \tag{PG-CAUCHY-retained}
\]
Here is an explicit finite certificate. At each quotiented party send local
basis indices $0,1,2,3$ to weights $0,1,1,2$, respectively, with coefficient
one. Index the resulting coordinates lexicographically with radix four at
$v$ and radix three elsewhere. The first identity holds in every one of
the $4\cdot3^5=972$ coordinates for each $v$. For the second assertion,
the following six-row minors, with columns in the same permutation order,
have the indicated determinants:
\[
\begin{array}{c|l|r}
 v&\text{zero-based row indices}&\det\\\hline
0&260,266,278,284,290,302&1\\
1&98,104,116,122,128,140&-1\\
2&44,50,143,149,155,248&1\\
3&53,86,125,158,191,230&-1\\
4&65,89,137,161,185,233&1\\
5&69,93,141,165,189,237&-1
\end{array}
\]
For example the first minor is
\[
\begin{pmatrix}
0&0&1&-1&1&-1\\
1&-1&0&0&-1&1\\
0&0&-1&0&-1&0\\
-1&1&-1&1&1&1\\
-1&0&0&0&0&-1\\
0&0&1&0&0&0
\end{pmatrix}.
\]
The standalone integer checker
\path{programs/pg_verify_cauchy_r3_commutant.py} constructs all source and
target coefficients directly from their boundary $\epsilon$ entries,
checks all six $972$-coordinate identities, and computes all six minors
exactly. Their full matrices are retained in
\path{outputs/pg_cauchy_r3_commutant_quotient.json}. Unit determinants,
not a characteristic-zero diagonalization of $S_v$, establish the
independence in every characteristic. An independent reconstruction in
\path{programs/ex_verify_crown_commutant.py} instead enumerates the $64$
internal edge assignments of each graph, quotients sparse support, and
computes determinants by their permutation expansion. It checks all $5832$
retained-party coordinates and reproduces the six unit minors; see
\path{outputs/ex_crown_commutant_independent.json}.

Now apply $Q_{-v}$ to (PG-CAUCHY-six). The normalized maps at the other
parties induce the identity on the symmetric quotient. Subtract the first
identity in (PG-CAUCHY-retained) to obtain
\[
 \sum_\sigma c_\sigma t_{\sigma v}(S_v-I)J_{\sigma v}=0.
\]
The second identity forces every $t_{\sigma v}=0$. This holds for all six
parties, so (PG-CAUCHY-six) would reduce to
$H_3=\sum_\sigma c_\sigma J_\sigma$. It does not: at the full local-basis
coordinate $(0,0,1,1,3,3)$, of flattened index $95$, the target coefficient
is zero and this bare sum has coefficient one. The same exact checker
verifies this discrepancy. Hence the proposed identity is impossible over
any field.
\end{proof}

\begin{proposition}[PG-CAUCHY-LIFT-NOGO: one permutation-algebra term per matching]
\label{pg:cauchy-lift-nogo}
Let $r\geq3$ and use the following canonical crown ports: the edge
$L_iR_{\sigma(\ell)}$, $i\ne\ell$, uses port $\ell$ at $L_i$ and port
$\sigma(i)$ at $R_{\sigma(\ell)}$. The ports at each party are labelled
by the other labels in its own part. Over every field, there is no identity
\[
 H_r=\sum_{\sigma\in S_r}
       \left(\bigotimes_v A_{\sigma v}\right)J_\sigma
 \tag{PG-CAUCHY-general-lift}
\]
with one term per matching and every $A_{\sigma v}$ in the local algebra
spanned by port permutations. More generally this remains impossible when
each local map satisfies
\[
 Q_v A_{\sigma v}=u_{\sigma v}Q_v
\]
for some scalar $u_{\sigma v}$, where
$Q_v:(\mathbb F^2)^{\otimes(r-1)}\to
\operatorname{Sym}^{r-1}(\mathbb F^2)$ is the symmetric quotient.
Thus a direct $r!$-term lift with these local maps already fails for every
$r\geq3$. Multiple terms in a matching class are not excluded.
\end{proposition}
\begin{proof}
Quotient all parties symmetrically. The unit minor in
Proposition~\ref{pg:cauchy-class} forces
\[
 \prod_vu_{\sigma v}=c_\sigma,
 \qquad c_\sigma=(-1)^{r(r-1)/2}\operatorname{sgn}(\sigma).
\]
These are nonzero in every field, so all the local scalars are nonzero.
Divide each local map by its scalar, writing $B_{\sigma v}$ for the result.
Then $Q_vB_{\sigma v}=Q_v$, and the coefficient outside the product of
normalized maps is $c_\sigma$.

Fix a matching $\sigma$ and a party $v$. Let $\rho$ be an arbitrary ordering
of the left labels, so the position of label $i$ is $\rho^{-1}(i)$.
Order the right labels by $\sigma\rho$ and prescribe right weight
$\rho^{-1}\sigma^{-1}(j)$ at $R_j$. At every party other than $v$, apply
the symmetric quotient and inspect the coefficient with these left and
right weights. Do not quotient party $v$.

Consider first a retained left party $L_i$. All right weights have been
specified. As every source edge has one bit at each endpoint and the two
bits sum to one, the specified right weights and the other left weights
force the unobserved left input weight to be $\rho^{-1}(i)$. Reorder rows
by $\rho$ and columns by $\sigma\rho$. The source edge-bit matrix now has
row sums $0,1,\ldots,r-1$ and column sums $r-1,r-2,\ldots,0$.
The staircase proof of Proposition~\ref{pg:cauchy-class} applies unchanged:
only the matching $\sigma$ can contribute, and its retained-party input is
exactly the basis vector $e_b$ whose one-ports are the labels preceding
$i$ in the order $\rho$. Its coefficient is $(-1)^{r(r-1)/2}$.
For a retained right party the same reasoning, using all left weights to
force its missing right weight, again isolates $\sigma$. The one-ports at
$R_j$ are the right labels preceding $j$ in the order $\sigma\rho$.

The target coefficient at the retained party is the same basis vector
$e_b$ multiplied by $\operatorname{sgn}(\sigma)$. Indeed its two clique
orientations have the transitive score orders $\rho$ and $\sigma\rho$;
their signs multiply to
$\operatorname{sgn}(\rho)\operatorname{sgn}(\sigma\rho)
=\operatorname{sgn}(\sigma)$. Equivalently this is the product of the two
Vandermonde coefficients in Proposition~\ref{pg:cauchy-class} with one
party left unquotiented. The total edge degree forces that party's missing
weight, so no additional tournament orientation contributes.

All normalized maps at the quotiented parties disappear because
$Q_wB_{\tau w}=Q_w$. Comparing the retained coefficient in
(PG-CAUCHY-general-lift) therefore yields
\[
 \operatorname{sgn}(\sigma)e_b
 =c_\sigma(-1)^{r(r-1)/2}B_{\sigma v}e_b,
 \qquad\text{hence}\qquad B_{\sigma v}e_b=e_b.
\]
Varying the order $\rho$ realizes every subset of the other labels as the
predecessors of the retained vertex. Thus $e_b$ ranges over every local
binary basis vector, forcing $B_{\sigma v}=I$ for all $\sigma,v$.
Any proposed identity must consequently be the bare signed sum
$H_r=\sum_\sigma c_\sigma J_\sigma$.

This bare identity fails for every $r\geq3$. It fails at $r=3$ by the
coefficient $(0,0,1,1,3,3)$ in
Proposition~\ref{pg:cauchy-commutant}. To reduce a larger $r$ to $r-1$,
contract both highest-labelled parties $L_{r-1},R_{r-1}$ to the all-one
basis covector, and contract the port labelled $r-1$ at every remaining
party to the zero basis covector. On the target this leaves exactly
$H_{r-1}$, with coefficient one. A crown with
$\sigma(r-1)\ne r-1$ contains the edge
$L_{r-1}R_{r-1}$ and is killed by its two one-endpoints. If
$\sigma(r-1)=r-1$, the contraction leaves the canonical $(r-1)$-vertex
crown, with scalar $(-1)^{r-1}$. Multiplying its original class sign gives
\[
 (-1)^{r-1}(-1)^{r(r-1)/2}\operatorname{sgn}(\sigma)
 =(-1)^{(r-1)(r-2)/2}\operatorname{sgn}(\sigma|_{[r-1]}).
\]
Thus the bare $r$-identity would imply the bare $(r-1)$-identity.
Iteration contradicts the checked $r=3$ coefficient. No field-dependent
division or semisimplicity has been used.
\end{proof}

\paragraph{A limit on using the unused matching as a literal selector.}
In an $N$-fold word of matching-deleted graphs, every party has exactly $N$
unused binary wire legs. Consequently any GHZ tensor obtained solely from
those unused wires has dimension at most $2^N$, by a one-party flattening.
A literal selector distinguishing all words in a type profile $\nu$ has
size $M_\nu=N!/\prod_j\nu_j!$. If $b$ external binary GHZ copies are also
allowed, this implementation therefore needs
\[
 b\geq\log_2 M_\nu-N.
\]
For a balanced profile of a fixed $A$-term identity, this becomes
$b\geq N(\log_2 A-1)-O_A(\log N)$. In particular, the unused matching
wires do not absorb the eighteen-letter profile entropy of
Proposition~\ref{pg:cauchy-r3-integral}. This bound concerns that literal
selector implementation; it does not exclude an algebraic compression
which avoids distinguishing every word.

\subsection{A crossing-operator construction and its missing regularization}

The crossing operator means of Bug\'ar and Vrana (2025) give a more
concrete R2 test, but their Schur-power limit must not be confused
with regularization under the tensor product in the tensor semiring.
Here all spaces and tensors are complex. Fix four parties in cyclic
order, the two bipartitions $b=12|34$ and $c=14|23$, and their
equal-weight two-variable operator geometric mean. Let $F_\alpha$,
$0<\alpha<1$, denote the resulting upper functional of order
$\alpha$ in their Example~4.10, equation~(121), p.~2817.
Its existence, subadditivity, submultiplicativity, scaling law, and
finite-block lower bound are Theorem~4.5, equations~(89)--(97),
pp.~2810--2811; the geometric-mean construction is
Proposition~4.8, pp.~2815--2816.

Define
\[
 p_0(T)=\limsup_{\alpha\downarrow0}F_\alpha(T),\qquad
 r(T)=\inf_{N\geq1}p_0(T^{\otimes N})^{1/N}.
 \tag{PG-OP-regularization}
\]
No existence of an ordinary $\alpha\downarrow0$ limit is assumed.
The second operation is the semiring power regularization; its
omission is a substantive gap, as the exact example below shows.

\begin{proposition}[PG-OP-GATE]
\label{pg:op-gate}
The functions $p_0$ and $r$ are upper functionals in the sense of
Anthropic \S A.4, and $r(T^{\otimes k})=r(T)^k$ for $k\geq1$.
Put $C_q=\operatorname{EPR}_{13,q}\otimes
\operatorname{EPR}_{24,q}$. Then
\[
 q\leq p_0(C_q)\leq q^2,\qquad
 2\leq r(C_2)\leq2^{\sqrt2}.
 \tag{PG-OP-bounds}
\]
The last upper bound uses the checked, conditional CAT(0) price-space
chain. If $r(C_2)>2$, the Euclidean square is a price over
$\mathbb C$, and consequently all finite planar Euclidean metrics
are prices. Whether $r(C_2)>2$ remains undecided; Theorem~\ref{pg:op-complete}
below proves that this inequality is equivalent to the square-price gate.
\end{proposition}
\begin{proof}
The bounds in Bug\'ar and Vrana, equation~(91), imply pointwise
boundedness of $F_\alpha$ as $\alpha\downarrow0$. Its normalization,
subadditivity, and submultiplicativity pass to the upper limit.
To check arbitrary linear restrictions, write any nonzero local
restriction map as a positive scalar times a contraction. Monotonicity
under local contractions and the order-$\alpha$ scaling law give
$F_\alpha(LT)\leq K^{2\alpha}F_\alpha(T)$, where $K$ is the
product of the local operator norms. Thus $p_0(LT)\leq p_0(T)$.
A finite cover is a restriction of the direct sum of its covering
tensors, so $p_0$ satisfies the cover inequality. The power
regularization step in Anthropic Proposition~A.5, pp.~20--21,
proves the asserted properties of $r$.

For a wire of dimension $q$ that crosses exactly one of $b,c$,
the second cut observable is the identity and the geometric mean
is the square root of the first. The bipartite Schur formula gives
$F_\alpha(\operatorname{EPR}_q)=q^{(1+\alpha)/2}$: apply
Bug\'ar and Vrana, Lemma~4.6, pp.~2813--2815, with
$\beta=2\alpha/(1+\alpha)$ and then raise its value to
$\alpha/\beta$.
If it crosses both cuts, the two observables coincide and
$F_\alpha(\operatorname{EPR}_q)=q$. Hence $p_0$ assigns wire
values $\sqrt q$ to the four cyclic sides and $q$ to both diagonals.
These statements concern individual wires, not simultaneous
multiplicativity on their products.

The dimension bound (91) gives $p_0(C_q)\leq q^2$. For the lower
bound, in the $n$-copy Schur decomposition select an arbitrary
partition $\lambda\vdash n$ with at most $q$ rows at parties
$1,3$, and the trivial partition $(n)$ at parties $2,4$.
The EPR-pairing tensor has a nonzero component in this subspace.
Both crossing observables there have the same eigenvalue
$2^{nH(\lambda/n)}$. The finite-block inequality (95), followed by
$\alpha\downarrow0$, yields $p_0(C_q)\geq2^{H(\lambda/n)}$.
Take partitions approaching the uniform distribution on $q$ rows
to get $p_0(C_q)\geq q$. Since $C_2^{\otimes N}\cong C_{2^N}$,
this proves $r(C_2)\geq2$.

By Lemma~\ref{pg:saturation}, a spectral point $\varphi\leq r$
satisfies $\varphi(C_2)=r(C_2)$. Its wire pseudometric has each
cyclic side at most $1/2$. CAT(0) roundness gives
$d_\varphi(1,3)^2+d_\varphi(2,4)^2\leq1$, and hence
\[
 \log_2 r(C_2)=d_\varphi(1,3)+d_\varphi(2,4)\leq\sqrt2.
\]
In a tree metric the four-point condition instead bounds this sum
by the larger of the two opposite-side sums, which is at most one.
Thus $r(C_2)>2$ would give a non-tree price. The square gate proved
above then gives the conclusion.
\end{proof}

\paragraph{An exact finite-block warning.}
At $q=2,n=4$ the crossing Schur projectors with partition
$(2,1,1)$ have a common component of $C_2^{\otimes4}$ of squared
norm $18$. The exact rational certificate is
\path{outputs/pg_crossing_n4_exact.json}, generated by
\path{programs/pg_crossing_operator.py}. It contains the
positive definite Gram matrix $G$, the two projector matrices
$P_b,P_c$, the vector $h$ of inner products with $C_2^{\otimes4}$,
and a rational vector $v$ satisfying
\[
 P_bv=P_cv=v,\qquad v^TGv=h^Tv=18.
 \tag{PG-OP-certificate}
\]
The Gram construction uses the permutation-pair states
$|\sigma\rangle_{13}\otimes|\tau\rangle_{24}$, their simultaneous
conjugacy orbits, and the exact inner products
$2^{\#\mathrm{cycles}(\sigma^{-1}\sigma')+
\#\mathrm{cycles}(\tau^{-1}\tau')}$. Character projectors use
the Murnaghan--Nakayama recursion. All certificate identities are
over $\mathbb Q$; floating-point eigenvalue calculations are not
needed for these identities. The separate verifier
\path{programs/pg_verify_crossing_certificate.py} rebuilds the
Gram and projector matrix elements using the elementary character
$\chi_{(2,1,1)}(\pi)=\operatorname{sgn}(\pi)
(\#\operatorname{Fix}(\pi)-1)$.
An independent full-coordinate check in
\path{programs/rt_verify_crossing_block.py} reconstructs the
actual $65536$-coordinate tensor, without using the Gram matrices
or character recursion, and verifies both projector equations and
the two norm identities. Its output is
\path{outputs/rt_crossing_block_verification.json}.

For a representation-theoretic description, take the individual
$S_4$-types $(3,1)$ at parties $1,3$ and $(2,2)$ at parties $2,4$.
Their tensor product decomposes as $(3,1)\oplus(2,1,1)$.
Swapping parties $2,4$ is the identity on the two-dimensional
invariant space: its negative eigenspace is the invariant part of
$(3,1)\otimes(3,1)\otimes\Lambda^2(2,2)$.
Here $\Lambda^2(2,2)=(1,1,1,1)$ is the sign representation, whereas
$(3,1)\otimes(3,1)=(4)\oplus(3,1)\oplus(2,2)\oplus(2,1,1)$
contains no sign representation.
The crossing projectors therefore coincide on this block and their
$(2,1,1)$ component has nonzero EPR-pairing projection. The exact
certificate also includes the block with these two local types
interchanged and gives the total squared norm $18$.

Since $H((2,1,1)/4)=3/2$, equation~(95) of the cited source gives
\[
 F_\alpha(C_2)\geq
 18^{\alpha/4}\,2^{(1-\alpha)3/2},
 \qquad p_0(C_2)\geq2^{3/2}>2^{\sqrt2}\geq r(C_2).
 \tag{PG-OP-strict-drop}
\]
The last two comparisons are conditional on the CAT(0) chain used
in Proposition~\ref{pg:op-gate}. They show that the semiring power
regularization is genuinely necessary here. In the Schur-copy limit
defining $F_\alpha(C_2)$, the permutation group acts on individual
copies; in $F_\alpha(C_2^{\otimes N})$ it acts on $N$-copy blocks
with local dimension $2^N$. These are different operators.
The fixed $q=2$ block alone supplies no lower bound exceeding $2^N$
for $p_0(C_{2^N})$. The high-precision values for blocks $n=3,4,5$
are retained as exploratory outputs, separately from the exact
rational certificate and from the unresolved regularized gate.

\begin{proposition}[PG-OP-RAW-ALLQ: a vanishing raw advantage]
\label{pg:op-raw-allq}
For every integer $q\geq2$, the unregularized endpoint satisfies
\[
 p_0(C_q)\geq q\,2^{1/q}.
 \tag{PG-OP-raw-allq}
\]
This estimate does not improve the proved lower bound $r(C_2)\geq2$.
\end{proposition}
\begin{proof}
Write $[\lambda]$ for the real Specht representation of $S_n$ of
shape $\lambda$. First take integers $a,b\geq2$, $n=ab$, and set
$V=[n-1,1]$, $W=[a^b]$. The permutation representation
$P=\mathbb R^n$ is $\mathbf1\oplus V$. The branching rule and
$P\otimes W\cong\operatorname{Ind}_{S_{n-1}}^{S_n}
\operatorname{Res}_{S_{n-1}}^{S_n}W$ give the multiplicity-free
decomposition
\[
 V\otimes W=[\lambda_+]\oplus[\lambda_-],\qquad
 \lambda_+=(a+1,a^{b-2},a-1),\quad
 \lambda_-=(a^{b-1},a-1,1).
 \tag{PG-OP-rectangle}
\]
Indeed, the rectangle has one removable corner, and the resulting
shape has exactly three addable corners. One restores $W$ and the
other two give the displayed summands. The branching and two-box
restriction facts used here are the cases $k=1,2$ of the Pieri rule,
together with Frobenius reciprocity; see Sam (2017),
Theorem~3.5.7, p.~24, \S4.1, p.~27, and
Proposition~4.3.2, p.~29.

We next show
\[
 \operatorname{Hom}_{S_n}(\Lambda^2V,W\otimes W)=0.
 \tag{PG-OP-no-exterior}
\]
Since
$\Lambda^2P\cong
\operatorname{Ind}_{S_{n-2}\times S_2}^{S_n}
(\mathbf1\boxtimes\operatorname{sgn})$, restriction and
Frobenius reciprocity give
\[
 \big\langle W\otimes W,\Lambda^2P\big\rangle
 =2\sum_{\nu\vdash n-2}
 c^{(a^b)}_{\nu,(2)}c^{(a^b)}_{\nu,(1,1)}=0.
\]
The two Littlewood--Richardson coefficients record removal of a
horizontal and a vertical two-box strip, respectively. From a
rectangle these removals are the last two boxes of its bottom row
and the bottom two boxes of its last column. They give different
remaining shapes, so every product in the sum vanishes.
As $\Lambda^2P=V\oplus\Lambda^2V$, nonnegativity of
multiplicities proves (PG-OP-no-exterior).

Consider the simultaneous $S_n$-invariants in
$V\otimes W\otimes V\otimes W$, with factors in party order
$1,2,3,4$. By the multiplicity-free decomposition
(PG-OP-rectangle), this invariant space has dimension two.
The simultaneous interchange $(1\,3)(2\,4)$ is the identity on
it: each summand is the symmetric invariant pairing of a real
irreducible representation with itself. A vector on which
$(2\,4)$ acts as $-1$ would therefore also be odd under
$(1\,3)$, and would belong to the invariant part of
$\Lambda^2V\otimes\Lambda^2W$. Equation
(PG-OP-no-exterior) shows that this subspace is zero.
Thus $(2\,4)$ is the identity on the invariant space.

The crossing Schur projectors for $12|34$ and $14|23$ are
conjugate by $(2\,4)$, so they coincide on this invariant space.
The EPR-pairing tensor has nonzero projection onto each of the two
summands: after grouping parties $12$ and $34$, it is the identity
pairing on $V\otimes W$, whose restriction to each irreducible
summand is nonzero. Therefore its $[\lambda_-]$ component is a
nonzero common eigenvector of both crossing observables.

Now take $a=2$, $b=q$, $n=2q$. Both local types $V$ and $W$
occur in $(\mathbb C^q)^{\otimes n}$, since their numbers of rows
are $2$ and $q$. The local Schur multiplicity spaces contribute
nonzero EPR pairings and do not change the crossing eigenvalues.
The common type is $\lambda_-=(2^{q-1},1,1)$, with entropy
\[
 H(\lambda_-/n)=\log_2q+\frac1q.
\]
If the squared norm of this nonzero projection is $c_q>0$, the
finite-block bound gives
\[
 F_\alpha(C_q)\geq c_q^{\alpha/(2q)}
 2^{(1-\alpha)(\log_2q+1/q)}.
\]
Letting $\alpha\downarrow0$ proves the assertion. At $q=2^N$
the resulting root lower bound is
$2\,2^{1/(N2^N)}$, which tends to $2$. Hence it supplies no
strict lower bound for the infimum defining $r(C_2)$.
\end{proof}

\begin{proposition}[PG-OP-MULTIPLICITYFREE: a uniform scope bound]
\label{pg:op-multiplicityfree}
Let $q\geq2$, and let $\lambda,\mu\vdash n$ each have at most $q$ rows.
If the entire Kronecker product $[\lambda]\otimes[\mu]$ is
multiplicity-free, then every constituent $[\nu]$ has at most $4q$ rows.
Consequently
\[
 H(\nu/n)\leq\log_2q+2.
 \tag{PG-OP-multiplicityfree}
\]
In particular, common-crossing certificates confined to such local Specht
products can supply entropy thresholds at most $\log_2q+2$. At
$q=2^N$, the $N$th roots of the corresponding endpoint lower bounds have
ceiling $2\,4^{1/N}\to2$. This entire multiplicity-free construction class
therefore cannot establish a uniform regularized lower bound
$r(C_2)>2$. This is not an upper bound for the full functional $p_0$, or
for blocks with genuine Kronecker multiplicities.
\end{proposition}
\begin{proof}
We use the complete classification of Bessenrodt and Bowman (2017),
Theorem~1.1, pp.~1--2 of the author version dated September~21, 2017.
Conjugation of either partition is allowed independently in that theorem,
so this must be tracked rather than discarded.

First, if a constituent $[\nu]$ occurs in
$[\lambda]\otimes[\mu]$, then
\[
 \ell(\nu)\leq\ell(\lambda)\ell(\mu).
\]
Indeed realize the two Specht modules inside tensor powers of spaces of
dimensions $\ell(\lambda)$ and $\ell(\mu)$, respectively. Their product
embeds in the $n$th tensor power of the tensor product of those spaces,
whose Schur--Weyl constituents have at most the product dimension many
rows.

In classification case~(1), a trivial factor gives the other partition;
a sign factor with at most $q$ rows forces $n\leq q$. In case~(2), the
standard factor either has two rows, giving at most $2q$ rows in a
constituent, or its conjugate has $n-1\leq q$ rows, giving
$n\leq q+1$. In cases~(3)--(4), a balanced two-line factor either has
two rows or, after conjugation, forces $n\leq2q$; if both factors remain
two-row shapes, the row-product bound is at most four. In case~(5), the
near-trivial factor has at most three rows, or its conjugate forces
$n-2\leq q$. Thus the respective bounds are $3q$ or $n\leq q+2$.
Finally case~(6) consists of the three exceptional pairs of degrees nine
and twelve. Every orientation of those pairs requires $q\geq3$, so
$n\leq12\leq4q$. These cases establish the uniform bound $4q$.

Entropy is at most the logarithm of the support size, proving
(PG-OP-multiplicityfree). In the four-party block with local types
$\lambda,\mu,\lambda,\mu$, both crossing cut observables have constituent
types from $[\lambda]\otimes[\mu]$. Their eigenvalue entropy thresholds
are therefore bounded as stated. Substituting $q=2^N$ gives the root
ceiling. No assertion concerning the other local Specht blocks has been
made.
\end{proof}

\begin{proposition}[PG-OP-PT: a positive partial-transpose test]
\label{pg:op-pt}
Let $V=[\lambda]$, $W=[\mu]$ be real orthogonal Specht models for $S_n$,
with $\ell(\lambda),\ell(\mu)\leq q$. Write $\Gamma$ for partial transpose
on the $W$ matrix indices in $\operatorname{End}(V\otimes W)$, and let
$P_{\geq h}$ be the sum of the constituent projectors with
$H(\nu/n)\geq h$. If a nonzero positive semidefinite operator $X$ satisfies
\[
 [X,\rho_V(g)\otimes\rho_W(g)]=0\quad(g\in S_n),\qquad
 \Gamma X=X,\qquad P_{\geq h}X=X,
 \tag{PG-OP-PT-test}
\]
then $p_0(C_q)\geq2^h$.

A sufficient test of this form is a rational linear feasibility problem on
central operators. Let $m_\nu=g(\lambda,\mu,\nu)>0$,
$d_\nu=\dim[\nu]$, and let $P_\nu$ be the associated isotypic projector.
Define
\[
 K_{\nu\tau}=\operatorname{Tr}(P_\nu\Gamma P_\tau),\qquad
 D_{\nu\tau}=2\bigl(\delta_{\nu\tau}d_\nu m_\nu-K_{\nu\tau}\bigr).
 \tag{PG-OP-PT-Gram}
\]
It suffices to find a nonzero real vector $x$ satisfying
\[
 x_\nu\geq0,\qquad Dx=0,\qquad
 x_\nu=0\ \text{when }H(\nu/n)<h.
 \tag{PG-OP-PT-linear}
\]
The entries of $K$ have the exact character formula
\[
 K_{\nu\tau}=
 \frac{d_\nu d_\tau}{(n!)^2}
 \sum_{g,h\in S_n}
 \chi_\nu(g)\chi_\tau(h)\chi_\lambda(gh)
 \chi_\mu(gh^{-1}).
 \tag{PG-OP-PT-character}
\]
These are sufficient finite-block tests. They provide a regularized gap
only if a lower bound stronger than $q$ is established uniformly at the
grouped dimensions required in Proposition~\ref{pg:op-gate}.
\end{proposition}
\begin{proof}
Under the real orthogonal pairings, identify
\[
 (V_A\otimes W_B\otimes V_C\otimes W_D)^{S_n}
 \cong\operatorname{End}_{S_n}(V\otimes W).
\]
The EPR-pairing vector corresponds to the identity operator. Interchanging
$B,D$ sends a matrix entry $X_{ab,cd}$ to $X_{ad,cb}$, which is exactly
$\Gamma$. The two crossing projectors are conjugate by this interchange.
Consequently a vector represented by $X$ in (PG-OP-PT-test) lies in both
high-entropy crossing subspaces. Its overlap with the EPR pairing is
$\operatorname{Tr}X>0$, because $X$ is nonzero and positive semidefinite.
The local Schur multiplicity spaces supply nonzero pairing factors and do
not alter this conclusion.

For each fixed $0<\alpha<1$, each crossing observable in the finite-block
formula dominates $2^{n(1-\alpha)h/\alpha}$ times the orthogonal projection
onto this common vector. Monotonicity of the operator geometric mean then
gives the same domination for their geometric mean. The finite-block
inequality of Bug\'ar and Vrana (2025), Theorem~4.5, equation~(95), therefore
gives $F_\alpha(C_q)\geq c^{\alpha/n}2^{(1-\alpha)h}$ for some $c>0$
independent of $\alpha$. Taking $\limsup_{\alpha\downarrow0}$ proves the
first assertion.

Partial transpose is a self-adjoint orthogonal involution for the
Hilbert--Schmidt inner product. The matrix $D$ is precisely the Gram
matrix
\[
 D_{\nu\tau}=\langle P_\nu-\Gamma P_\nu,
                         P_\tau-\Gamma P_\tau\rangle.
\]
Thus, for $X=\sum_\nu x_\nu P_\nu$ with real coefficients,
$\Gamma X=X$ is equivalent to $Dx=0$. The projectors are mutually
orthogonal, so $X\geq0$ is equivalent to all $x_\nu\geq0$, and the support
condition is the displayed vanishing condition. Zero-multiplicity
constituents have been omitted, so a nonzero nonnegative vector gives
$\operatorname{Tr}X>0$.

Finally expand the central projectors as
\[
 P_\nu=\frac{d_\nu}{n!}\sum_{g\in S_n}
 \chi_\nu(g)\rho_V(g)\otimes\rho_W(g).
\]
Since the models are real orthogonal, partial transpose replaces
$\rho_W(h)$ by $\rho_W(h^{-1})$. Taking the trace of the product gives
$\chi_\lambda(gh)\chi_\mu(gh^{-1})$, proving
(PG-OP-PT-character). Characters of $S_n$ are integer-valued and agree on
inverse permutations, so the resulting matrix is rational. The summand is
invariant under simultaneous conjugation of $g,h$. Hence one may sum $g$
over the $p(n)$ conjugacy-class representatives with their class-size
weights and retain the full $h$ sum, reducing the exact computation to
$p(n)n!$ permutation pairs. This does not assume that $\Gamma$ preserves
the central subspace: the Gram equation tests equality in the entire
commutant.
\end{proof}



\begin{lemma}[PG-OP-SIGNED: positivity is sufficient, not necessary]
\label{pg:op-signed}
In Proposition~\ref{pg:op-pt}, replace positive semidefiniteness of $X$
by the condition $\operatorname{Tr}X\ne0$. Retain real symmetry,
commutation with the diagonal $S_n$ action, $\Gamma X=X$, and
$P_{\geq h}X=X$. Then the same conclusion $p_0(C_q)\geq2^h$ holds.
\end{lemma}
\begin{proof}
Vectorization gives a nonzero common high-crossing vector $v$ whose EPR
pairing is $\operatorname{Tr}X\ne0$. Each positive crossing observable
dominates $2^{n(1-\alpha)h/\alpha}$ times the rank-one \emph{orthogonal
projector} onto $v$. The signs of the eigenvalues of the matrix $X$ do
not affect this statement about the vector $v$. The geometric-mean and
finite-block argument in Proposition~\ref{pg:op-pt} therefore applies
unchanged.
\end{proof}

\begin{proposition}[PG-OP-SOS-DUALITY: an exact finite obstruction]
\label{pg:op-sos-duality}
Fix $q,n\geq1$ and $\lambda\vdash n$ with $\ell(\lambda)\leq q$.
Put $V=[\lambda]$, $M=\dim V$, and
\[
 B_{\lambda}(h)=
 \sum_{\substack{\nu\vdash n\\H(\nu/n)<h}}
       d_\nu g_{\lambda\lambda\nu}.
\]
If $\sigma_{\mathbb C}(M)>B_{\lambda}(h)$, there exists a real symmetric
$S_n$-equivariant operator $X$ on $V\otimes V$ with
\[
 \Gamma X=X,\qquad P_{\geq h}X=X,\qquad
 \operatorname{Tr}X\ne0.
 \tag{PG-OP-SOS-witness}
\]
Consequently $p_0(C_q)\geq2^h$. No positivity or centrality of $X$ is
asserted.
\end{proposition}
\begin{proof}
Use a real orthogonal Specht model. Let $\mathcal E$ be the finite
real inner product space of real symmetric operators on $V\otimes V$
which commute with the diagonal $S_n$ action. Its inner product is
$\langle A,B\rangle=\operatorname{Tr}(AB)$. Partial transpose $\Gamma$
is a self-adjoint orthogonal involution on $\mathcal E$: it commutes
with total transpose, and partial transpose of conjugation by
$\rho(g)\otimes\rho(g)$ is conjugation by the same real orthogonal
operator. Write $\mathcal E_+$ and $\mathcal E_-$ for its two
eigenspaces.

Every element of the commutant is block diagonal with respect to the
isotypic decomposition
\[
 V\otimes V=\bigoplus_\nu [\nu]\otimes
                              \mathbb R^{g_{\lambda\lambda\nu}}.
\]
Let $\mathcal L$ consist of its symmetric blocks supported on
$H(\nu/n)<h$, and let $\mathcal H=\mathcal L^\perp$ be the corresponding
high-support space. Suppose no $X\in\mathcal H\cap\mathcal E_+$ has
nonzero trace. The identity operator is then perpendicular to this
intersection. Orthogonal-complement duality in a finite-dimensional
real inner product space gives
\[
 I\in(\mathcal H\cap\mathcal E_+)^\perp
   =\mathcal L+\mathcal E_-.
\]
Thus $I=Y+Z$ for real symmetric $Y\in\mathcal L$ and
$Z\in\mathcal E_-$, and applying $\Gamma$ gives
\[
 Y+\Gamma Y=2I.
 \tag{PG-OP-SOS-dual}
\]
In particular, $\operatorname{rank}Y\leq B_\lambda(h)$.

For arbitrary real $a,b\in V$, partial transpose leaves evaluation on
$a\otimes b$ unchanged. Therefore (PG-OP-SOS-dual) implies the polynomial
identity
\[
 (a\otimes b)^TY(a\otimes b)
       =(a^Ta)(b^Tb).
 \tag{PG-OP-SOS-polynomial}
\]
If $R=\operatorname{rank}Y$, a real orthogonal diagonalization of $Y$,
followed by complex square roots of its nonzero eigenvalues, gives
$Y=B^TB$ with $B$ an $R\times M^2$ complex matrix. Each coordinate of
$B(a\otimes b)$ is bilinear in $a,b$. Hence (PG-OP-SOS-polynomial)
is a sum-of-squares identity of length $R$ over $\mathbb C$; equality on
real variables already gives equality of the complex polynomials.
This proves
$\sigma_{\mathbb C}(M)\leq R\leq B_\lambda(h)$, contrary to the
hypothesis. The desired $X$ therefore exists, and
Lemma~\ref{pg:op-signed} applies. This argument is complex-specific at
the rank-factorization step; no characteristic-two conclusion is drawn.
\end{proof}

\begin{theorem}[PG-OP-COMPLETE: the regularized operator gate is exact]
\label{pg:op-complete}
Conditional on the checked price-space and planar-shortening chains used
above, the operator functional of Proposition~\ref{pg:op-gate} satisfies
\[
 r(C_2)=
 \begin{cases}
 2,&\text{if the Euclidean square is not a price over }\mathbb C,\\
 2^{\sqrt2},&\text{if the Euclidean square is a price over }\mathbb C.
 \end{cases}
 \tag{PG-OP-exact-gate}
\]
In the second case, for every integer $q\geq2$ one has the uniform
bound $p_0(C_q)\geq q^{\sqrt2}$. Thus the regularized crossing-operator
gate agrees with the universal cycle-cover gate. This identifies the
two gates; it does not decide which value occurs.
\end{theorem}
\begin{proof}
Only the square-present implication needs a new argument. Fix an integer
$q\geq2$. For each $n$, choose a partition $\lambda_n\vdash n$ with at
most $q$ rows whose Specht dimension $M_n=d_{\lambda_n}$ is maximal
among those partitions. The needed elementary representation estimates
are Schur--Weyl duality and the Weyl and hook dimension bounds; see
Christandl et al. (2018), \S3.1, equation~(11), p.~17, and
Remark~3.7, equations~(12)--(13), p.~19. They give
\[
 q^n=\sum_{\ell(\lambda)\leq q}
          \dim S_\lambda(\mathbb C^q)\,d_\lambda,
 \qquad
 M_n\geq\frac{q^n}{(n+1)^{q+q(q-1)/2}}.
 \tag{PG-OP-max-Specht}
\]

Regrouping the two copies of $(\mathbb C^q)^{\otimes n}$ and comparing
Schur--Weyl multiplicities gives, for each $\nu\vdash n$,
\[
 \dim S_\nu(\mathbb C^{q^2})
 =\sum_{\ell(\alpha),\ell(\beta)\leq q}
       g_{\alpha\beta\nu}
       \dim S_\alpha(\mathbb C^q)
       \dim S_\beta(\mathbb C^q).
 \tag{PG-OP-multiplicity-bound}
\]
All summands are nonnegative integers. Thus a constituent of
$[\lambda_n]\otimes[\lambda_n]$ has at most $q^2$ rows and
\[
 g_{\lambda_n\lambda_n\nu}
 \leq\dim S_\nu(\mathbb C^{q^2})
 \leq(n+1)^{q^2(q^2-1)/2}.
\]
The hook dimension bound also gives
$d_\nu\leq n!/\prod_i\nu_i!\leq2^{nH(\nu/n)}$. Since there are at
most $(n+1)^{q^2}$ possible partitions, for every real $h$,
\[
 B_{\lambda_n}(h)
 \leq (n+1)^{q^2+q^2(q^2-1)/2}\,2^{nh}.
 \tag{PG-OP-low-dimension}
\]
The constants and polynomial exponents here may depend on the fixed
$q$; no bound uniform in $q$ is needed while $n$ tends to infinity.

If the square is a price, Proposition~\ref{rt:sos} gives
\[
 \log_2\sigma_{\mathbb C}(M)
       \geq\sqrt2\log_2M-O(\sqrt{\log M})
       \qquad(M\longrightarrow\infty).
\]
Together with (PG-OP-max-Specht), this yields
\[
 \log_2\sigma_{\mathbb C}(M_n)
       \geq\sqrt2\,n\log_2q-O_q(\sqrt n+\log n).
\]
Fix $h<\sqrt2\log_2q$. The last bound eventually exceeds the
logarithm of (PG-OP-low-dimension). Proposition~\ref{pg:op-sos-duality}
therefore produces, for all sufficiently large $n$, a high-support
$\Gamma$-fixed vector with nonzero EPR overlap, and gives
$p_0(C_q)\geq2^h$. Letting $h$ increase to $\sqrt2\log_2q$ proves
$p_0(C_q)\geq q^{\sqrt2}$ for this $q$. Since $q$ was arbitrary, this
holds at every grouped dimension $q=2^N$. Consequently
\[
 r(C_2)=\inf_{N\geq1}p_0(C_{2^N})^{1/N}
       \geq2^{\sqrt2}.
\]
The opposite inequality is Proposition~\ref{pg:op-gate}.

If the square is absent, the same proposition already says that
$r(C_2)>2$ would force its presence, while $r(C_2)\geq2$ always.
Thus $r(C_2)=2$. Comparing with the proved two-value cycle-cover gate
gives the final assertion. No finite-block estimate was tensor-amplified:
for each fixed grouped dimension, the common vector was obtained anew
from its full growing-row Schur decomposition.
\end{proof}


\begin{corollary}[PG-OP-DETERMINISTIC: a prescribed rational matrix family]
\label{pg:op-deterministic}
For every dyadic integer $q=2^N$, $N\geq1$, make the following fixed
choices. Put
\[
 A_q=q(q+1)/2,\qquad B_q=q^2(q^2+1)/2.
\]
For $n$ a power of two, define the integers
\[
 L_n=\left\lceil\frac{q^n}{(n+1)^{A_q}}\right\rceil,
 \quad m_n=\lfloor\log_2L_n\rfloor,
 \quad t_n=\left\lceil\sqrt{2m_n-1}\right\rceil
\]
when $m_n\geq1$. Let $n(q)$ be the least such power of two for which
$k_n=m_n-1-t_n>0$ and
\[
 2^{28k_n}>(n+1)^{20B_q}q^{25n}.
 \tag{PG-OP-integer-size}
\]
Choose the lexicographically first maximum-dimension partition
$\lambda\vdash n(q)$ with at most $q$ rows. These choices are
well-defined and computable using integer arithmetic.

Take a rational Specht model $R$ for $[\lambda]$ and a rational positive
definite invariant Gram matrix $G$. Form the rational linear system in
an unknown matrix $X\in\operatorname{End}(V\otimes V)$:
\[
\begin{split}
 &(G\otimes G)X=X^T(G\otimes G),\\
 &[X,R(s_i)\otimes R(s_i)]=0\quad(1\leq i<n(q)),\\
 &\Gamma_G X=X,\qquad P_{<h_q}X=0,\qquad\operatorname{Tr}X=1,\\
 &h_q=\tfrac54\log_2q.
\end{split}
 \tag{PG-OP-rational-system}
\]
where $s_i=(i,i+1)$ and
$\Gamma_GX=(I\otimes G^{-1})X^{T_2}(I\otimes G)$.
Then, conditional on the same checked chains as
Theorem~\ref{pg:op-complete}, the Euclidean square is absent over
$\mathbb C$ if and only if (PG-OP-rational-system) is inconsistent for
at least one $N\geq1$.
\end{corollary}
\begin{proof}
For fixed $q$, $m_n=n\log_2q-O_q(\log n)$ and $t_n=O_q(\sqrt n)$.
After taking binary logarithms, the leading terms in
(PG-OP-integer-size) are $28n\log_2q$ on the left and
$25n\log_2q$ on the right. Thus the inequality eventually holds even
when $n$ is restricted to powers of two, proving termination of the
specified search. The maximum Specht dimension $M$ is at least $L_n$.
If the square is a price, the finite estimate in
Proposition~\ref{rt:sos}, monotonicity in the number of variables, and
$\sqrt2>7/5$ imply
\[
 \log_2\sigma_{\mathbb C}(M)
 \geq\sqrt2\bigl(m_n-1-\sqrt{2m_n-1}\bigr)
 >\tfrac75 k_n.
\]
Taking twentieth roots in (PG-OP-integer-size) shows that
\[
 \sigma_{\mathbb C}(M)>(n+1)^{B_q}q^{5n/4}
             \geq B_\lambda(h_q).
\]
Proposition~\ref{pg:op-sos-duality} therefore gives a solution after
rescaling its nonzero trace to one.

All coefficients in the displayed system are rational. One may obtain
$R$ from the usual rational polytabloid construction, and take
$G=\sum_{g\in S_n}R(g)^TR(g)$. The central isotypic projectors are
\[
 P_\nu=\frac{d_\nu}{n!}\sum_{g\in S_n}
       \chi_\nu(g)R(g)\otimes R(g).
\]
Their high/low classification is also exact: at this value of $h_q$,
\[
 H(\nu/n)\geq h_q
 \quad\Longleftrightarrow\quad
 n^{4n}\geq q^{5n}\left(\prod_i\nu_i^{\nu_i}\right)^4.
\]
Conjugating by $G^{1/2}\otimes G^{1/2}$ transforms the rational
symmetry and partial-transpose equations into the orthonormal equations
of Proposition~\ref{pg:op-sos-duality}; the trace is unchanged. A
rational linear system is consistent over $\mathbb R$ if and only if it
is consistent over $\mathbb Q$. Thus the formulation introduces no
algebraic square-root ambiguity.

Conversely, if the square is absent, Theorem~\ref{pg:op-complete}
gives $r(C_2)=2$. Since $2<2^{5/4}$, its defining infimum supplies
some $N\geq1$ with $p_0(C_{2^N})<(2^N)^{5/4}$. At that grouped
dimension, any solution of (PG-OP-rational-system) would contradict
Lemma~\ref{pg:op-signed}. The corresponding prescribed system is
therefore inconsistent. This proves both implications. The system may
be extremely large; the result replaces arbitrary cover searches by a
specified finite rational family, rather than asserting computational
tractability or deciding the branch.
\end{proof}

The first sizes in the prescribed family are
$n(2)=2048$, $n(4)=8192$, and $n(8)=131072$. The integer-only checker
\path{programs/rt_operator_gate_sizes.py}, with its input and output
of the same stem, records every preceding power of two and verifies the
first passing inequalities. These are sizes of the specified maximum-dimension systems, whose
feasibility has not been solved here. Proposition~\ref{pg:op-determinant}
below solves the corresponding first three compressed and rectangular
systems by explicit polynomials.


\begin{proposition}[PG-OP-COMPRESSED: a symmetric-power encoding and direct dual]
\label{pg:op-compressed}
Let $q,n\geq1$ and work in the rational space
\[
 \mathcal V_{q,n}=\operatorname{Sym}^n((\mathbb Q^q)^{\otimes4}),
 \qquad
 D_{q,n}=\dim\mathcal V_{q,n}
        =\binom{n+q^4-1}{q^4-1}.
\]
For $b=12|34$ and $c=14|23$, let $P_{b,<h}$ and $P_{c,<h}$ denote the
sums of the crossing Schur projectors with $H(\nu/n)<h$, restricted to
this space. Put $\Psi=C_q^{\otimes n}$. The rational linear system
\[
 P_{b,<h}v=0,\qquad P_{c,<h}v=0,
 \qquad\langle\Psi,v\rangle=1
 \tag{PG-OP-compressed-system}
\]
has the following alternatives.
\begin{enumerate}
\item If it is consistent, then $p_0(C_q)\geq2^h$.
\item If it is inconsistent, rational linear algebra constructs tensors
$S,U$ on the four grouped parties, each of dimension $q^n$, such that
\[
 C_q^{\otimes n}=S+U,
 \qquad
 \max\{\operatorname{rank}_{12|34}S,
          \operatorname{rank}_{14|23}U\}
 \leq(n+1)^{B_q}2^{nh},
 \tag{PG-OP-compressed-dual}
\]
where $B_q=q^2(q^2+1)/2$.
\end{enumerate}
At the prescribed $n=n(q)$ and $h=(5/4)\log_2q$ from
Corollary~\ref{pg:op-deterministic}, square absence over $\mathbb C$ is
equivalent to inconsistency of (PG-OP-compressed-system) for at least
one dyadic $q$. This is another encoding of the same gate, with a larger
witness space but a smaller ambient representation. It does not decide
any instance.
\end{proposition}
\begin{proof}
The crossing isotypic projectors commute with simultaneous permutation
of all $n$ tensor copies and therefore preserve $\mathcal V_{q,n}$.
They are rational orthogonal projectors. A solution lies in both high
crossing subspaces and has nonzero EPR overlap, so the rank-one
projection argument of Lemma~\ref{pg:op-signed} proves the first
alternative; no explicit partial-transpose symmetry is needed here.

For the second alternative, let $\mathcal L_b$ and $\mathcal L_c$ be
the low-projector images inside $\mathcal V_{q,n}$. Inconsistency means
that $\Psi$ is orthogonal to
$\ker P_{b,<h}\cap\ker P_{c,<h}$. Thus
\[
 \Psi\in\mathcal L_b+\mathcal L_c.
\]
All spaces and the inner product have rational matrices, so a
decomposition $\Psi=S+U$ with $S\in\mathcal L_b$ and
$U\in\mathcal L_c$ is obtained by rational Gaussian elimination.
On the row space of the $b$ flattening, the low projector has rank
\[
 \sum_{\substack{\ell(\nu)\leq q^2\\H(\nu/n)<h}}
     \dim S_\nu(\mathbb C^{q^2})\,d_\nu
 \leq(n+1)^{B_q}2^{nh},
\]
using the same Weyl, hook, and partition bounds as above. Since the row
space of $S$ lies in this image, this bounds its flattening rank. The
same argument applies to $U$ across $c$. Regrouping $n$ copies identifies
$C_q^{\otimes n}$ with $C_{q^n}$, giving the claimed two-cut
certificate.

If the square is a price, the finite bound of
Proposition~\ref{rt:sos} forbids (PG-OP-compressed-dual) at the prescribed
size: $q^{n(q)}\geq L_{n(q)}\geq2^{m_{n(q)}}$, and
(PG-OP-integer-size) already makes the lower bound at that smaller power
of two exceed $(n+1)^{B_q}q^{5n/4}$. Thus every prescribed compressed
system is consistent. If the square is absent, $r(C_2)=2$ supplies some
dyadic $q$ with $p_0(C_q)<q^{5/4}$, which rules out a solution for any
$n$, including the prescribed one. This proves the equivalence.

For an explicit rational encoding, use the averaged orbit tensors
$e_\alpha$, indexed by occupation vectors
$\alpha\in\mathbb N^{q^4}$ with $|\alpha|=n$. The tensor $e_\alpha$ is
the average of the $n!/\prod_a\alpha_a!$ distinct words of that type.
Its Gram entry is $\prod_a\alpha_a!/n!$. If $v=\sum_\alpha v_\alpha
e_\alpha$, the last equation of (PG-OP-compressed-system) is simply
\[
 \sum_{\substack{|\alpha|=n\\
 \operatorname{supp}\alpha\subseteq
 \{(a,b,a,b):1\leq a,b\leq q\}}}v_\alpha=1.
\]
Thus even the normalization row has only zero and one entries. This
proves the rational encoding claim. Finally, the argument itself gives
another proof of the square-present branch in
Theorem~\ref{pg:op-complete}: for each fixed $q$ and
$h<\sqrt2\log_2q$, the two-cut lower bound eventually exceeds
$(n+1)^{B_q}2^{nh}$. It avoids the intermediate sum-of-squares
factorization, while Proposition~\ref{pg:op-sos-duality} remains a
separate exact finite implication.
\end{proof}

\begin{proposition}[PG-OP-CAPELLI: exact construction without permutation enumeration]
\label{pg:op-capelli}
For fixed $q$, the two projectors in (PG-OP-compressed-system), and the
consistency of that system at $h=(5/4)\log_2q$, can be computed with a
number of exact bit operations polynomial in $n$. The degree and
constants depend on $q$ and are large. This is not a polynomial bound in
$\log n$, nor a uniform polynomial bound when $q$ varies.
\end{proposition}
\begin{proof}
Put $d=q^2$. For one of the two crossings, identify the $q^4=d^2$
variables with a $d\times d$ matrix $(x_{ia})$. In the averaged-orbit
basis, identified with the degree-$n$ monomials, the left
$\mathfrak{gl}_d$ action is
\[
 E_{ij}=\sum_{a=1}^d x_{ia}\frac{\partial}{\partial x_{ja}}.
 \tag{PG-OP-differential}
\]
The other crossing is obtained by permuting the four physical indices
before using the same construction.

Use the Capelli central elements
\[
 \mathcal C_r=
 \sum_{1\leq i_1<\cdots<i_r\leq d}
 \operatorname{column\!\det}
 \bigl(E_{II}+\operatorname{diag}(r-1,r-2,\ldots,0)\bigr),
 \qquad 1\leq r\leq d.
 \tag{PG-OP-Capelli}
\]
Here the column determinant multiplies the selected entries in increasing
column order. Centrality of these elements is the standard Capelli
statement recorded in Itoh (2012), \S4.1, equation~(4.1), pp.~13--14.
We give the eigenvalue and separation calculation explicitly, so no
quantum-immanant formula with a different convention is needed.
On highest weight $\mu=(\mu_1,\ldots,\mu_d)$, the eigenvalue is
\[
 c_r(\mu)=\sum_{i_1<\cdots<i_r}
                    \prod_{k=1}^r(\mu_{i_k}+r-k),\qquad c_0(\mu)=1.
 \tag{PG-OP-Capelli-eigenvalues}
\]
Indeed, in a nonidentity determinant term let $k$ be the largest moved
column. Then its selected row is smaller than $k$, so that factor is a
raising operator. All factors to its right are diagonal and it
annihilates a highest-weight vector. Only the identity permutation
survives, giving the displayed product.

These $d$ eigenvalues distinguish all partitions of $n$ into at most
$d$ parts. To see this directly, with
$z^{\underline s}=z(z-1)\cdots(z-s+1)$, one has
\[
 \sum_{r=0}^d(-1)^{d-r}z^{\underline{d-r}}c_r(\mu)
       =\prod_{i=1}^d(\mu_i+d-i-z).
 \tag{PG-OP-Capelli-characteristic}
\]
For a proof, separate the subsets in (PG-OP-Capelli-eigenvalues)
according as they contain $d$. This gives
\[
 c_r^{(d)}(\mu)=c_r^{(d-1)}(\mu_1,\ldots,\mu_{d-1})
       +\mu_d\,c_{r-1}^{(d-1)}(\mu_1+1,\ldots,\mu_{d-1}+1).
\]
The two contributions to the left side of
(PG-OP-Capelli-characteristic) are respectively $-z$ and $\mu_d$
times $\prod_{i<d}(\mu_i+d-i-z)$, by induction. This proves the
identity. Its roots $\mu_i+d-i$ are strictly decreasing, so their
multiset determines the partition.

For distinct partitions $\nu,\mu\vdash n$ with at most $d$ rows, choose
the first $r=r(\nu,\mu)$ with $c_r(\nu)\ne c_r(\mu)$. The exact
projector onto cut type $\nu$ is
\[
 P_\nu=
 \prod_{\substack{\mu\vdash n,\,\ell(\mu)\leq d\\\mu\ne\nu}}
 \frac{\mathcal C_{r(\nu,\mu)}-c_{r(\nu,\mu)}(\mu)I}
      {c_{r(\nu,\mu)}(\nu)-c_{r(\nu,\mu)}(\mu)}.
 \tag{PG-OP-Capelli-projector}
\]
The central operators commute and act by the displayed scalar
characters on each Schur summand; the product is one on type $\nu$
and zero on every other type. Schur--Weyl duality identifies these
summands with the restrictions of the crossing projectors to
$\mathcal V_{q,n}$. Thus summing the appropriate $P_\nu$ constructs the
required low projector. Its orthogonality is with respect to the
averaged-orbit Gram matrix, not necessarily the unweighted coordinate
inner product.

For fixed $d$, (PG-OP-Capelli) is a sum of a fixed number of products of
at most $d$ first-order operators. Their matrices in degree $n$ have
integer entries with $O_d(\log(n+1))$ bits and can be assembled in time
polynomial in $D_{q,n}$. There are at most $(n+1)^d$ partitions, so
(PG-OP-Capelli-projector) uses polynomially many matrix operations.
Each eigenvalue and each nonzero denominator has $O_d(\log(n+1))$
bits. Products of polynomially many such matrices and denominators have
polynomial bit length, as follows by bounding each coefficient by the
product of the operator entry bounds and matrix dimensions. Summing
the projectors and using fraction-free rational elimination therefore
has polynomial bit complexity in $n$ for fixed $q$. The exact entropy
classification is the integer comparison in
Corollary~\ref{pg:op-deterministic}. No $n!$-term character sum or
floating-point spectral threshold is required by this construction.
\end{proof}


\begin{corollary}[PG-OP-RECTANGLE: a fixed local Schur type]
\label{pg:op-rectangle}
The deterministic family in Corollary~\ref{pg:op-deterministic} remains
valid if its maximum-dimension choice is replaced by the rectangle
\[
 \lambda=\bigl((n(q)/q)^q\bigr).
 \tag{PG-OP-rectangle}
\]
This is a separate prescribed family; consistency of one rectangle
instance does not assert consistency of the earlier maximum-dimension
instance.
\end{corollary}
\begin{proof}
First, the prescribed integer inequality implies $n(q)>2q^2$. Indeed,
$L_n\leq q^n$ and $k_n<m_n\leq n\log_2q$, so it implies
\[
 3n\log_2q>20B_q\log_2(n+1)
                 \geq10q^2(q^2+1).
\]
For $n\leq2q^2$ this would require
$6\log_2q>10(q^2+1)$, which is false for $q\geq2$.
Both $n(q)$ and $q$ are powers of two, so in particular
$2q^2\mid n(q)$.

Write $n=qm$. The balanced multinomial coefficient is maximal among all
$q$-tuples of nonnegative integers of sum $n$, hence
$n!/(m!)^q\geq q^n/(n+1)^q$. The hook formula for the rectangle gives
\[
 d_{(m^q)}=
 \frac{n!\prod_{j=0}^{q-1}j!}{\prod_{j=0}^{q-1}(m+j)!}
 \geq\frac{n!/(m!)^q}{(n+1)^{q(q-1)/2}}
 \geq\frac{q^n}{(n+1)^{A_q}}.
\]
The first inequality uses the $q(q-1)/2$ factors $m+t$ in the
factorial ratios and the fact that each is at most $n+1$; the numerator
factorials are at least one. Thus every lower-dimension estimate used in
Corollary~\ref{pg:op-deterministic} holds for this fixed rectangle, and
the same proof applies.
\end{proof}

\begin{proposition}[PG-OP-DETERMINANT: explicit full-space common vectors]
\label{pg:op-determinant}
Let $d=q^2$, and let $D_b(X)$ and $D_c(X)$ be the determinants of the
$d\times d$ flattenings $12|34$ and $14|23$ of a four-party variable
tensor $X$. For any integer $a\geq1$, the polynomial
\[
 f_{q,a}(X)=\bigl(D_b(X)D_c(X)\bigr)^a,
 \qquad n=2ad,
 \tag{PG-OP-det-vector}
\]
has nonzero evaluation at $C_q$. Its symmetric tensor lies in both
crossing high spaces at threshold
\[
 h_d=H\!\left(\frac{d+1,1^{d-1}}{2d}\right)
     =\log_2(2d)-\frac{d+1}{2d}\log_2(d+1).
 \tag{PG-OP-det-entropy}
\]
Consequently
\[
 p_0(C_q)\geq 2^{h_d}
          =\frac{2q^2}{(q^2+1)^{(q^2+1)/(2q^2)}}
          =2q\,(1-o(1)).
\]
The same polynomial, with $a=n(q)/(2q^2)$, explicitly solves the
compressed prescribed systems at $q=2,4,8$. It also supplies a witness
in the rectangular Schur family of Corollary~\ref{pg:op-rectangle}.
The next proposition extends these explicit solutions further. No
assertion that the $q=16$ system is inconsistent is made.
\end{proposition}
\begin{proof}
The $b$ flattening of $C_q$ is the identity, while its $c$ flattening is
the permutation matrix interchanging the two $q$-dimensional indices.
Thus
\[
 D_b(C_q)=1,\qquad D_c(C_q)=(-1)^{q(q-1)/2},
\]
and the evaluation of $f_{q,a}$ is a sign. Identify homogeneous
polynomials with symmetric tensors in the averaged-orbit basis of
Proposition~\ref{pg:op-compressed}; this evaluation is exactly the EPR
overlap used there.

For the $b$ crossing, multiplication by $D_b^a$ twists every
$\mathrm{GL}_d$ Schur type of the remaining degree-$ad$ polynomial
$D_c^a$ by $\det^a$. Every resulting partition of $n=2ad$ therefore has
all $d$ parts at least $a$. Among such partitions the vector
$(a(d+1),a^{d-1})$ majorizes all others. Schur concavity of entropy gives
the lower bound (PG-OP-det-entropy). The same argument with $b,c$
interchanged proves both high-support conditions. Thus the common
vector has nonzero EPR overlap, proving the functional bound.

The displayed minimum entropy is actually attained by this polynomial's
cut support. Substitute a rank-one $b$ flattening
$X_{ab,cd}=u_{ab}v_{cd}$. Its $c$ determinant is, up to sign,
$(\det u)^q(\det v)^q$, which is not the zero polynomial. Therefore
$D_c^a$ has a nonzero $b$ component of one-row type $(ad)$; the
coordinate ring of the rank-one matrices in degree $ad$ is
$\operatorname{Sym}^{ad}(\mathbb C^d)\otimes
\operatorname{Sym}^{ad}(\mathbb C^d)$. Twisting by $D_b^a$ gives the
extreme partition $(a(d+1),a^{d-1})$. Hence this particular family does
not conceal a higher minimum support entropy.

Under any one physical $\mathrm{GL}_q$ action, each flattening
determinant transforms by $(\det g)^q$. The polynomial $f_{q,a}$
therefore has local type $( (2aq)^q )=( (n/q)^q )$ at every party.
In this rectangle the local Schur multiplicity space
$S_{(n/q)^q}(\mathbb C^q)$ is one-dimensional, so the polynomial gives
a vector in the corresponding four-Specht invariant space. Averaging
it with the $B,D$ interchange and the simultaneous $A,C$ and $B,D$
interchange preserves both high supports and its EPR overlap, since
$C_q$ is fixed by both permutations. The latter permutation is total
transpose under the matrix identification. Hence the averaged vector
has a real symmetric, partial-transpose-fixed matrix representative
with nonzero trace, which satisfies the rectangular version of
(PG-OP-rational-system) after rescaling its trace.

For $q=2^N$, the comparison $h_d>(5/4)\log_2q$ is exactly
\[
 2^{(4+3N)d}>(d+1)^{2(d+1)}.
 \tag{PG-OP-det-integer}
\]
Integer arithmetic verifies it for $q=2,4,8$. At the prescribed sizes,
the exponents $a$ are respectively $256,256,1024$, and all three
polynomials evaluate to $1$ at $C_q$, so no normalization denominator
is needed. The checker \path{programs/pg_verify_determinant_gate.py}
and output \path{outputs/pg_determinant_gate_verification.json}
preserve the exact powers in (PG-OP-det-integer), verify divisibility,
recheck the prescribed-size inequalities, and reconstruct both
flattening permutation signs.

For $q=16$ the reverse strict inequality holds, so this particular
polynomial does not reach the prescribed threshold $5$. More generally
$h_{q^2}<\log_2q+1$, and for all dyadic $q\geq16$ this is at most
$(5/4)\log_2q$. Thus taking further powers of this determinant product
does not solve the remaining systems. The bound tends to $2q$, whose
$N$th root at $q=2^N$ tends to $2$; no regularized gap follows from this
explicit family alone.
\end{proof}


\begin{proposition}[PG-OP-SLICE-DETERMINANT: six prescribed instances]
\label{pg:op-slice-determinant}
Define the slice-minor product
\[
 W_q(X)=\prod_{i=1}^q\det\bigl(X_{i,b,i,d}\bigr)_{1\leq b,d\leq q},
 \qquad
 F_q(X)=D_b(X)^{21}D_c(X)^{21}W_q(X)^{22}.
 \tag{PG-OP-slice-polynomial}
\]
Put $d=q^2$, $n_0=64d$, and
\[
 \eta_q=
 \bigl(21(d+1)+22q,\ (21+22q)^{q-1},\ 21^{d-q}\bigr)
 \vdash n_0.
 \tag{PG-OP-slice-shape}
\]
Then $F_q(C_q)\ne0$ and every crossing Schur type occurring in $F_q$
has entropy at least $H(\eta_q/n_0)$, for both crossings. The same
holds for $F_q^t$ with scaled type $t\eta_q$.

For each $q=2,4,8,16,32,64$, a suitable power of $F_q$ gives an explicit
solution of the prescribed compressed system. Its projection to the
local rectangle also has nonzero EPR overlap and solves the prescribed
rectangular system. Therefore the first unresolved dyadic parameter in
these two families is now $q=128$. No statement of inconsistency at
$q=128$, or of feasibility for the earlier maximum-dimension Specht
family, is implied.
\end{proposition}
\begin{proof}
Each slice determinant is a $q\times q$ minor for both the $b$ and $c$
flattenings, up to transposition. It is a polynomial of Cauchy type
$(1^q)$ for either cut (Howe \& Lee, 2012, exterior Pieri rule,
p.203). Tensor-product highest-weight dominance therefore
bounds every type in $W_q^{22}$ by $( (22q)^q )$. Indeed, every weight
of a highest-weight module is its highest weight minus a nonnegative
integer combination of positive roots. A highest weight occurring in
a tensor product is a sum of weights of its factors, and is therefore
dominated by the sum of their highest weights. This proves the needed
dominance consequence directly; compare the Littlewood--Richardson
rule in Howe and Lee (2012, \S4, p.202).

For cut $b$, every type in $D_c^{21}$ is a partition of $21d$, hence
is dominated by $(21d)$. Multiplication by $D_b^{21}$ adds $21$ to all
$d$ parts. Applying the same highest-weight dominance once more gives
precisely the upper partition $\eta_q$ in (PG-OP-slice-shape). Since
entropy is Schur concave, all occurring types have entropy at least
$H(\eta_q/n_0)$. The argument is symmetric in $b,c$. Taking powers
scales the displayed dominating partition, preserving its normalized
entropy. At $C_q$, every slice matrix is the identity, while the two
full determinants have the signs computed above; hence $F_q(C_q)$ is
a sign and is nonzero.

The rectangle assertion requires more than balance of the diagonal
torus weights. Let
\[
 \operatorname{res}f(A,B)=f(A\otimes B),\qquad
 (A\otimes B)_{ab,cd}=A_{ac}B_{bd}.
\]
This restriction of polynomials is equivariant for all four local
$\mathrm{GL}_q$ actions. For arbitrary nonnegative integers $a,\ell$,
it gives
\[
 \operatorname{res}(D_b^aD_c^aW_q^\ell)
 =\pm\det(A)^{2aq}
        \left(\prod_{i=1}^q A_{ii}^{\ell q}\right)
        \det(B)^{(2a+\ell)q}.
 \tag{PG-OP-slice-restriction}
\]
In applying the formula below take $a=21t$ and $\ell=22t$.

Project the polynomial to the rectangle
$\lambda=((n/q)^q)$ at each physical party, where
$n=(2a+\ell)q^2$. Restriction commutes with these isotypic projections.
The $B$ factor in (PG-OP-slice-restriction) already has the required
one-dimensional determinant type. On the $A$ factor, multiplication by
$\det(A)^{2aq}$ is an injective determinant twist. Thus the desired
component is nonzero exactly when the projection of
$\prod_i A_{ii}^{\ell q}$ onto $\det(A)^{\ell q}$ is nonzero. In the
unitarily invariant Fischer inner product on matrix polynomials,
defined on monomials by
$\langle x^\alpha,x^\beta\rangle=\delta_{\alpha\beta}\alpha!$
(M\"ollers, 2013, p.811, immediately before Eq.(1.16)),
\[
 \left\langle\det(A)^{\ell q},\prod_i A_{ii}^{\ell q}\right\rangle
       =\bigl((\ell q)!\bigr)^q>0,
\]
because the diagonal monomial has coefficient one in the determinant
power. Consequently the restricted projected polynomial is a nonzero
multiple of
$\det(A)^{n/q}\det(B)^{n/q}$, and its evaluation at $A=B=I$ is
nonzero. The physical isotypic projections commute with both crossing
projectors, so the entropy support conditions survive. Averaging total
transpose and the $B,D$ interchange as in
Proposition~\ref{pg:op-determinant}, and normalizing the trace, gives
the required rectangular witness. All these polynomial and projector
operations are rational.

It remains to check the prescribed sizes and entropy thresholds. For
$q\geq8$, the inequality $n(q)\leq64q^2$ would contradict
\[
 3n(q)\log_2q>10q^2(q^2+1),
\]
because $192\log_2q<10(q^2+1)$ for dyadic $q\geq8$; the latter
inequality holds at $q=8$ and persists upon doubling. For $q=2,4$ the
already computed sizes are also greater than $64q^2$. Thus
$64q^2\mid n(q)$ in every case, and one can take
$t=n(q)/(64q^2)$.

For $q=2^N$, let $u=21(d+1)+22q$ and $v=21+22q$.
The exact entropy comparison at $n_0=64d$ is
\[
 2^{\,n_0(\log_2 n_0-5N/4)}
       >u^u\,v^{v(q-1)}\,21^{21(d-q)}.
 \tag{PG-OP-slice-integer}
\]
Its exponent is an integer. The checker
\path{programs/pg_verify_decorated_determinants.py} verifies this strict
integer inequality for $N=1,\ldots,6$; its output
\path{outputs/pg_decorated_determinants_verification.json} records the
exact exponent, right-side bit length and hash, shape, and evaluation
sign. Independent reconstruction of the partition and exact integer
comparison appears in \path{programs/ex_verify_decorated_entropy.py},
with result \path{outputs/ex_decorated_entropy_independent.json}.
A second independent checker,
\path{programs/rt_verify_decorated_entropy.py}, certifies the same signs
using rational upper and lower logarithm bounds; its result is
\path{outputs/rt_decorated_entropy_verified.json}.
At $q=64$ the certified entropy is approximately $7.5059004$,
strictly greater than $7.5$. For $q=128$ this particular dominance bound
is approximately $8.5401358$, below $8.75$; that failure of the bound
is not a proof that the polynomial or the system fails.
The Fischer argument proves that some crossing components survive the
physical rectangle projection. It does not identify those components,
and in particular does not assert that the extremal dominating type
$\eta_q$ survives that projection.

The family still gives only a constant-factor raw improvement. In fact,
\[
 H(\eta_q/(64q^2))
   =\log_2q+H(21/64,22/64,21/64)+o(1),
\]
so the resulting lower bound is $(c-o(1))q$, with
$c=2^{H(21/64,22/64,21/64)}<3$. Its grouped-dimension root therefore
tends to $2$. A superconstant improvement remains necessary to select
the square-present branch by this construction.
\end{proof}

\paragraph{PG-OP-TOP-PARITY: a limitation of the projection argument.}
The top crossing component need not survive the physical rectangle
projection. For example, let $q$ and $\ell$ both be odd, and consider
$W_q^\ell$ at the $b$ crossing. Put $U=A\otimes B$ and
$p_i=\bigwedge_{j=1}^q(e_i\otimes e_j)\in\bigwedge^q U$.
In the highest Cauchy component
$\beta=((\ell q)^q)$, the row vector is proportional to the Cartan
product $\prod_i p_i^\ell$; multiplication into this component factors
through the separate row and column Cartan maps. Permutation matrices
on $A$ permute the $p_i$ without changing the order of the $B$ wedge
factors. Commutativity of the Cartan product therefore fixes this row
vector. But its required $A$ rectangle has character
$\det_A^{\ell q}$, on which a transposition acts as $-1$.
Equivariant projection to that rectangle consequently vanishes.
This argument is over characteristic zero and gives no conclusion
for dyadic $q$. It illustrates why nonzero total rectangle projection,
or a positive rectangular Kronecker multiplicity, cannot replace a
check of the particular crossing component.

\begin{proposition}[PG-OP-PT-FINITE: exact limits of six central tests]
\label{pg:op-pt-finite}
For each pair in the following table, let $h_{\mathrm{cent}}$ be the
largest $h$ for which (PG-OP-PT-linear) has a nonzero solution. The listed
values are exact. Here $H(\tau/n)$ has its usual Shannon meaning.
\[
\begin{array}{c|c|c|c|c}
 n&\lambda&\mu&q_{\min}&h_{\mathrm{cent}}\\ \hline
 4&(3,1)&(2,2)&2&H((2,1,1)/4)=3/2\\
 4&(2,2)&(2,2)&2&H((2,2)/4)=1\\
 9&(3,3,3)&(3,3,3)&3&H((3,3,3)/9)=\log_2 3\\
 9&(3,3,3)&(4,3,2)&3&H((5,4)/9)\\
 10&(3,3,2,2)&(3,3,2,2)&4&H((5,4,1)/10)\\
 10&(3,3,3,1)&(3,3,2,2)&4&H((5,5)/10)=1
\end{array}
\]
In particular, none of the four tested multiplicity-bearing blocks gives
an improvement over the baseline $q_{\min}$ by this central positive
test. This is a statement about the specified finite cones. It is not an
upper bound for noncentral positive operators, indefinite common
crossing vectors, other local Schur types, or $p_0$ itself.
\end{proposition}
\begin{proof}
The exact rational certificates contain both a lower witness $x$ and an
upper witness $y$. For a proposed boundary partition $\tau$, define
\[
 J=\{\nu:g(\lambda,\mu,\nu)>0,
                    \ H(\nu/n)>H(\tau/n)\}.
\]
The lower witness satisfies $x\geq0$, $Dx=0$,
$\sum_\nu d_\nu m_\nu x_\nu=1$, and has minimum support entropy exactly
$H(\tau/n)$. The integer upper witness satisfies
\[
 (Dy)_\nu>0\qquad(\nu\in J).
 \tag{PG-OP-PT-dual}
\]
If $J$ is empty the upper assertion is immediate. Otherwise any nonzero
$x'\geq0$ supported in $J$ would satisfy
$y^TDx'=\sum_{\nu\in J}(Dy)_\nu x'_\nu>0$, contradicting $Dx'=0$.
Since the constituent set is finite, this proves the exact optimum.
No floating-point entropy comparison is needed: for fixed $n$,
\[
 H(\nu/n)>H(\tau/n)
 \quad\Longleftrightarrow\quad
 \prod_i\nu_i^{\nu_i}<\prod_i\tau_i^{\tau_i}.
\]

The discovery program \path{programs/pg_op_pt_central.py} constructs the
character table by Murnaghan--Nakayama, verifies character orthogonality,
and reconstructs $K,D$ from (PG-OP-PT-character). Its integer permutation
enumerator is \path{programs/pg_op_pair_counts.cpp}. It uses one canonical
representative for each conjugacy class of $g$, all permutations $h$,
and their exact class-size weights. The enumeration checks its total
$n!$, both identity-class specializations, and the symmetry obtained by
interchanging $g,h$. The two $n=9$ calculations each evaluate
$30\cdot9!=10{,}886{,}400$ permutation pairs; each $n=10$ calculation
evaluates $42\cdot10!=152{,}409{,}600$ pairs.

The separate program \path{programs/pg_verify_op_pt_duals.py} reconstructs
all character entries by the Jacobi--Trudi determinant in the power-sum
basis, using $h_k=\sum_{\rho\vdash k}p_\rho/z_\rho$. Thus its character
calculation does not reuse the rim-hook recursion. It recomputes $K,D$
over $\mathbb Q$, verifies the lower witnesses, and checks the integer
dual inequalities exactly. A numerical linear program is used only to
find candidate vectors; its infeasibility reports are not used as proof.

The six input/output stems are
\path{pg_op_pt_n4_31_22}, \path{pg_op_pt_n4_22_22},
\path{pg_op_pt_n9_333_333}, \path{pg_op_pt_n9_333_432},
\path{pg_op_pt_n10_3322_3322}, and
\path{pg_op_pt_n10_3331_3322}.
For each stem, \path{outputs/}\emph{stem}\path{.json} contains the full
character and Gram data and the lower witness;
\path{outputs/}\emph{stem}\path{_dual_verified.json} contains the input
hash, boundary partition, integer dual, and every coordinate of $Dy$.
Independent reconstruction of the group counts was completed in
\path{programs/ex_verify_pt_word_counts.cpp} and its Python wrapper,
using Heap permutations and reversed cycle representatives. The output
\path{outputs/ex_pt_word_counts_independent.json} checks all six cases.
A further root reconstruction and exact primal/dual audit is preserved in
\path{outputs/rt_pt_certificates_verified.json}. Thus the result does
not rely on one shared enumeration implementation.
\end{proof}


\paragraph{An exact distinction between positive and signed central tests.}
For $n=9$ and $\lambda=\mu=(3,3,3)$, a signed central witness has trace
one and minimum support entropy $H((4,3,1,1)/9)>\log_2 3$, whereas
Proposition~\ref{pg:op-pt-finite} proves that the positive central optimum
is only $\log_2 3$. Its nonzero coefficients $x_\nu$ are
\[
\begin{gathered}
\begin{array}{c|rrrr}
\nu&(4,3,1,1)&(4,2,2,1)&(4,2,1^3)&(4,1^5)\\ \hline
x_\nu&1/1512&-1/432&-2/189&-3/112
\end{array}\\[5pt]
\begin{array}{c|rrrr}
\nu&(3,3,1^3)&(2^3,1^3)&(2^2,1^5)&(1^9)\\ \hline
x_\nu&1/140&3/112&23/378&3/14
\end{array}
\end{gathered}
\]
The exact certificate is
\path{outputs/pg_op_pt_n9_333_333_signed_verified.json}; the checker
is \path{programs/pg_verify_op_pt_signed.py}. It verifies $Dx=0$ and
$\sum_\nu d_\nu m_\nu x_\nu=1$. It also gives a rational vector $y$
with $(Dy)_\nu=d_\nu m_\nu$ on every strictly higher-entropy constituent.
For a signed kernel vector supported there, pairing with $y$ would force
its trace to zero. This proves the signed central optimum as well.
The five other tabulated blocks have the same signed and positive
optima, with corresponding certificates. The $n=9$ signed gain is a
finite scope distinction and is weaker than the earlier standard-type
raw bound at $q=3$; it does not itself imply a regularized gain.

\paragraph{What a uniform operator construction must change.}
For arbitrary local types with row counts $a=\ell(\lambda)$ and
$b=\ell(\mu)$, every constituent has at most $ab$ rows, by the
Schur--Weyl argument used above. Consequently any common-crossing
threshold inside this block is at most $\log_2(ab)$. In particular,
if $a,b\leq q$ and $\min(a,b)=q^{o(1)}$ along grouped dimensions, this
construction can give at most $q^{1+o(1)}$. A proposed uniform
$q^{1+\varepsilon}$ lower bound therefore needs both local row counts
at least $q^\varepsilon$, as well as $n\geq q^{1+\varepsilon}$.
This necessary condition applies independently of centrality or
positivity. It explains why increasing only the rectangle height while
keeping a bounded-row companion cannot resolve the regularized gate.
No general entropy ceiling has been proved for central positive
partial-transpose-fixed operators with both row counts growing.
Theorem~\ref{pg:op-complete} resolves the uniform construction
conditionally by using the full symmetric commutant and allowing
indefinite operators; it does not establish such a ceiling.

\subsection{A universal cycle-cover formulation with two possible values}

The operator functional above is one upper functional with prescribed
side-wire bounds. Taking all covers by side-wire products gives a
universal upper functional with those bounds. A normalization argument
is essential for using a square price in this formulation.

\begin{lemma}[PG-CYCLE-NORMALIZATION]
\label{pg:cycle-normalization}
Let $G_q$ be the product of the four dimension-$q$ wires on the
cycle $12,23,34,41$. Over every field, $\widetilde Q(G_2)=4$.
Consequently, every ordinary normalized
four-party spectral point $\psi$ satisfies
\[
 \sum_{e\in\{12,23,34,41\}}d_\psi(e)\geq2.
\]
If the unit-side Euclidean square is a price $c\,d_\psi$, then
$c\leq2$. In particular, its rescaling to side length $1/2$ is
itself the wire metric of the enlarged spectral point $\psi^{c/2}$.
\end{lemma}
\begin{proof}
The minimum-cut formula is Vrana and Christandl (2017),
Theorem~1, arXiv version~2, p.~3; cycles are discussed explicitly
on p.~6. We give the field check for the four-cycle, since its use
over $\mathbb F_2$ cannot rely only on the quantum interpretation.
Number the cyclic edges $1,2,3,4$ and take integer vectors
\[
 c_1=(1,0),\quad c_2=(1,1),\quad
 c_3=(0,1),\quad c_4=(1,-1).
\]
Nonincident edges have orthogonal vectors. For labels
$0\leq i_e<q$, the value of $\sum_e c_e i_e$ ranges over at
most $9q^2$ integer pairs. Some fiber therefore contains at least
$q^2/9$ labelings. Fix its value $g$. The integer energy
$\|\sum_e c_e i_e-g\|^2$ is a sum of functions of the labels
incident to individual vertices: the only cross terms involve
incident edges. Local diagonal Laurent-monomial maps consequently
degenerate $G_q$ to the sum over this fiber, precisely as in the
energy construction of Vrana and Christandl, equations~(11)--(14),
p.~4, and Theorem~2, p.~5. At each vertex its two incident labels
determine the other two, because their two coefficient vectors
are linearly independent over $\mathbb R$. Thus the fiber tensor
is a diagonal tensor of size $M\geq q^2/9$.

This is a coefficient-$1$ monomial degeneration over every field.
Its integer exponents and combinatorial label selections are not
reduced modulo the field characteristic. For completeness, fix $q$
before using degeneration monotonicity: after a common Laurent
shift, the desired tensor is a fixed leading coefficient of the
product of four polynomial local maps. In its $N$th tensor power,
that coefficient is a sum of $O_q(N^3)$ restrictions of
$G_q^{\otimes N}$, by expanding the four grouped polynomial maps
and summing their coefficient degrees. The cover inequality and
$N$th roots prove $\psi(G_q)\geq M$ without any restriction on
field size. Equivalently, the polynomial diagonal overhead can be
absorbed in $o(N)$ extra copies using the spanning-tree
$\operatorname{GHZ}_2$ restriction. Only after this fixed-$q$
argument take $q=2^m$ and let $m\to\infty$. It gives
$\widetilde Q(G_2)\geq4$; the singleton flattening rank gives the
reverse inequality. This ordering of limits avoids using an
exponentially growing interpolation degree in one step.

Now $\log_2\psi(G_2)=\sum_e d_\psi(e)\geq2$.
For a unit-side square $c\,d_\psi$ the left side is $4/c$, so
$c\leq2$. For $0<\gamma\leq1$, raising a normalized spectral
point to the power $\gamma$ preserves multiplicativity,
restriction monotonicity and the cover inequality, since
$(\sum_j x_j)^\gamma\leq\sum_j x_j^\gamma$.
Taking $\gamma=c/2$ proves the last assertion.
\end{proof}

For $\mathbf k=(k_{12},k_{23},k_{34},k_{41})\in\mathbb N_0^4$,
put
\[
 S_{\mathbf k}=\prod_{e\in\{12,23,34,41\}}
 \operatorname{EPR}_{e,2^{k_e}},\qquad
 w(\mathbf k)=2^{\frac12\sum_e k_e}.
\]
The empty product is the unit tensor. Define
\[
 u_F(T)=\inf\left\{\sum_{j=1}^s w(\mathbf k_j):
           (S_{\mathbf k_j})_{j=1}^s\text{ covers }T\right\},
 \qquad
 U_F(T)=\inf_{N\geq1}u_F(T^{\otimes N})^{1/N}.
 \tag{PG-CYCLE-cover}
\]
The infima include arbitrary linear restrictions in the definition
of a cover. No assertion about attainment of the first infimum is
required.

\begin{theorem}[PG-CYCLE-GAP]
\label{pg:cycle-gap}
Assume the checked price-space and square-gate chain. Over every
field $F$, $u_F$ is a normalized upper functional, and
\[
 U_F(C_2)=
 \begin{cases}
 2,&\text{if the Euclidean square is not a price over }F,\\
 2^{\sqrt2},&\text{if the Euclidean square is a price over }F.
 \end{cases}
 \tag{PG-CYCLE-two-values}
\]
In particular, there are no intermediate values. This theorem does
not decide which value occurs over $\mathbb C$ or $\mathbb F_2$.
\end{theorem}
\begin{proof}
Rank-one decompositions cover every tensor by finitely many unit
tensors, so $u_F$ is finite. Empty covers give $u_F(0)=0$, and
every nonempty covering term has cost at least one, giving
$u_F(1)=1$. Composing restrictions and concatenating covers prove
restriction monotonicity and the cover inequality. Tensoring
two covers gives products $S_{\mathbf k}S_{\mathbf l}
=S_{\mathbf k+\mathbf l}$ with multiplying costs, so
$u_F$ is submultiplicative. The uniform upper support functional
$\zeta^\theta$, $\theta_i=1/4$, has value
$w(\mathbf k)$ on each $S_{\mathbf k}$ by the graph calculation
in Proposition~\ref{pg:star}. It is below $u_F$ by its cover
inequality. Since $\zeta^\theta(\operatorname{GHZ}_r)=r$,
while the rank-one cover has cost $r$, we also have
$u_F(\operatorname{GHZ}_r)=r$.

An enlarged spectral point $\varphi$ is below $u_F$ if and only if
its four side-wire values are at most $\sqrt2$. The forward
implication follows by the one-term wire covers. For the reverse
implication, multiplicativity bounds each
$\varphi(S_{\mathbf k})$ by $w(\mathbf k)$, and the cover
inequality gives $\varphi(T)\leq u_F(T)$. Such a point is also
below $U_F$, by taking $N$th roots. Conversely, the power
regularization is an upper functional, and
Lemma~\ref{pg:saturation} produces a spectral point below it
which attains its value at any prescribed tensor. Hence
\[
 U_F(C_2)=\max\left\{\varphi(C_2):
 d_\varphi(12),d_\varphi(23),d_\varphi(34),d_\varphi(41)
 \leq\tfrac12\right\}.
 \tag{PG-CYCLE-dual}
\]

CAT(0) roundness bounds the sum of the two diagonal lengths by
$\sqrt2$, so the maximum is at most $2^{\sqrt2}$.
A crossing cut metric scaled by $1/2$ gives value $2$ and meets
all four side bounds. If the square is absent, every price is a
tree metric by Theorem~\ref{pg:square-gate}. The tree four-point
condition bounds the diagonal sum by one, giving $U_F(C_2)=2$.

If the square is present, the finite-price normalization theorem
writes its unit-side metric as $c\,d_\psi$ with $\psi$ an
ordinary normalized spectral point. Lemma~\ref{pg:cycle-normalization}
shows that $\varphi=\psi^{c/2}$ is an allowed enlarged spectral
point. Its sides have length $1/2$ and its diagonals have length
$1/\sqrt2$, so $\varphi(C_2)=2^{\sqrt2}$. This attains the
universal upper bound and proves the second case.
\end{proof}

\begin{corollary}[PG-CYCLE-FINITE: a finite algebraic certificate]
\label{pg:cycle-finite}
Under the hypotheses of Theorem~\ref{pg:cycle-gap}, the square is
absent if and only if, for some $m\geq1$ and $q=2^m$, there is a
finite cover of $C_q$ by the resources $S_{\mathbf k_j}$ satisfying
\[
 \sum_j2^{\frac12\sum_e k_{e,j}}<q^{\sqrt2}.
 \tag{PG-CYCLE-finite}
\]
Equivalently, such a cover exists if and only if a planar strict
shortening exists. No geometric auxiliary vertices occur in this
certificate.
\end{corollary}
\begin{proof}
If the cover exists, then
$U_F(C_2)\leq u_F(C_{2^m})^{1/m}<2^{\sqrt2}$, so the
two-values theorem selects $U_F(C_2)=2$ and excludes the square.
Conversely, in the absence of the square,
$u_F(C_{2^m})^{1/m}\to2<2^{\sqrt2}$, which gives such a cover
for sufficiently large $m$. The equivalence with planar shortening
was proved in the fixed-dimensional square gate.

The costs in a finite cover belong to
$\mathbb N_0+\sqrt2\,\mathbb N_0$, a locally finite subset of
$[0,\infty)$. Thus the infimum defining $u_F(T)$ is attained.
For fixed $m$ and the displayed budget, there are only finitely
many resource formats: each term costs at least one and each
total edge exponent is less than $2m\sqrt2$.
For each format, the existence of the local maps is a finite
system of polynomial equations with integer coefficients,
obtained by equating the entries of the sum with $C_q$.
This describes a concrete exhaustive algebraic search, without
asserting a bound on the required $m$ or practical feasibility.
\end{proof}

\begin{proposition}[PG-CYCLE-FINITE-LOWER]
\label{pg:cycle-finite-lower}
Over every field and for every integer $q\geq2$,
\[
 u_F(C_q)\geq\frac{q(3q-1)}{q+1}.
 \tag{PG-CYCLE-rank-lower}
\]
Consequently no certificate (PG-CYCLE-finite) exists at dyadic
width $m=1,2,3$. The first dyadic width not excluded by this
bound is $m=4$; this is not an assertion that a certificate exists
there.
\end{proposition}
\begin{proof}
Discard any zero output terms from a cycle cover. Let $r_j,s_j$
be the actual ranks of its $j$th output tensor across the two
crossing cuts. Both are between $1$ and $q^2$, and its resource
cost is at least $\sqrt{r_js_j}$. For $1\leq r,s\leq q^2$, put
$t=\sqrt{\max(r,s)/\min(r,s)}$. Then $1\leq t\leq q$ and
$(t-1)(q-t)\geq0$, giving
\[
 \sqrt{rs}\geq
 \frac{r+s+(q-1)\min(r,s)}{q+1}.
\]
The two crossing ranks of $C_q$ are both $q^2$, so rank
subadditivity gives $\sum_jr_j,\sum_js_j\geq q^2$.
Fixing the same basis coordinate at parties $1,3$ leaves the
$q\times q$ identity matrix on parties $2,4$. Each sliced term
has rank at most both $r_j$ and $s_j$, and consequently
$\sum_j\min(r_j,s_j)\geq q$. Summing the displayed inequality
and taking the cover infimum proves the bound.

At $q=2,4,8$ the lower bounds are respectively $10/3$, $44/5$,
and $184/9$. Since $\sqrt2<10/7$, the thresholds $q^{\sqrt2}$
are respectively less than $3$, $8$, and $20$. For the last
comparison, $2^{30}<20^7$ gives $8^{10/7}<20$.
Thus all three strict certificate inequalities are excluded.
At $q=16$ the displayed lower bound is $752/17$, below
$16^{\sqrt2}$, so it does not exclude that width.
\end{proof}

\begin{lemma}[PG-CROSSCUT-RESIDUAL]
\label{pg:crosscut-residual}
If $U$ has rank at most two across $14|23$, then over every field
\[
 \operatorname{rank}_{12|34}(C_q-U)\geq\frac{q(q-1)}2.
\]
The analogous statement holds with the two crossing cuts exchanged.
The bound is sharp already for a rank-one tensor on the other cut.
\end{lemma}
\begin{proof}
In the $12|34$ flattening, $C_q$ is the identity on the vector
space of $q\times q$ matrices. A tensor of rank at most two
across $14|23$ acts, after harmless transpose conventions, as
\[
 X\longmapsto M_1X^TN_1+M_2X^TN_2.
\]
We show that its fixed space has dimension at most $q(q+1)/2$.
Extend the field to an algebraic closure, which preserves ranks.

First, any two square matrices can be made upper triangular by
simultaneous independent changes of domain and codomain bases.
Indeed, if $M_1$ is singular, take a nonzero $v$ in its kernel;
then $M_1v,M_2v$ span at most a line. If $M_1$ is invertible,
take an eigenvector of $M_1^{-1}M_2$ for the same conclusion.
Choose a codomain line containing both images, arbitrarily if
both are zero. The two maps descend to maps between the quotient
spaces by these domain and codomain lines. Induction on the
common dimension, with each chosen line first in its basis,
makes both matrices upper triangular.

An invertible change $X=P Y Q$ replaces the two $M_\ell$ by
$P^{-1}M_\ell Q^T$, allowing precisely these independent basis
changes. Write the resulting upper triangular matrices as
$A_1,A_2$ and the other two matrices as $B_1,B_2$. The fixed
point equation is
\[
 Y=A_1Y^TB_1+A_2Y^TB_2.
\]
For a strict-lower entry $Y_{ij}$, with $i>j$, its right-hand
side is a linear combination of entries $Y_{tk}$ with $k\geq i$.
If $t\leq k$ this entry is on or above the diagonal. If $t>k$,
it is strictly below the diagonal in row $t>i$. Process the
rows from bottom to top. Every strict-lower entry is thus
determined by the upper triangular entries, including the
diagonal, and entries already determined in larger rows.
Projection onto the upper triangular entries is injective on
the fixed space, whose dimension is therefore at most
$q(q+1)/2$. Rank-nullity proves the assertion. No division by
two is used, so characteristic two is included.

For sharpness, take $M_1=N_1=I$ and $M_2=0$. The fixed space is
exactly the symmetric matrices, of dimension $q(q+1)/2$ in every
characteristic.
\end{proof}

\begin{proposition}[PG-CYCLE-DYADIC-LOWER]
\label{pg:cycle-dyadic-lower}
For every dyadic $q\geq4$, over every field,
\[
 u_F(C_q)\geq q+\frac{5q(q-1)}{2(q+1)}.
 \tag{PG-CYCLE-dyadic-lower}
\]
In particular, no finite cycle certificate exists at $q=16$.
Together with Proposition~\ref{pg:cycle-finite-lower}, this excludes
all dyadic certificate widths $m\leq4$. This proposition alone first leaves open $m=5$, or $q=32$;
the next proposition excludes that case as well.
\end{proposition}
\begin{proof}
For each nonzero output term, let $R,S$ be its two cycle-resource
cut capacities, capped at $q^2$. Its cost is at least $\sqrt{RS}$.
First suppose a resource has $(R,S)=(q^2,1)$, or the reversed
pair. Its cost is at least $q$, and its output $U$ has rank one
across the second cut. After removing this term, the sums of the
actual remaining ranks are at least $q(q-1)/2$ across the first
cut, by Lemma~\ref{pg:crosscut-residual}, and at least $q^2-1$
across the second cut, by rank subadditivity. The fixed-coordinate
slice of $U$ has matrix rank at most one. The remaining slices
therefore sum to a matrix of rank at least $q-1$, giving a lower
bound $q-1$ on the sum of the minima of the two remaining cut
ranks. Apply the term inequality from
Proposition~\ref{pg:cycle-finite-lower} to these remaining
outputs. Their total resource cost is at least
\[
 \frac{q(q-1)/2+(q^2-1)+(q-1)^2}{q+1}
 =\frac{5q(q-1)}{2(q+1)}.
\]
Adding the removed resource proves the stated bound in this case.

Otherwise no resource has either extreme pair. Since $q$ and
the capped capacities are powers of two,
$\max(R,S)/\min(R,S)\leq q^2/2$.
The same elementary term inequality now uses
$h=q/\sqrt2$ in place of $q$:
\[
 \sqrt{RS}\geq\frac{R+S+(h-1)\min(R,S)}{h+1}.
\]
The two capacity sums are at least $q^2$ and the minimum-capacity
sum is at least $q$. Summing yields
\[
 u_F(C_q)\geq q+\frac{2q(q-1)}{q/\sqrt2+1}
 \geq q+\frac{5q(q-1)}{2(q+1)}.
\]
For the last comparison it suffices that
$4(q+1)\geq5(q/\sqrt2+1)$, which holds for $q\geq4$.
This exhausts all resource formats.

At $q=16$, the resulting bound is $872/17>51$.
Meanwhile $\sqrt2<17/12$ and $2^{17}<51^3$ imply
$16^{\sqrt2}<2^{17/3}<51$. Hence the strict finite-cycle
certificate inequality cannot hold at this width. At $q=32$
the bound is below $32^{\sqrt2}$, so no assertion is made there.
\end{proof}

\begin{proposition}[PG-CYCLE-Q32]
\label{pg:cycle-q32}
Over every field, $u_F(C_{32})\geq135>32^{\sqrt2}$.
Consequently every strict finite cycle certificate has dyadic
width $m\geq6$.
\end{proposition}
\begin{proof}
Suppose a cover has cost less than $135$. Cap each resource's
crossing capacities $R,S$ at $1024$; they remain powers of two,
and its cost is at least $\sqrt{RS}$. Call a resource high if
$\max(R,S)/\min(R,S)>256$. Its capped pair belongs to
\[
 (512,1),\ (1024,1),\ (1024,2),\
 (1,512),\ (1,1024),\ (2,1024).
 \tag{PG-Q32-high-types}
\]
All remaining resources have ratio at most $256$. For these the
term inequality with $h=16$ is
\[
 17\sqrt{RS}\geq R+S+15\min(R,S).
 \tag{PG-Q32-low-inequality}
\]

Here is a finite rational check of the high types. For their
multiset $H$, let $H_R,H_S,H_{\min}$ be the respective sums of
$R,S,\min(R,S)$. Let $K_R$ be the maximum sum of $R$ over a
submultiset whose sum of $S$ is at most two; the empty submultiset
is allowed. Define $K_S$ symmetrically. Such a submultiset has at
most two terms, and its output sum has rank at most two on the
small cut. Lemma~\ref{pg:crosscut-residual} applied to this sum
therefore gives the following lower bounds on the low-resource
capacity sums:
\[
 \begin{split}
 L_R&=\max(0,1024-H_R,496+K_R-H_R),\\
 L_S&=\max(0,1024-H_S,496+K_S-H_S),\\
 L_{\min}&=\max(0,32-H_{\min}).
 \end{split}
 \tag{PG-Q32-profile-bounds}
\]
The last bound again comes from slicing parties $1,3$ at the
same coordinate, which produces an identity matrix of rank $32$.

Use the strict rational lower estimate $\sqrt2>7/5$ for the high
costs. In the order (PG-Q32-high-types) these lower costs are
\[
 112/5,\quad32,\quad224/5,\quad112/5,\quad32,\quad224/5.
\]
There are exactly $205$ nonnegative integer count vectors whose
sum of these lower costs is less than $135$. Summing
(PG-Q32-low-inequality) gives the lower bound
\[
 \underline c(H)+\frac{L_R+L_S+15L_{\min}}{17}.
 \tag{PG-Q32-enumerated-bound}
\]
The following table gives its minimum for each number of high
terms, excluding the single count vector $(1,0,0,1,0,0)$:
\begin{center}
\begin{tabular}{c|r|c}
High terms & Count vectors & Minimum rational lower bound\\\hline
0&1&$2528/17$\\
1&6&$11904/85$\\
2&20&$11904/85$\\
3&56&$12261/85$\\
4&76&$12261/85$\\
5&38&$13374/85$\\
6&7&$13374/85$
\end{tabular}
\end{center}
Every table entry is greater than $135$. The enumeration uses
only integer arithmetic and rational comparisons; its complete
$205$ rows and the defining expressions are recorded in
\path{outputs/pg_cycle_q32_exact.json}, generated by
\path{programs/pg_cycle_q32_exact.py}. An independent verifier,
\path{programs/rt_verify_cycle_q32.py}, reconstructs the six
types from all dyadic capacities and enumerates multisets by a
different recursive algorithm; its output
\path{outputs/rt_cycle_q32_verification.json} matches every
row, the grouped minima, and the sole exception. The table is an
exhaustive finite check of the explicit bounds, not an assertion
that any count vector is realized by tensor maps.

It remains to exclude one high resource $(512,1)$ and one
$(1,512)$. The low resources then have sums of $R$ and $S$ at
least $511$, and sum of minima at least $30$. Write
\[
 M=\sum\min(R,S),\qquad
 X=\sum_{R>S}(\sqrt{RS}-S),\qquad
 Y=\sum_{S>R}(\sqrt{RS}-R).
\]
Their capped cost is $W=M+X+Y$. We claim $W\geq90$.
Suppose otherwise. Since $M$ is an integer, $30\leq M\leq89$.
The oriented excess capacities satisfy
\[
 E_R=\sum_{R>S}(R-S)\geq511-M\geq422,
 \qquad E_S\geq422.
\]
For a term with $R>S$, put $\delta=\sqrt{RS}-S$. Then
$(R-S)/\delta=1+\sqrt{R/S}$. We show that $X<30$ would force
$E_R<422$.

If no ratio-$256$ term occurs in this orientation, then
$E_R<(1+8\sqrt2)30<422$. Otherwise the contribution of such a
term is $15S$. Since $S$ is a positive integer power of two,
$X<30$ permits exactly one such term, necessarily $(256,1)$.
It contributes $255$ to $E_R$ and $15$ to $X$.
If no ratio-$128$ term remains, the remaining ratio is at most
$64$, so $E_R<255+9\cdot15=390$. If a ratio-$128$ term
remains, its excess cost is $(8\sqrt2-1)S$. The remaining budget
below $15$ permits exactly one, namely $(128,1)$. The residual
excess-cost budget is less than $16-8\sqrt2<7$, excluding every
ratio-$64$ term, whose excess cost is at least seven.
The remaining ratios are therefore at most $32$, giving
\[
 E_R<255+127+(1+4\sqrt2)(16-8\sqrt2)
     =334+56\sqrt2<422.
\]
These contradictions show $X\geq30$. The symmetric argument
shows $Y\geq30$, and hence $W\geq M+60\geq90$, proving the
claim. Adding the two high resources yields a cover cost at least
$90+32\sqrt2>135$.

Thus every cover costs at least $135$. Finally,
$\sqrt2<99/70$ and the integer comparison $2^{99}<135^{14}$
give $32^{\sqrt2}<2^{99/14}<135$. Together with the earlier
width exclusions, this proves the assertion.
\end{proof}

\begin{lemma}[PG-CROSSCUT-RESOURCE]
\label{pg:crosscut-resource}
Suppose $U$ is a restriction of two opposite cycle wires of
dimensions $a,b$, the other two cycle dimensions being one.
Put $r=\min(q,a,b)$. Over every field,
\[
 \operatorname{rank}_{12|34}(C_q-U)
 \geq q^2-\frac{r(r+1)}2,
\]
where $12|34$ is the cut crossed by both nontrivial wires.
If $a,b$ are powers of two with uncapped product $R=ab<q^2$,
one may replace $r$ by
$\min(q,2^{\lfloor(\log_2R)/2\rfloor})$.
\end{lemma}
\begin{proof}
The output acts as $X\mapsto MX^TN$, with
$\operatorname{rank}M\leq\min(q,a)$ and
$\operatorname{rank}N\leq\min(q,b)$. If $\rho=\operatorname{rank}M$,
independent changes $X=P Y Q$ put the first factor into
$D=\operatorname{diag}(I_\rho,0)$. The fixed equation
$Y=DY^TN'$ forces all bottom rows to vanish. Its remaining blocks
satisfy $A=A^TN'_{11}$ and $B=A^TN'_{12}$. Over an algebraic
closure, choose $S$ so that $S N'_{11}S^{-1}$ is upper triangular
and write $A=S^T ZS$. Every strict-lower entry of $Z$ is then
determined by strict-upper entries, giving a fixed-space bound
$\rho(\rho+1)/2$. Apply the same argument to the transposed
fixed equation to replace $\rho$ by $\operatorname{rank}N$ if
that rank is smaller. Rank-nullity proves the first assertion.
For dyadic $a,b$, their minimum is at most
$2^{\lfloor(\log_2(ab))/2\rfloor}$, proving the refinement.
Capping a larger raw product at $q^2$ does not justify any smaller
bound on $\min(a,b)$.
\end{proof}

\begin{proposition}[PG-SLICE-CLIFFORD]
\label{pg:slice-clifford}
Let $q\geq3$ and suppose $C_q=\sum_jT_j$ over an arbitrary field.
Then
\[
 \sum_j\min\bigl(\operatorname{rank}_{12|34}T_j,
                   \operatorname{rank}_{14|23}T_j\bigr)\geq q+1.
 \tag{PG-SLICE-strict}
\]
In particular, $\chi_F(q)\geq\lceil(q+1)/2\rceil$.
\end{proposition}
\begin{proof}
Write each summand as a sum of rank-one tensors across its
cheaper crossing cut, and let $r$ be the total number of these
terms. Contract parties $1,3$ against independent variable vectors
$x,z\in F^q$. A term of the first cut becomes an outer product
of a vector linear in $x$ and a vector linear in $z$; for the
other cut the roles of $x,z$ are interchanged. Concatenating
these columns and rows gives homogeneous linear matrices
$L(x,z)$ of size $q\times r$ and $R(x,z)$ of size $r\times q$
with
\[
 L(x,z)R(x,z)=f(x,z)I_q,
 \qquad f(x,z)=\sum_{i=1}^q x_i z_i.
\]
The usual slice rank argument gives $r\geq q$. Suppose $r=q$.
Over the fraction field $F(x,z)$, the square matrices $L,R$ are
invertible; consequently $RL=fI_q$ there and hence as a
polynomial identity. Thus the homogeneous linear block matrix
\[
 D(x,z)=\begin{pmatrix}0&L(x,z)\\R(x,z)&0\end{pmatrix}
       =\sum_i x_iA_i+\sum_i z_iB_i
\]
acts on $V=F^{2q}$ and satisfies $D(x,z)^2=f(x,z)I_{2q}$.
Coefficient comparison, including in characteristic two, gives
\[
 \begin{gathered}
 A_i^2=B_i^2=0,\qquad
 A_iB_j+B_jA_i=\delta_{ij}I,\\
 A_iA_j+A_jA_i=0,
 \qquad B_iB_j+B_jB_i=0\quad(i\ne j).
 \end{gathered}
\]
Put $P_i=A_iB_i$ and $K_i=A_i+B_i$. The displayed relations
show that the $P_i$ are pairwise commuting idempotents,
$K_i^2=I$, and
\[
 K_iP_iK_i=I-P_i,
 \qquad K_iP_j=P_jK_i\quad(j\ne i).
\]
Commuting idempotents decompose $V$ as the direct sum of their
$2^q$ simultaneous eigenspaces, indexed by $\{0,1\}^q$.
The invertible map $K_i$ interchanges the spaces whose $i$th bits
differ and leaves all other bits unchanged. All $2^q$ spaces
therefore have the same dimension. Since $V\ne0$, this forces
$\dim V=2q\geq2^q$, a contradiction for $q\geq3$.
This elementary argument uses no division by two or assumption
that the field is infinite. It proves $r\geq q+1$.
For a two-cut decomposition with both prescribed ranks at most
$\chi_F(q)$, the sum of the two cheaper ranks is at most
$2\chi_F(q)$, giving the stated corollary.
\end{proof}

\begin{proposition}[PG-RECTANGULAR-FACTOR]
\label{pg:rectangular-factor}
Let $\rho_F(q)$ be the least inner dimension $r$ of homogeneous
linear matrices $L(x,z)$ and $R(x,z)$ of sizes $q\times r$ and
$r\times q$ satisfying $LR=(x\cdot z)I_q$. Over every field,
\[
 \chi_F(q)\leq\rho_F(q)\leq2\chi_F(q).
\]
Over $\mathbb C$ one also has
\[
 \tfrac12\sigma_{\mathbb C}(q)\leq\rho_{\mathbb C}(q)
 \leq2\sigma_{\mathbb C}(q).
\]
For $q\geq4$, every field satisfies $\rho_F(q)\geq q+2$.
Consequently the left-hand side of (PG-SLICE-strict) is at least
$q+2$ when $q\geq4$.
\end{proposition}
\begin{proof}
A two-cut decomposition with both ranks at most $\chi_F(q)$
produces the linear slice factorization from
Proposition~\ref{pg:slice-clifford}, with inner dimension at most
$2\chi_F(q)$. Conversely, split a general factorization as
$L=L_x+L_z$ and $R=R_x+R_z$. Comparing bidegrees gives
\[
 L_xR_z+L_zR_x=(x\cdot z)I_q.
\]
Contracting this matrix identity against the other two variable
vectors yields a two-cut decomposition, each cut rank at most
$r$, so $\chi_F(q)\leq r$.

For the complex comparison, setting $z=x$ and contracting against
$y$ on both sides expresses
$(x\cdot x)(y\cdot y)$ as a sum of $r$ products of bilinear
forms. Each product is a sum of two complex squares, giving
$\sigma_{\mathbb C}(q)\leq2r$. In the reverse direction, a
$s$-square identity can be written
$A(x)^TA(x)=(x\cdot x)I_q$, where $A(x)$ is an $s\times q$
linear matrix. Polarization gives
\[
 A(x)^TA(z)+A(z)^TA(x)=2(x\cdot z)I_q.
\]
Taking $L=[A(x)^T\ A(z)^T]$ and
$R=\tfrac12[A(z)^T\ A(x)^T]^T$ proves $\rho\leq2s$.
These two complex comparisons use division by two and are not
asserted for sums of squares in characteristic two.

It remains to exclude $r=q+1$; the equality $r=q$ was excluded
in Proposition~\ref{pg:slice-clifford}. Write $f=x\cdot z$ and
put
\[
 K=RL-fI_{q+1}.
\]
Over the fraction field, $F(x,z)^{q+1}$ is the direct sum of
$\operatorname{im}R$ and $\ker L$. The matrix $K$ vanishes on
the former and acts as $-f$ on the latter, so it has rank one.
Moreover $\operatorname{tr}K=-f$ in every characteristic.

A homogeneous rank-one quadratic matrix over the polynomial
unique factorization domain factors as $K=uv^T$, with $u,v$
homogeneous of degrees $a,2-a$, for some $a\in\{0,1,2\}$.
For completeness, divide a nonzero homogeneous column by the
greatest common divisor of its entries to obtain a primitive
homogeneous column $u$. Every other column is a rational multiple
of $u$; primitiveness forces this multiplier to be polynomial.
Homogeneity gives the indicated common degrees.
The identities $LK=KR=0$ and the trace give
\[
 Lu=0,\qquad v^TR=0,\qquad v^Tu=-f.
\]
If $u$ or $v$ is constant, a constant invertible change of the
inner basis removes one unused coordinate, reducing to an inner
dimension $q$ factorization, already impossible. Otherwise both
are linear, and there is a square linear completion
\[
 \begin{pmatrix}L\\v^T\end{pmatrix}
 \begin{pmatrix}R&-u\end{pmatrix}
 =fI_{q+1}.
\]
The creation--annihilation argument from
Proposition~\ref{pg:slice-clifford} applies to any square
factorization of $f$ of size $N$: its doubled space has dimension
$2N$, divisible by $2^q$. Here $N=q+1$, whereas
$2(q+1)<2^q$ for $q\geq4$. This contradiction excludes $r=q+1$.
The rank-one factorization argument is specific to this defect;
no analogous linear factorization of a higher-rank quadratic
matrix is assumed.
\end{proof}

Thus proving a superlinear lower bound for these rectangular
linear factorizations would already prove a superlinear lower
bound for the complex sums-of-squares length. The square-matrix
spinor obstruction alone cannot supply such a bound. In
particular, $RL=fI_r$ is false for a general rectangular
factorization: over the fraction field, $RL$ has rank $q<r$.
The algebraic comparisons above identify this route with the
same unresolved exponent gate, up to constant factors.

\begin{proposition}[PG-SINGH-CROSSING: independently audited linear bounds]
\label{pg:singh-crossing}
Let $P_F(q)$ be the minimum sum of the two prescribed crossing
ranks in a decomposition $C_q=S+U$. Equivalently,
\[
 P_F(q)=\min_{C_q=\sum_jT_j}
 \sum_j\min\bigl(\operatorname{rank}_{12|34}T_j,
                 \operatorname{rank}_{14|23}T_j\bigr).
\]
The bounds of Singh (2026, Theorems~1--3, pp.~3--5 and
Appendices~A--B), audited below, give over every field
\[
 P_F(q)\geq
 \begin{cases}
 \lceil4q/3\rceil,&q\geq3,\\
 \lceil15q/11\rceil,&q\geq5,\\
 \lceil67q/49\rceil,&q\geq7.
 \end{cases}
 \tag{PG-SINGH-bounds}
\]
Each applicable line is a lower bound. These remain linear
bounds and do not decide the exponent alternative in
Corollary~\ref{pg:crossing-exponent-gap}.
\end{proposition}
\begin{proof}
Group a decomposition's terms according to their cheaper cut
and use rank subadditivity to obtain one two-cut decomposition.
Conversely, decompose the two crossing matrices into rank-one
matrices. This proves the equality defining $P_F(q)$.
Interchanging parties $3,4$ sends our tensor $\delta_{13}\delta_{24}$
to Singh's $\delta_{14}\delta_{23}$, preserves the first allowed
cut, and sends the second to $13|24$. Thus his matrix realignment
quantity is exactly $P_F(q)$, with no change of normalization.

Here is an independent derivation of the first line. A decomposition
of cut ranks $r,s$ gives $q\times r$ matrices $X_i,Y_i$ and
$q\times s$ matrices $Z_i,W_i$ satisfying
\[
 X_iY_j^T+Z_jW_i^T=\delta_{ij}I_q.
 \tag{PG-SINGH-slices}
\]
Write $U_i=\operatorname{col}X_i$, $V_i=\operatorname{col}Z_i$,
$\widehat U_i=\operatorname{col}Y_i$, and
$\widehat V_i=\operatorname{col}W_i$. For $i\ne j$, the product
$[X_i\ Z_j][Y_j\ W_i]^T$ is zero, so
\[
 \dim(U_i+V_j)+\dim(\widehat U_j+\widehat V_i)\leq r+s.
 \tag{PG-SINGH-zero-product}
\]
For a diagonal slice, rank subadditivity also gives
$q\leq\dim\widehat U_i+\dim V_i$ and
$q\leq\dim U_i+\dim\widehat V_i$.
Choose distinct $a,b,c$. Two applications of submodularity of
subspace dimension give the path inequality
\[
 q+\dim V_b+\dim U_c
 \leq\dim(U_a+V_b)+\dim(U_c+V_b)+\dim(U_c+V_a).
\]
For the transposed slice system they give
\[
 q+\dim\widehat U_b+\dim\widehat V_c
 \leq\dim(\widehat U_b+\widehat V_a)
   +\dim(\widehat U_b+\widehat V_c)
   +\dim(\widehat U_a+\widehat V_c).
\]
Pair the three right-hand terms using
(PG-SINGH-zero-product) for $(i,j)=(a,b),(c,b),(c,a)$.
The two diagonal mixed-rank bounds make the left-hand sum at
least $4q$, while the right-hand sum is at most $3(r+s)$.
This proves $4q\leq3(r+s)$ over every field.

For the two stronger lines, we independently checked every
inequality and the exact weighted sums in Singh's printed
certificates. The five-index certificate has $189$ rows:
$75$ submodularity instances, $40$ monotonicity instances,
$10$ diagonal spanning inequalities, $38$ intersection-rank
inequalities, and $26$ applications of Sylvester's inequality.
The latter two types follow directly from, for $i\ne j$,
\[
 \operatorname{col}(X_iY_j^T)\subseteq U_i\cap V_j,
 \qquad X_iY_j^T=-Z_jW_i^T,
\]
and the transposed identities. Its exact rational weighted sum is
$\tfrac{15}{11}q-r-s\leq0$; every auxiliary dimension and rank
coefficient cancels exactly.

The seven-index certificate consists of $71$ submodularity or
monotonicity inequalities averaged over all $7!=5040$ index
permutations, together with six averaged diagonal,
intersection-rank, and Sylvester inequalities. The independent
checker enumerates the permutations, verifies every stated orbit
cardinality and the uniform averaging on subspace types, checks
the six remaining inequalities against the slice identities, and
obtains the exact sum $\tfrac{67}{49}q-r-s\leq0$.
These rational averages concern integer dimensions, not field
coefficients, so they do not require the field characteristic to
avoid any denominator.

The complete normalized mathematical certificate is
\path{inputs/pg_singh_rank_certificates.json}; the independent
semantic and arithmetic checker is
\path{programs/pg_verify_singh_certificates.py}, with result
\path{outputs/pg_singh_certificate_verification.json}.
The original author's Python checker was neither imported nor
executed. We did not run the Lean proofs and make no Lean-kernel
verification claim. In particular, the separate small-case claims
whose full proofs are only in the repository's Lean files are not
needed for the displayed bounds.
\end{proof}

The source was found after the independent
Propositions~\ref{pg:slice-clifford} and
\ref{pg:rectangular-factor} were derived. Its bounds are stronger
for the total crossing-rank quantity in the ranges above; the
rectangular one-sided factorization statements and their proofs
have been retained with their separate scope. The source's prose
about the signs of auxiliary coefficients has no effect on the
audited certificates: those coefficients are all exactly zero.

\begin{corollary}[PG-SINGH-RECTANGULAR]
\label{pg:singh-rectangular}
For positive integers $p,q$, let
$C_{p,q}=\operatorname{EPR}_{13}(p)\otimes
\operatorname{EPR}_{24}(q)$. Define
\[
 \beta(n)=
 \begin{cases}
 1,&n=1,2,\\
 4/3,&n=3,4,\\
 15/11,&n=5,6,\\
 67/49,&n\geq7.
 \end{cases}
\]
Over every field, any decomposition $C_{p,q}=\sum_jT_j$ obeys
\[
 \sum_j\min\{\operatorname{rank}_{12|34}T_j,
              \operatorname{rank}_{14|23}T_j\}
 \geq
 \left\lceil\max\{p\beta(q),q\beta(p)\}\right\rceil.
 \tag{PG-SINGH-rectangular}
\]
In particular, a two-cut decomposition with prescribed ranks
$r,s$ satisfies $r+s\geq4p/3$ when $q\geq3$,
$r+s\geq15p/11$ when $q\geq5$, and
$r+s\geq67p/49$ when $q\geq7$.
The analogous bounds with $p,q$ interchanged also hold.
\end{corollary}
\begin{proof}
First consider a two-cut decomposition. Contracting the two
$q$-dimensional parties with coordinate vectors indexed by
$i,j\in\{1,\ldots,q\}$ gives
\[
 X_iY_j^T+Z_jW_i^T=\delta_{ij}I_p,
 \tag{PG-SINGH-rectangular-slices}
\]
where $X_i,Y_i$ have size $p\times r$ and $Z_i,W_i$ have size
$p\times s$. Thus the ambient dimension of all column spaces is
$p$, while the number of available slice labels is $q$.
A single diagonal identity already gives $r+s\geq p$.

Put $U_i=\operatorname{col}X_i$, $V_i=\operatorname{col}Z_i$,
and let $\widehat U_i,\widehat V_i$ be the column spaces of
$Y_i,W_i$. For distinct labels $i,j$, zero-product rank gives
\[
 \dim(U_i+V_j)+\dim(\widehat U_j+\widehat V_i)\leq r+s.
\]
The diagonal identities give
$p\leq\dim\widehat U_i+\dim V_i$ and
$p\leq\dim U_i+\dim\widehat V_i$.
For three distinct labels $a,b,c$, submodularity gives
\begin{align*}
 p+\dim V_b+\dim U_c
 &\leq\dim(U_a+V_b)+\dim(U_c+V_b)+\dim(U_c+V_a),\\
 p+\dim\widehat U_b+\dim\widehat V_c
 &\leq\dim(\widehat U_b+\widehat V_a)
       +\dim(\widehat U_b+\widehat V_c)
       +\dim(\widehat U_a+\widehat V_c).
\end{align*}
Adding, pairing the three right-hand terms by the zero-product
inequality, and using the two diagonal mixed-rank bounds yields
$4p\leq3(r+s)$. Only the existence of three distinct labels was
used; the number of labels was not identified with $p$.

The same separation of roles applies to every row of the
independently audited five- and seven-label certificates in
Proposition~\ref{pg:singh-crossing}. Submodularity and monotonicity
are dimension-free. The diagonal spanning rows have right-hand
side $p$. Each off-diagonal intersection row follows from
\[
 \operatorname{col}(X_iY_j^T)\subseteq U_i\cap V_j,
 \qquad X_iY_j^T=-Z_jW_i^T,
\]
and the transposed identities. The Sylvester rows use the inner
matrix dimensions $r,s$ and the column-space dimensions; they
contain no occurrence of the number of labels. The averages over
five or seven labels simply relabel a chosen subset of the $q$
indices. Therefore the same exact rational weighted sums give
$\tfrac{15}{11}p-r-s\leq0$ and
$\tfrac{67}{49}p-r-s\leq0$, respectively. All auxiliary
coefficients cancel exactly as in the audited square case.
These averages are of integer ranks and dimensions, so no field
characteristic is excluded.

A cyclic relabeling of the four parties exchanges $p,q$ and
interchanges the two allowed cuts, giving the other orientation.
For an arbitrary decomposition, group each summand by its
cheaper cut. Rank subadditivity gives a two-cut decomposition
whose prescribed rank sum is at most the displayed sum of
minimum ranks. Applying both orientations and using integrality
proves the assertion.
\end{proof}

\begin{proposition}[PG-BORDER-SIGN: an integral border crossing identity]
\label{pg:border-sign}
Over every field, the four-dimensional crossing tensor is a
border sum of two tensors of ranks at most three across its two
prescribed cuts. In particular this holds in characteristic two.
This follows from an integral sign normalization of the tables
in Singh (2026, Remark~15 and the accompanying
\path{border/border_n4.json}); the excluded primes of the
original integer normalization are unnecessary for this identity.
\end{proposition}
\begin{proof}
Use the source's party order, so the target is
$S_4=\sum_{x,y=0}^3 e_{xyyx}$, and write $e_{ij}$ for the
standard basis of a sixteen-dimensional pair space.
The following twelve vectors give a complete integral formula:
\begin{align*}
 u_1&=e_{00}-t^2e_{33}+t^3e_{11},&
 v_1&=e_{11}-t^2e_{22}+t^3e_{00},\\
 u_2&=-t^2e_{23}+t^2e_{30}-t^3e_{12},&
 v_2&=-e_{21}+te_{03}-te_{32},\\
 u_3&=te_{03}-t^2e_{21}-t^2e_{32},&
 v_3&=-te_{12}-te_{23}+t^2e_{30},\\
 g_1&=-e_{01}+t^2e_{32}+t^3e_{10},&
 h_1&=e_{01}-t^2e_{32}-t^3e_{10},\\
 g_2&=t^2e_{22}-t^2e_{31}-t^3e_{13},&
 h_2&=-e_{31}-te_{13}+te_{22},\\
 g_3&=-te_{02}-t^2e_{20}-t^2e_{33},&
 h_3&=-te_{02}-te_{33}-t^2e_{20}.
\end{align*}
Set
\[
 B(t)_{xyuv}=\sum_{j=1}^3(u_j)_{xy}(v_j)_{uv},\qquad
 C(t)_{xyuv}=\sum_{j=1}^3(g_j)_{xu}(h_j)_{yv}.
\]
Direct multiplication gives the exact polynomial identity
\[
 B(t)+C(t)=t^3S_4+t^4E_4+t^5E_5,
 \tag{PG-BORDER-sign}
\]
where
\begin{align*}
 E_4&=-e_{2130}-e_{1203}+e_{2200}+e_{1133},\\
 E_5&=-e_{3300}-e_{1122}-e_{3120}-e_{1302}.
\end{align*}
The first tensor has rank at most three across $12|34$;
the second has rank at most three across $13|24$.
The leading coefficient is exactly one, and all entries are
integers. Therefore the same formal degeneration works over
any field, with no division and no restriction on characteristic.
Interchanging parties $3,4$ gives the stated crossing convention.
This is a border identity; the higher-degree error tensors do
not vanish in characteristic two.

The original table's low-degree coordinate equations have only
two forms: one positive product equal to the common leading
coefficient, or two products of opposite signs adding to zero.
Replacing every nonzero integer coefficient by its sign therefore
preserves these equations with leading coefficient one. The
formula above also verifies this directly, independently of the
original large coefficients.

The complete sparse integral arrays are in
\path{inputs/pg_singh_border_sign_normalized.json}.
The independent checker \path{programs/pg_verify_border_signs.py}
convolves and realigns their products, verifies all $1024$
coordinates in degrees zero through three, and records the entire
polynomial in \path{outputs/pg_singh_border_sign_verification.json}.
A separate reconstruction by
\path{programs/ex_verify_singh_signs.py} reproduced all these
coefficients, including the four errors in each of degrees four
and five; its result is \path{outputs/ex_singh_sign_check.json}.
The original author's checker was not executed.
\end{proof}

\paragraph{PG-BORDER-CYCLE-COST: explicit covers obtained by coefficient extraction.}
The border formula by itself does not give a cheap cycle cover.
The following direct calculation includes the non-dyadic bridge
ranks and all extraction overhead. After the party permutation,
$B(t)$ is a restriction of the path with cycle edge dimensions
$(4,3,4,1)$, and $C(t)$ of the rotated path with dimensions
$(1,4,3,4)$. Put the first pair factor at the inner vertex of each
path, the second at the other inner vertex, and use identity maps
at the two ends. All polynomial coefficients are then on parties
$2,3,4$; party $1$ has a constant map.

In the $N$th tensor power, a word containing $k$ first terms and
$N-k$ second terms is a restriction of the cycle resource with
edge dimensions
\[
 (4^k,\;3^k4^{N-k},\;4^k3^{N-k},\;4^{N-k}).
\]
Its three polynomial local maps have degrees at most
$3k,3N,3(N-k)$, respectively. Extracting the coefficient of
$t^{3N}$ from this word uses at most
$(3k+1)(3(N-k)+1)$ restrictions of the same resource: select
the degrees of the first and third polynomial maps, and the
middle degree is then determined. The full identity's leading
coefficient is $S_4^{\otimes N}$, so summing over words gives
an exact cover over the original field. No interpolation points,
field extensions, or divisions are needed.

Pad the two powers of three separately to the next powers of
two. The integer cycle widths are
\[
 \bigl(2k,\;2(N-k)+\lceil k\log_2 3\rceil,\;
 2k+\lceil(N-k)\log_2 3\rceil,\;2(N-k)\bigr).
\]
Each padded resource has cost at most
$2(4\sqrt3)^N$. The number of covering terms is at most
\[
 A_N=\sum_{k=0}^N\binom Nk(3k+1)(3(N-k)+1)
     =2^{N-2}(9N^2+3N+4).
\]
Consequently the exact finite covers satisfy
\[
 u_F(C_{4^N})\leq
 \frac{9N^2+3N+4}{2}(8\sqrt3)^N.
 \tag{PG-BORDER-cycle-cost}
\]
The extraction and dyadic padding cost only a polynomial factor
and a constant factor, respectively, but the $2^N$ choice of
crossing orientation remains in the exponential base.
Since $8\sqrt3>4^{\sqrt2}$, this construction does not meet the
strict finite-cycle gate. No tensor-power multiplicativity of
sums-of-squares identities or of the two-cut rank is assumed.

\begin{proposition}[PG-CROSSING-UNBALANCED]
\label{pg:crossing-unbalanced}
If $C_q=S+U$ has the two prescribed crossing ranks $r,s$, then,
over every field, for $q\geq2$,
\[
 r+3s\geq2q,\qquad 3r+s\geq2q.
\]
For every integer $1\leq t\leq q$ it also satisfies
\[
 r+t^2s\geq tq,\qquad t^2r+s\geq tq.
 \tag{PG-CROSSING-unbalanced}
\]
Consequently every cycle cover with capped crossing capacities
$(R_j,S_j)$ satisfies
\[
 \sum_j\min(R_j,3S_j)\geq2q,\qquad
 \sum_j\min(3R_j,S_j)\geq2q,
\]
and, for $1\leq t\leq q$,
\[
 \sum_j\min(R_j,t^2S_j)\geq tq,\qquad
 \sum_j\min(t^2R_j,S_j)\geq tq.
 \tag{PG-CYCLE-unbalanced}
\]
\end{proposition}
\begin{proof}
Use the slice matrices in (PG-SINGH-slices). A diagonal identity
implies $\operatorname{rank}X_i,\operatorname{rank}Y_i\geq q-s$.
For $i\ne j$, the identity
$X_iY_j^T=-Z_jW_i^T$ gives
$\operatorname{rank}(X_iY_j^T)\leq s$. Sylvester's inequality
then yields $2(q-s)-r\leq s$, which is the first assertion.
Exchanging the two cut orientations gives its transpose.

For the general family, select $t$ distinct indices. Stack the
corresponding $X_i$ vertically, and likewise the $Y_i$.
Their product $M$, a $tq\times tq$ block matrix with blocks
$M_{ij}=X_iY_j^T$, has rank at most $r$. The slice identity gives
$M+N=I_{tq}$, where $N_{ij}=Z_jW_i^T$.
Each of the $t^2$ blocks of $N$ has rank at most $s$.
Writing $N$ as the sum of the matrices supported on its
individual blocks therefore gives $\operatorname{rank}N\leq t^2s$.
Rank subadditivity now yields $tq\leq r+t^2s$; exchanging cuts
gives the other inequality.

Finally, assign each term in a cover to either of the two cut
orientations and sum within the two classes. Rank subadditivity
bounds the resulting two ranks by the sums of the assigned
capacities. Minimizing a weighted sum over all assignments acts
independently on each term and gives the displayed sums of
minima. Thus these linear profile inequalities include every
partition of the covering terms; no restriction to a single
chosen partition is being imposed.
\end{proof}

\paragraph{PG-CYCLE-Q64-PROFILE: correction of the earlier relaxation.}
The previously displayed profile
\[
 (4096,4),\quad(4,4096),\quad(64,64),\quad(16,16)
\]
has cost $336$ and minimum-capacity sum $88$, but it does not
satisfy Proposition~\ref{pg:crossing-unbalanced}. Assign the first
term to its rank-four cut and all remaining terms to the other
cut. This gives $s\leq4$, $r\leq84$, whereas $r+3s\geq128$ is
required. The stronger inequality with $t=8$ would require
$r\geq256$. Hence this profile is rigorously excluded; it is
not a surviving candidate for a cover.

The exploratory mixed-integer program first suggested
infeasibility at $q=64$. The following exact certificate replaces
that numerical diagnostic.

\begin{proposition}[PG-CYCLE-Q64]
\label{pg:cycle-q64}
Over every field,
\[
 u_F(C_{64})\geq359>64^{\sqrt2}.
\]
Consequently the finite cycle certificate in
Corollary~\ref{pg:cycle-finite} cannot have dyadic width $m\leq6$.
The first width not excluded by these results is $m=7$, or $q=128$.
\end{proposition}
\begin{proof}
Suppose a cover has cost below $359$, and cap each of its two
crossing capacities at $q^2=4096$. Capping only decreases the
cost lower bound $\sqrt{RS}$. Set $\alpha=140/99<\sqrt2$ and,
for dyadic capacities, define
\[
 w_\alpha(R,S)=
 \begin{cases}
 2^k,&RS=2^{2k},\\
 \alpha2^k,&RS=2^{2k+1}.
 \end{cases}
\]
This rational lower bound is at most $\sqrt{RS}$.
There are exactly six capacity pairs of ratio greater than
$1024$ that can occur:
\[
 (2048,1),(4096,1),(4096,2),
 (1,2048),(1,4096),(2,4096).
 \tag{PG-Q64-high}
\]
Their rational costs are respectively
$4480/99,64,8960/99$, repeated for the transposed pairs.
There are exactly $540$ nonnegative integer multiplicity vectors
whose rational high-resource cost is below $359$.

We give the complete rational inequality used on the remaining
resources. Put
$L_k(R,S)=\min(R,kS)+\min(kR,S)$. For every remaining capacity
pair, with ratio at most $1024$ and rational cost below $359$,
\begin{align*}
 F(R,S)={}&\frac{4201}{169740}(R+S)
 +\frac{17}{99}\min(R,S)+\frac8{33}L_3(R,S)\\
 &+\frac2{231}L_9(R,S)+\frac5{84}L_{16}(R,S)
 +\frac{32}{495}L_{49}(R,S)+\frac{40}{2829}L_{81}(R,S)
 \leq w_\alpha(R,S).
\end{align*}
The exact certificate checks this inequality on all $127$
possible remaining pairs. Its coefficients are nonnegative.
The terms $L_9,L_{16},L_{49},L_{81}$ use only
$t=3,4,7,9$ from Proposition~\ref{pg:crossing-unbalanced}.

Here are the precise right-hand sides for each high-resource
branch $H$. Before subtracting its contributions, the thirteen
individual rows in $F$ have lower bounds
\[
 (4096,4096,88,128,128,192,192,256,256,
   448,448,576,576).
\]
The first two are the crossing flattening ranks, the third is
$\lceil67\cdot64/49\rceil=88$ from
Proposition~\ref{pg:singh-crossing}, and the others are
(PG-CYCLE-unbalanced). Strengthen the first bound to
\[
 B_R(H)=\max\left\{
 4096,\;
 \max_{\substack{J\subseteq H\\\sum_{j\in J}S_j\leq2}}
       \left(2016+\sum_{j\in J}R_j\right),\;
 \max_{\substack{j\in H\\S_j=1}}
       \left(4096-\frac{\rho_j(\rho_j+1)}2+R_j\right)
 \right\},
\]
where
$\rho_j=\min\{64,2^{\lfloor(\log_2R_j)/2\rfloor}\}$;
use the transposed expression for $B_S(H)$.
These are exactly the rank-at-most-two removal bound and the
resource-sensitive rank-one bound. For a capped value
$R_j=4096$, using $\rho_j=64$ is valid regardless of the
uncapped product. For $R_j<4096$, the uncapped product is
$R_j$, which justifies the smaller value of $\rho_j$.
Since every capacity is at least one, each subcollection in the
second maximum has at most two terms.

For each of the thirteen rows, subtract the high-resource
contribution from its stated or strengthened lower bound and
replace a negative result by zero. Denote the resulting vector
by $b(H)$, and the thirteen coefficients of $F$ by $\lambda$.
Every cover with this high multiset then has cost at least
\[
 \sum_{j\in H}w_\alpha(R_j,S_j)+\lambda\cdot b(H).
 \tag{PG-Q64-branch}
\]
Fraction arithmetic verifies all $540$ branches and gives
\[
 \min_H\left\{\sum_{j\in H}w_\alpha(R_j,S_j)
                       +\lambda\cdot b(H)\right\}
 =\frac{422523176}{1089165}>387>359.
\]
The numbers of branches with zero through seven high resources
are respectively
\[
 1,\ 6,\ 21,\ 56,\ 121,\ 174,\ 121,\ 40.
\]
This contradicts the assumed cost below $359$.

The complete executable certificate is
\path{programs/pg_cycle_q64_simple_certificate.py}; it
reconstructs all capacity pairs, checks the single rational
inequality, enumerates every high multiplicity vector, and checks
every branch lower bound using only integers and fractions.
Its complete table is
\path{outputs/pg_cycle_q64_simple_exact.json}.
A separate implementation,
\path{programs/rt_verify_cycle_q64.py}, independently
reconstructs the capacity pairs, recursively enumerates the
multiplicity vectors by a different algorithm, checks all
removal subsets, and reproduces every bound and the same exact
minimum; its result is
\path{outputs/rt_cycle_q64_verification.json}.
No numerical optimizer or source program is used by this
certificate. The earlier multi-dual discovery certificate is
also retained as an independent implementation of the rational
bounds, but is not needed for this proof.
Finally $\sqrt2<99/70$ and the exact integer inequality
$2^{297}<359^{35}$ give
$64^{\sqrt2}<2^{297/35}<359$.
\end{proof}

\paragraph{PG-CYCLE-Q128-PROFILE: the next relaxation boundary.}
The same necessary-capacity constraints admit a profile at
$q=128$ of cost
\[
 712+128\sqrt2<900<128^{\sqrt2}.
\]
Its two capacity sums are both $17408$, and its minimum-capacity
sum is $176=\lceil67\cdot128/49\rceil$.
The full thirteen-type list and exact checks of every
$t=1,\ldots,128$ unbalanced inequality and every applicable
rank-at-most-two or refined rank-one removal are in
\path{programs/pg_verify_cycle_q128_profile.py} and
\path{outputs/pg_cycle_q128_profile_exact.json}.
The strict cost comparison uses $\sqrt2<10/7$,
$\sqrt2>24/17$, and $900^{17}<2^{168}$.
This is a feasible point of the stated necessary relaxation,
not a tensor cover; no local maps realizing it are known.

\paragraph{PG-CYCLE-AMPLIFICATION: direct conversion of a finite cover.}
The forward implication also has a direct construction which does
not use the CAT(0) price-space theorem. Write the number of covering
terms as $s$, their weights as $w_j$, and $Z=\sum_jw_j<q^{\sqrt2}$.
In the $N$th tensor power, group words by their multiplicities
$\nu=(\nu_1,\ldots,\nu_s)$. The sum for one profile is a
restriction of
\[
 \operatorname{GHZ}_{M_\nu}\otimes S_\nu,
 \qquad M_\nu=\binom{N}{\nu_1,\ldots,\nu_s},
\]
where $S_\nu$ has cycle widths $\sum_j\nu_j k_{e,j}$.
The selector orders the tensor factors locally for each word and
applies its corresponding local maps. This construction uses each
word once and requires no division by $M_\nu$.

Fix $0<\eta<2$. From $\widetilde Q(G_2)=4$ there is a constant
$B_\eta$ such that every $\operatorname{GHZ}_M$ is a restriction
of $G_2^{\otimes t}$ for some integer
\[
 t\leq\frac{\log_2M}{2-\eta}+B_\eta.
\]
Indeed, choose a sufficiently large fixed block with subrank per
copy above $2^{2-\eta}$, take its powers, and enlarge the constant
to handle the remaining bounded values of $M$.
Thus every profile is a restriction of one cycle-wire network.
Place its four vertices at $e_1,e_2,-e_1,-e_2$, so each cycle edge
has length $\sqrt2$. Its cost is at most
\[
 2\sqrt2\sum_j\nu_j\log_2w_j
 +\frac{4\sqrt2}{2-\eta}\log_2M_\nu+4\sqrt2 B_\eta.
\]
The multinomial expansion gives
$M_\nu\prod_jw_j^{\nu_j}\leq Z^N$, and
$\log_2M_\nu\leq N\log_2s$. Therefore the maximum profile
cost is at most
\[
 2\sqrt2 N\log_2Z
 +\frac{2\sqrt2\eta}{2-\eta}N\log_2s
 +4\sqrt2 B_\eta.
\]
Choose $\eta$ small enough that the coefficient of $N$ is less
than $4m$. There are at most $(N+1)^s$ profiles, whose external
selector toll is at most $4s\log_2(N+1)$. For sufficiently large
$N$ the resulting integer-width layouts satisfy the strict
planar shortening inequality for width $mN$. All their vertices
are the four terminals. The additional cycle widths supplying
$\operatorname{GHZ}_{M_\nu}$ are included in the displayed cost;
this does not assert that a profile can be implemented using an
unchanged source graph with unchanged widths.

\begin{corollary}[PG-CYCLE-SPECIALIZATION]
\label{pg:cycle-specialization}
Under the hypotheses of Corollary~\ref{pg:cycle-finite}, the
following statements are equivalent:
\begin{enumerate}
\item the Euclidean square is not a price over $\mathbb C$;
\item there is an integer $M$ such that, for infinitely many primes
$p$, a certificate (PG-CYCLE-finite) over $\mathbb F_p$ exists with
$m\leq M$;
\item there is an integer $M$ such that, for all sufficiently large
primes $p$, such a certificate exists with $m\leq M$.
\end{enumerate}
In particular, this statement does not determine the characteristic-two
case. If the square is a complex price but is absent in infinitely
many positive characteristics, their minimal certificate widths
must escape every fixed bound.
\end{corollary}
\begin{proof}
Suppose a complex certificate has dyadic dimension $q=2^m$ and
cost $Z<q^{\sqrt2}$. Its local-map equations have integer
coefficients. The existence of a complex solution therefore gives
one over a number field $K$, of some finite degree $d$. After
excluding finitely many primes to clear the denominators of the
finitely many coefficients, reduction at a prime of $K$ gives the
same cover over a residue extension $E/\mathbb F_p$ of degree
at most $d$. Its $N$th tensor power has cover cost $Z^N$.

The coefficient-descent argument from EX-13 applies to each of
its four-party restrictions with an overhead at most $d^3$.
To see the constant directly, choose an $F$-basis $\beta_b$ of
$E$ and an $F$-linear map $\ell:E\to F$ with $\ell(1)=1$.
Expand the first three grouped local maps in this basis. For each
triple of basis indices, apply $\ell$ entrywise to the fourth map
multiplied by the product of the corresponding three basis
elements. Summing at most $d^3$ resulting $F$-restrictions is
exactly the application of $\ell$ to the original identity.
The grouped tensor-power maps still have coefficients in $E$,
so the overhead remains $d^3$, independent of $N$.
Consequently the descended cycle cover has cost at most
$d^3Z^N$. Choose one fixed $N$ with
$d^3Z^N<q^{N\sqrt2}$. It works at every remaining prime and
has the uniform dyadic width $mN$. This proves (1)$\Rightarrow$(3),
and (3)$\Rightarrow$(2) is immediate.

For (2)$\Rightarrow$(1), the bounded widths and costs allow only
finitely many resource formats, as in
Corollary~\ref{pg:cycle-finite}. Some fixed format therefore has
solutions at infinitely many primes. If its integer-coefficient
local-map equations had no complex solution, the
Nullstellensatz would put $1$ in their ideal over $\mathbb Q$.
Clearing denominators would put a nonzero integer in the ideal
over $\mathbb Z$, excluding solutions at all primes except its
finitely many prime divisors. This is a contradiction. The
complex solution has exactly the same strict cost bound, so
Corollary~\ref{pg:cycle-finite} excludes the complex square.
\end{proof}

\begin{corollary}[PG-CROSSING-EXPONENT-GAP]
\label{pg:crossing-exponent-gap}
Assume the checked source chains used in
Theorem~\ref{pg:cycle-gap} and Propositions~\ref{rt:sos}
and~\ref{rt:sos-upper}. For the two-cut rank $\chi_F(q)$ defined
in Section~\ref{sec:root-sos}, the limit exists and satisfies
\[
 \lim_{q\to\infty}\frac{\log\chi_F(q)}{\log q}
 =
 \begin{cases}
 1,&\text{if the Euclidean square is not a price over }F,\\
 \sqrt2,&\text{if the Euclidean square is a price over }F.
 \end{cases}
 \tag{PG-CROSSING-exponent}
\]
Over $\mathbb C$, the same alternative holds for the bilinear
sums-of-squares length $\sigma_{\mathbb C}(q)$. The latter
assertion is not made in characteristic two.
\end{corollary}
\begin{proof}
For one cycle resource $S_{\mathbf k}$, the two crossing
flattening ranks are
\[
 R=2^{k_{23}+k_{41}},\qquad
 S=2^{k_{12}+k_{34}},\qquad w(\mathbf k)=\sqrt{RS}.
\]
In an arbitrary cover of $C_q$, put each restricted term into
the first sum if $R\leq S$, and into the second sum otherwise.
Rank subadditivity on the respective cuts gives a two-term
decomposition whose maximum cut rank is at most
$\sum_j\min(R_j,S_j)\leq\sum_j\sqrt{R_jS_j}$.
Taking the infimum proves
\[
 \chi_F(q)\leq u_F(C_q).
\]
For the elementary reverse exponent bound, fix the same basis
coordinate at parties $1,3$ in $C_q$. The remaining matrix on
parties $2,4$ is the $q\times q$ identity. A summand separated
across either crossing cut becomes rank one after this slice.
Thus any two-cut decomposition with both cut ranks at most $R$
satisfies $q\leq2R$, proving $\chi_F(q)\geq q/2$ over every field.

If the square is absent, submultiplicativity and
$U_F(C_2)=2$ give
$u_F(C_{2^N})=2^{N+o(N)}$. Hence
$\chi_F(2^N)=2^{N+o(N)}$.
For arbitrary $q$, restrict from the next power of two; its
factor at most two does not change the exponent. Together with
the lower bound $q/2$, this proves the first case.

If the square is present, its equivalence with absence of a
planar strict shortening and Proposition~\ref{rt:sos} give
$\chi_F(q)\geq q^{\sqrt2-o(1)}$.
The integral crossing identity, in the grouped form of
Proposition~\ref{rt:sos-upper}, gives the matching upper bound
$\chi_F(q)\leq q^{\sqrt2+o(1)}$. Finally,
$\chi_{\mathbb C}\leq\sigma_{\mathbb C}\leq4\chi_{\mathbb C}$
transfers the two exponent values to the complex sums-of-squares
length. No polarization in characteristic two is used.
\end{proof}

Thus finding any fixed positive improvement below the
$\sqrt2$ exponent would, through these conditional chains,
force the crossing exponent all the way to one. The argument
does not supply such an improvement, nor does a planar square
alone establish flatness in every dimension.

\subsection{A uniform multi-cut consequence of the crossing identity}

For the matching tensor
$T_{n,d}=\bigotimes_{i=1}^d\operatorname{EPR}_{i^+i^-,n}$,
call a bipartition \emph{transversal} if each side contains exactly one
endpoint of every matching edge. Let $P_F(n,d)$ be the minimum number
of tensors separated across transversal bipartitions whose sum is
$T_{n,d}$. Equivalently, group all summands with the same bipartition
and minimize the sum of their respective flattening ranks. This is
precisely the partial-transpose rank of $I_{n^d}$ in Rossman and Zhu
(2026, Definition~3.4 and Lemma~3.5, p.113:8).

\begin{proposition}[PG-MULTICUT-UPPER]
\label{pg:multicut-upper}
The checked integral crossing identity implies that there is an
absolute constant $C$ such that, over every field and for all integers
$n\geq2$, $d\geq1$,
\[
 P_F(n,d)\leq
 2^{d-1} n^{\sqrt d}\bigl(\log(n+2)\bigr)^{C\sqrt d}.
 \tag{PG-MULTICUT-upper}
\]
Here and in this proof logarithms are natural. In particular,
$P_F(n,d)\leq n^{\sqrt d+o(1)}$ for each fixed $d$. Uniformly for
$1\leq d\leq O(\log n/\log\log n)$, the upper bound is
$n^{(1+o(1))\sqrt d}$. The same bound holds for a Kronecker product
of $d$ arbitrary $n\times n$ matrices.
\end{proposition}
\begin{proof}
First quantify the rectangular version of RT-SOS-UPPER. Frobenius
symbols of energy $2k$ are in bijection with partitions of $k$, so
\[
 P(E)=\sum_{0\leq k\leq E/2}p(k),\qquad
 \log P(E)=\kappa\sqrt E+O(\log(E+2)),\quad
 \kappa=\pi/\sqrt3.
 \tag{PG-MULTICUT-count}
\]
The latter follows directly from the Hardy--Ramanujan formula
$p(k)\sim e^{\pi\sqrt{2k/3}}/(4\sqrt3 k)$ (Hardy \& Ramanujan,
1918, Eq.(1.41), p.79), bounding the sum between its last summand
and $(E+1)$ times its largest summand. Let $Q(E)$ be the largest
cardinality of a class of fixed size and energy at most $E$.
There are at most $(E+1)^2$ classes, whence
$P(E)/(E+1)^2\leq Q(E)\leq P(E)$.
For $p\geq2$, choose the least $E_p$ with $Q(E_p)\geq p$.
These estimates imply
\[
 E_p=O((\log(p+2))^2),\qquad
 \kappa\sqrt{E_p}\leq\log p+O(\log\log(p+2)).
\]
For the second inequality use $Q(E_p-1)<p$ and absorb the bounded
change in the square root. All constants here are absolute.

Use classes of energies $e_p\leq E_p$ and $e_q\leq E_q$ for the
crossing wires of dimensions $p,q$. Group the coefficients of
Anthropic Proposition~B.2 at $N\leq e_p+e_q$ and at larger $N$.
Each grouped tensor has rank at most
\[
 (e_p+e_q+1)^2P(e_p+e_q)
\]
on its corresponding crossing cut. Restrict the four index sets
to dimensions $p,q,p,q$. If $R(p,q)$ is the larger of these two
rank bounds, the preceding estimates and the Euclidean triangle
inequality give
\[
 \log R(p,q)\leq
 \sqrt{(\log p)^2+(\log q)^2}
 +C_0\log(\log p+\log q+2)
 \tag{PG-MULTICUT-merge}
\]
for an absolute $C_0$. If either dimension is one, simply retain
the other wire and use no error term or additional orientation.
All coefficient identities are integral before reduction to the field.

Now combine groups of matching edges on a binary tree. For a fixed
transversal cut of a child group, a rank-$p$ tensor has a factorization
$(L\otimes R)\operatorname{EPR}_p$, with $L,R$ arbitrary maps to the
respective endpoint groups. For another child write similarly
$(A\otimes B)\operatorname{EPR}_q$. Apply the two-cut decomposition
above to these two intermediate wires. After composing the four maps,
the new cuts are
\[
 (L\cup A)\mid(R\cup B),\qquad
 (L\cup B)\mid(R\cup A),
\]
both transversal on the union of the child groups. Each resulting
rank is at most $R(p,q)$. Fixing a canonical side for each child cut
removes complement duplication: every parent transversal cut determines
its two child cuts and the relative flip. Starting with one rank-$n$
wire at each leaf therefore yields at most $2^{d-1}$ patterns, each
with the rank bound obtained by applying
(PG-MULTICUT-merge) at every active internal node. This argument uses
rank factorizations and actual tensor equalities; no polarization of
squares or division by two is involved.

For completeness, the error bound is uniform in the number of leaves.
Pad the consecutive $d$ active leaves to the next power of two with
dimension-one wires and use the balanced binary tree. Put $m=\log n$.
At height $j$ from the leaves, let $\boldsymbol\ell_j$ be the vector
of logarithmic rank bounds. Replacing each pair by its Euclidean norm
preserves the $\ell^2$ norm of the whole vector. The additional error
vector at height $j$ has at most $2d/2^j$ nonzero entries: only a node
with two nonempty children needs a crossing decomposition. A crude
induction from (PG-MULTICUT-merge) bounds the largest logarithmic rank
at height $j$ by $K^j(m+K)$ for an absolute $K>1$. Hence each error
entry is at most
$K'(\log(m+2)+j)$, for another absolute constant. Minkowski's inequality
at each level now gives
\[
 \log R_{\rm root}
 \leq\sqrt d\,m+
 K'\sqrt{2d}\sum_{j\geq1}2^{-j/2}
                         (\log(m+2)+j)
 \leq\sqrt d\,\bigl(m+K''\log(m+2)\bigr).
\]
Both geometric series converge. Expanding the final rank factorizations
and multiplying by the number $2^{d-1}$ of patterns proves
(PG-MULTICUT-upper), after increasing its constant to cover small $n$.
Any $n\times n$ matrix is obtained from $I_n$ by a linear map on one
endpoint, and such maps preserve every transversal rank bound. This
proves the last assertion.
\end{proof}

Rossman and Zhu (2026, Proposition~6.2 and Theorem~6.3, p.113:18)
obtain $n^{O(d^{0.7})}$ from Hrube\v{s}'s bilinear sums-of-squares
construction in characteristic different from two. Their introduction
(p.113:4) leaves the identity matrix over $\mathbb F_2$ as a possible
candidate for PT-rank $n^{d-o(\log d)}$ in the growing-$d$ regime.
The checked integral crossing identity, through the proposition above,
excludes this particular candidate in that regime as well. This is a
consequence of the supplied crossing identity, not an independent
verification of the paper's claimed complexity separation.
The resulting logarithmic rank is $\sqrt d\log n$ to leading order,
which is exactly the critical two-hub scale for $d$ orthogonal
antipodal demands. The factor $2^{d-1}$ counts orientation patterns;
it is not a free coherent selector. No strict shortening or new price
is obtained from this upper bound.


\begin{corollary}[PG-MULTICUT-FLAT: conditional sharpness]
\label{pg:multicut-flat}
Fix $d\geq2$. If no strict shortening exists in dimension $d$, then,
over the field in question,
\[
 \lim_{n\to\infty}\frac{\log P_{\mathbb F}(n,d)}{\log n}=\sqrt d.
 \tag{PG-MULTICUT-flat}
\]
In particular this conclusion holds if all finite Euclidean
$d$-dimensional configurations are prices. No converse is asserted.
\end{corollary}
\begin{proof}
Proposition~\ref{pg:multicut-upper} gives the upper bound. Suppose instead
that some fixed $r<\sqrt d$ and an unbounded sequence of integers $n$
satisfy $P_{\mathbb F}(n,d)\leq n^r$.
Put $m=\lfloor\log_2n\rfloor$ and restrict each wire to dimension $2^m$.
Grouping a minimum transversal decomposition by its cut expresses this
restricted matching tensor as a sum of at most $A=2^{d-1}$ tensors, each
with rank at most $n^r$ on its selected cut. Each such tensor is a
restriction of a two-hub network: its $2d$ terminal wires have width $m$,
and its bridge has width at most
$k=\lceil r\log_2n\rceil$. Thus $k/m\to r$ along the sequence.

Place the terminals at $\pm e_1,\ldots,\pm e_d$.
A transversal cut chooses signs $s_i\in\{1,-1\}$, with selected terminals
$s_ie_i$ on one side. Put its two hubs at
$\pm\tau(s_1,\ldots,s_d)/\sqrt d$, for fixed $0<\tau<1$.
Every terminal-to-hub distance is
$L(\tau)=\sqrt{1+\tau^2-2\tau/\sqrt d}$, and the bridge length is $2\tau$.
Consequently every term has wire cost at most
\[
 2dmL(\tau)+2\tau k.
\]
The derivative at $\tau=0$ of $2dL(\tau)+2r\tau$ is
$2(r-\sqrt d)<0$. Choose a sufficiently small fixed $\tau$ so that this
expression is less than $2d$. Since $k/m\to r$, the same strict
inequality, with a fixed margin, holds for the normalized actual costs
at all sufficiently large $m$. The selector toll
$2d\lceil\log_2A\rceil$ is constant in $m$ and is absorbed by that margin.
This is a strict shortening in dimension $d$, a contradiction.
The final assertion uses the previously proved implication from a
Euclidean $d$-price configuration to absence of strict $d$-shortening.
\end{proof}
